% Options for packages loaded elsewhere
\PassOptionsToPackage{unicode}{hyperref}
\PassOptionsToPackage{hyphens}{url}
\PassOptionsToPackage{dvipsnames,svgnames,x11names}{xcolor}
%
\documentclass[
  10pt,
]{article}
\usepackage{amsmath,amssymb}
\usepackage{iftex}
\ifPDFTeX
  \usepackage[T1]{fontenc}
  \usepackage[utf8]{inputenc}
  \usepackage{textcomp} % provide euro and other symbols
\else % if luatex or xetex
  \usepackage{unicode-math} % this also loads fontspec
  \defaultfontfeatures{Scale=MatchLowercase}
  \defaultfontfeatures[\rmfamily]{Ligatures=TeX,Scale=1}
\fi
\usepackage{lmodern}
\ifPDFTeX\else
  % xetex/luatex font selection
\fi
% Use upquote if available, for straight quotes in verbatim environments
\IfFileExists{upquote.sty}{\usepackage{upquote}}{}
\IfFileExists{microtype.sty}{% use microtype if available
  \usepackage[]{microtype}
  \UseMicrotypeSet[protrusion]{basicmath} % disable protrusion for tt fonts
}{}
\makeatletter
\@ifundefined{KOMAClassName}{% if non-KOMA class
  \IfFileExists{parskip.sty}{%
    \usepackage{parskip}
  }{% else
    \setlength{\parindent}{0pt}
    \setlength{\parskip}{6pt plus 2pt minus 1pt}}
}{% if KOMA class
  \KOMAoptions{parskip=half}}
\makeatother
\usepackage{xcolor}
\usepackage[margin=2.3cm]{geometry}
\usepackage{longtable,booktabs,array}
\usepackage{calc} % for calculating minipage widths
% Correct order of tables after \paragraph or \subparagraph
\usepackage{etoolbox}
\makeatletter
\patchcmd\longtable{\par}{\if@noskipsec\mbox{}\fi\par}{}{}
\makeatother
% Allow footnotes in longtable head/foot
\IfFileExists{footnotehyper.sty}{\usepackage{footnotehyper}}{\usepackage{footnote}}
\makesavenoteenv{longtable}
\setlength{\emergencystretch}{3em} % prevent overfull lines
\providecommand{\tightlist}{%
  \setlength{\itemsep}{0pt}\setlength{\parskip}{0pt}}
\setcounter{secnumdepth}{-\maxdimen} % remove section numbering
% Extra preamble for the pandoc LaTeX template (pdflatex).
\usepackage{amsmath,amssymb}
\usepackage{graphicx}
\usepackage{cite}                 % numeric citations, ranges compressed: [3--5]
\usepackage{caption}
\captionsetup{font=small,labelsep=period,justification=raggedright,singlelinecheck=false}
\usepackage{etoolbox}
\usepackage{framed}
\definecolor{shadecolor}{gray}{0.95}
\AtBeginEnvironment{longtable}{\small}
\setlength{\emergencystretch}{3em}
% allow line breaks after underscores (column names in \texttt)
\DeclareRobustCommand{\_}{\textunderscore\hspace{0pt}}
% characters used in the text that the utf8 input encoding does not define
\DeclareUnicodeCharacter{2212}{\ensuremath{-}}
\DeclareUnicodeCharacter{2192}{\ensuremath{\rightarrow}}
\DeclareUnicodeCharacter{03C1}{\ensuremath{\rho}}
\DeclareUnicodeCharacter{03B2}{\ensuremath{\beta}}
\DeclareUnicodeCharacter{0394}{\ensuremath{\Delta}}
\DeclareUnicodeCharacter{03C3}{\ensuremath{\sigma}}
\DeclareUnicodeCharacter{03BA}{\ensuremath{\kappa}}
\DeclareUnicodeCharacter{2265}{\ensuremath{\geq}}
\DeclareUnicodeCharacter{2264}{\ensuremath{\leq}}
\DeclareUnicodeCharacter{2248}{\ensuremath{\approx}}
\DeclareUnicodeCharacter{00B3}{\ensuremath{{}^{3}}}
\DeclareUnicodeCharacter{2011}{\mbox{-}}
\DeclareUnicodeCharacter{202F}{\,}
\ifLuaTeX
  \usepackage{selnolig}  % disable illegal ligatures
\fi
\IfFileExists{bookmark.sty}{\usepackage{bookmark}}{\usepackage{hyperref}}
\IfFileExists{xurl.sty}{\usepackage{xurl}}{} % add URL line breaks if available
\urlstyle{same}
\hypersetup{
  pdftitle={{[}Title{]}},
  colorlinks=true,
  linkcolor={blue},
  filecolor={Maroon},
  citecolor={Blue},
  urlcolor={blue},
  pdfcreator={LaTeX via pandoc}}

\title{{[}Title{]}}
\usepackage{etoolbox}
\makeatletter
\providecommand{\subtitle}[1]{% add subtitle to \maketitle
  \apptocmd{\@title}{\par {\large #1 \par}}{}{}
}
\makeatother
\subtitle{Supplementary
Information}
\author{}
\date{}

\begin{document}
\maketitle

\section{1 About
these tables}

This document
describes the 49
supplementary tables
(\protect\hyperlink{S1}{S1}
to
\protect\hyperlink{S49}{S49})
and the seven data
tables of the
Results (T5 to T11)
of the paper. It
does not reprint
them. Every one of
these tables is
released as a CSV
file (Tables 1 to 4,
in the Methods, are
descriptive and are
not released as
files); each entry
below gives the
path, the number of
rows and columns,
what a row is, what
the key columns
mean, a three-row
excerpt, the places
in the paper that
cite the table, and
caveats.
Supplementary tables
are numbered in the
order in which the
Methods and Results
first cite them.

\begin{verbatim}
README.md          what is released, what was changed, licence
INDEX.csv          one row per table: id, file, title, rows, columns, licence flag, sha256
main/              T5 ... T11    Tables 5 to 11 of the paper
supplementary/     S1 ... S49    supplementary tables, one CSV per table
\end{verbatim}

\textbf{How to read
an entry.}
\emph{File} is the
path of the CSV, its
size and its licence
flag.
\emph{Supports} is
the part of the
paper the table
backs. \emph{Cited
in} lists the
sections of the
Methods, Results and
Discussion that cite
it (a link opens the
paper at that
section). \emph{A
row is} says what
one row represents.
\emph{Key columns}
explains the columns
that matter; the
shared provenance
block (below) is not
repeated. The
\emph{excerpt} shows
real rows from the
file.

\section{2
Conventions shared
by the tables}

\textbf{Shared
provenance columns.}
The result tables
(T5 to T11 and the
supplementary tables
whose entry says so)
carry the same block
of up to 12 columns,
so that a row can be
traced to its
denominator:
\texttt{claim\_type},
\texttt{unit} and
\texttt{n\_units}
(what was counted
and how many),
\texttt{n\_clusters}
and
\texttt{cluster\_def}
(the resampling
unit: enzyme
sub-subclass for
natural enzymes,
campaign and plate
well or backbone
cluster for designs,
position for BglB),
\texttt{k},
\texttt{N} and
\texttt{k\_of\_N} (a
count with its
denominator),
\texttt{denominator\_def},
\texttt{headline\_ok}
(1 if the row may be
quoted as a headline
result; 0 for
per-metric detail,
descriptive
analyses, reference
rows and
wrong-ligand
AUROCs),
\texttt{status} and
\texttt{licence\_flag}.

\textbf{Claim
types.}
\texttt{detection}
(can the metric see
damage),
\texttt{specificity}
(catalytic against
matched control),
\texttt{equivalence}
(panel ratio against
1),
\texttt{discrimination}
(a real contrast,
dead against
active),
\texttt{ranking}
(against a measured
label),
\texttt{enrichment},
\texttt{baseline} (a
reference row),
\texttt{design},
\texttt{audit},
\texttt{not\_measurable}.

\textbf{Applicability
status} (in
\protect\hyperlink{S3}{S3},
\protect\hyperlink{S4}{S4}
and the
\texttt{applicability\_status}
columns).
\texttt{OK} can be
tested;
\texttt{OK-KNOCKON}
tested but reported
as displacement of
the other catalytic
residues, with no
specificity verdict;
\texttt{OK-PROVISIONAL}
the raw status of
the interface terms
in the re-predicted
ladder, resolved by
testability;
\texttt{GLOBAL}
tested but has no
site input, so it is
a comparator;
\texttt{REF} a fixed
reference row;
\texttt{X-BOOKKEEP}
a counter or
constant;
\texttt{X-NOINPUT}
the input cannot
reach the changed
quantity;
\texttt{X-ARM} the
metric belongs to
the other kind
(structure space or
prediction) than the
test;
\texttt{X-UNDEF-COV}
defined on too few
units;
\texttt{X-WITHDRAWN}
the test was
withdrawn;
\texttt{X-NOCTRL} no
matched control
exists;
\texttt{X-CONFOUND}
the control is
confounded;
\texttt{NOT-RUN};
\texttt{DUP}
identical to another
test and read once.

\textbf{Verdicts in
the specificity
tables.}
\texttt{specific}
(response
significant,
interval of the
specificity ratio
above 1, control
adequate);
\texttt{non\_specific}
(responds, not more
than the control);
\texttt{blind} (does
not respond);
\texttt{invariant}
(does not change);
and counts-floor
labels for cells
with too few units.

\textbf{Why some
metrics are not in
the main set.} A
metric with no
residue-set input
(whole-protein
energies, sequence
composition,
whole-prediction
confidence) cannot
respond
\emph{specifically}
to catalytic
residues, and a
metric whose input
cannot reach the
changed quantity
cannot register the
lesion at all
(sequence
composition under a
rotamer change, for
example); such
metrics are kept as
comparators. One
further quantity
passes these tests
but is kept out of
the main set: the
Chai-1 combined
score, which Chai-1
uses to rank its own
models and which
equals 0.2 × pTM +
0.8 × ipTM − 100 ×
(inter-chain clash
flag), so that it is
a rescaled ipTM and
not an independent
metric. It is still
scored, and its role
and reason are in
\protect\hyperlink{S1}{S1}.

\textbf{Names that
recur.}
\texttt{canonical\_key}
is the distinct
metric after literal
aliases are merged
(\texttt{metric\_name}
may be an alias);
\texttt{in\_main\_set}
is 1 for the 33
reported metrics;
\texttt{role} is
\texttt{main},
\texttt{reference}
or a supplement role
(see
\protect\hyperlink{S1}{S1});
\texttt{source} is
structure-space,
prediction-based,
PLACER, baseline,
natural or denovo as
the table says;
\texttt{lesion\_test}
is the catalytic
lesion or the
second-shell lesion;
\texttt{level} or
\texttt{lesion\_step}
is the lesion step
(isosteric 15.3,
non-isosteric 36.4,
Ala 53.2, Gly 80.5
ų median side-chain
volume change;
\texttt{second\_shell\_1/2/4}
is the second-shell
dose);
\texttt{ame\_flavour}
is \texttt{crystal}
(against the
deposited structure)
or
\texttt{self\_design}
(against a design's
own model) and the
two are never
pooled;
\texttt{sr\_*} is a
specificity ratio
with its interval;
\texttt{auroc} is
direction-free
unless a column says
\texttt{signed} or
\texttt{directed};
\texttt{rho} is a
Spearman
correlation.

\textbf{Test names.}
Tables that refer to
a test use the names
of
\protect\hyperlink{S6}{S6},
a name followed by
its class in square
brackets, for
example
\texttt{main\ set\ {[}catalytic\ lesion,\ structure-space{]}}.

\textbf{Licence.}
Tables flagged
\texttt{D2DCure\_aggregate}
in
\texttt{INDEX.csv}
are built from the
D2DCure BglB data,
whose licence is
unstated; they hold
metric-level
aggregates only, no
per-variant row. All
other tables are
\texttt{public}.

\textbf{Reproducibility.}
The tables were
generated by script
from result files
with seeds and input
hashes recorded;
\texttt{INDEX.csv}
gives the sha256 of
every released file.

\section{3 Where do
I find \ldots?}

{\small

\begin{longtable}[]{@{}
  >{\raggedright\arraybackslash}p{(\columnwidth - 2\tabcolsep) * \real{0.5455}}
  >{\raggedright\arraybackslash}p{(\columnwidth - 2\tabcolsep) * \real{0.4545}}@{}}
\toprule\noalign{}
\begin{minipage}[b]{\linewidth}\raggedright
I want to know
\end{minipage} &
\begin{minipage}[b]{\linewidth}\raggedright
Go to
\end{minipage} \\
\midrule\noalign{}
\endhead
\bottomrule\noalign{}
\endlastfoot
Which 33 metrics are
reported, and why
the rest are not &
\protect\hyperlink{S1}{S1},
\protect\hyperlink{S2}{S2};
then
\protect\hyperlink{S3}{S3}
and
\protect\hyperlink{S4}{S4}
for the reason in
each test \\
What a metric reads,
and over which
residues it is
scored &
\protect\hyperlink{S2}{S2} \\
The 30 tests, with
comparator, control
type and role &
\protect\hyperlink{S6}{S6} \\
The specificity
ratio of one metric
at one lesion step &
\protect\hyperlink{S33}{S33}
(comparators:
\protect\hyperlink{S34}{S34}) \\
Whether the controls
were good &
\protect\hyperlink{S26}{S26},
\protect\hyperlink{S27}{S27}
(structure space);
\protect\hyperlink{S37}{S37},
\protect\hyperlink{S39}{S39},
\protect\hyperlink{S40}{S40},
\protect\hyperlink{S41}{S41},
\protect\hyperlink{S38}{S38}
(predictor
ladder) \\
How fragile the
structure-space
specificity result
is &
\protect\hyperlink{S31}{S31},
\protect\hyperlink{S42}{S42},
\protect\hyperlink{S29}{S29} \\
How a metric
responds to the
substrate swap, and
the role of ligand
size & Table 7,
\protect\hyperlink{S14}{S14},
\protect\hyperlink{S15}{S15} \\
The distance control
for the interface
terms & Table 10,
\protect\hyperlink{S30}{S30},
\protect\hyperlink{S37}{S37},
\protect\hyperlink{S39}{S39},
\protect\hyperlink{S40}{S40},
\protect\hyperlink{S41}{S41} \\
Every metric on the
21-pair evaluation
set (complete Table
5) &
\protect\hyperlink{S11}{S11} \\
Every metric on the
192 plated designs,
active against no
active (complete
Table 6) &
\protect\hyperlink{S13}{S13} \\
Every quantity
analysed on the 432
BglB variants
(complete Table 8) &
\protect\hyperlink{S16}{S16} \\
Every metric on the
16 plated designs
with a measured
activity (complete
Table 9) &
\protect\hyperlink{S21}{S21} \\
Hits per plate when
a metric fills a
96-well plate &
\protect\hyperlink{S20}{S20} \\
The wider set of 49
pairs, the 28
re-predicted pairs,
the pair-set
definitions &
\protect\hyperlink{S7}{S7},
\protect\hyperlink{S8}{S8},
\protect\hyperlink{S9}{S9} \\
BglB: every metric,
prediction-based
metrics, sensitivity
analyses &
\protect\hyperlink{S17}{S17},
\protect\hyperlink{S19}{S19},
\protect\hyperlink{S18}{S18} \\
Plated designs:
enrichment and
ordering for every
metric &
\protect\hyperlink{S23}{S23},
\protect\hyperlink{S24}{S24},
\protect\hyperlink{S22}{S22} \\
The reference rows
(resolution, length,
tyrosine fraction,
distance) &
\protect\hyperlink{S12}{S12} \\
Experiments that
could not be made &
\protect\hyperlink{S48}{S48},
\protect\hyperlink{S49}{S49} \\
\end{longtable}

}

\section{4 Table of
contents}

{\footnotesize

\begin{longtable}[]{@{}
  >{\raggedright\arraybackslash}p{(\columnwidth - 6\tabcolsep) * \real{0.0818}}
  >{\raggedright\arraybackslash}p{(\columnwidth - 6\tabcolsep) * \real{0.5636}}
  >{\raggedright\arraybackslash}p{(\columnwidth - 6\tabcolsep) * \real{0.2182}}
  >{\raggedright\arraybackslash}p{(\columnwidth - 6\tabcolsep) * \real{0.1364}}@{}}
\toprule\noalign{}
\begin{minipage}[b]{\linewidth}\raggedright
Table
\end{minipage} &
\begin{minipage}[b]{\linewidth}\raggedright
Title
\end{minipage} &
\begin{minipage}[b]{\linewidth}\raggedright
Topic
\end{minipage} &
\begin{minipage}[b]{\linewidth}\raggedright
Rows × cols
\end{minipage} \\
\midrule\noalign{}
\endhead
\bottomrule\noalign{}
\endlastfoot
\protect\hyperlink{S1}{S1}
& Main metric set:
the role of every
distinct metric, and
why & Metric audit &
191 × 16 \\
\protect\hyperlink{S2}{S2}
& Metric contracts:
what each named
quantity reads, over
which residues, and
where it is valid &
Metric audit & 216 ×
19 \\
\protect\hyperlink{S3}{S3}
& Applicability of
every metric in
every test (long
format) & Metric
audit & 6,480 ×
11 \\
\protect\hyperlink{S4}{S4}
& Applicability by
class of metric and
test & Metric audit
& 340 × 5 \\
\protect\hyperlink{S5}{S5}
& Counts restated on
the eligible metrics
(denominators before
and after the audit)
& Metric audit & 58
× 14 \\
\protect\hyperlink{S6}{S6}
& Test inventory:
the 30 tests, with
comparator, control
type, role and
caveats & Metric
audit & 30 × 14 \\
\protect\hyperlink{S7}{S7}
& Wider set of 49
zymogen-mature
pairs: all 136
scored metrics, with
the same-state null
& Activity
discrimination (3.1)
& 136 × 25 \\
\protect\hyperlink{S8}{S8}
& Zymogen pairs
re-predicted with
Chai-1, restated
with the definedness
gate & Activity
discrimination (3.1)
& 33 × 22 \\
\protect\hyperlink{S9}{S9}
& Definitions of the
pair sets & Activity
discrimination (3.1)
& 6 × 4 \\
\protect\hyperlink{S10}{S10}
& Substrate-trapping
mutants checked
against the
literature (tier
retired) & Activity
discrimination (3.1)
& 6 × 4 \\
\protect\hyperlink{S11}{S11}
& Zymogen-mature
evaluation set: all
158 metrics ranked
by AUROC (complete
Table 5) & Activity
discrimination (3.1)
& 158 × 27 \\
\protect\hyperlink{S12}{S12}
& Reference rows in
full, across the
activity tests &
Activity
discrimination (3.1)
& 35 × 9 \\
\protect\hyperlink{S13}{S13}
& Plated designs:
all 128 metrics
ranked by the AUROC
of active against no
active designs
(complete Table 6) &
Activity
discrimination (3.1)
& 128 × 29 \\
\protect\hyperlink{S14}{S14}
& Substrate swap on
all systems: all 33
prediction-based
quantities, natural
enzymes and de novo
designs & Substrate
discrimination (3.2)
& 242 × 17 \\
\protect\hyperlink{S15}{S15}
& Substrate swap
stratified by ligand
size and charge &
Substrate
discrimination (3.2)
& 1,136 × 27 \\
\protect\hyperlink{S16}{S16}
& BglB: every
analysed quantity
ranked by rank
correlation with the
measured impairment
(complete Table 8) &
Activity ranking
(3.3) & 114 × 35 \\
\protect\hyperlink{S17}{S17}
& BglB: all 84
structure-space and
comparator metrics &
Activity ranking
(3.3) & 84 × 29 \\
\protect\hyperlink{S18}{S18}
& BglB: sensitivity
analyses & Activity
ranking (3.3) & 11 ×
22 \\
\protect\hyperlink{S19}{S19}
& BglB:
prediction-based
metrics (Chai-1 with
the assay substrate)
& Activity ranking
(3.3) & 28 × 29 \\
\protect\hyperlink{S20}{S20}
& Plated designs:
hits per 96-well
plate for every
metric (filter-style
analysis) & Activity
ranking (3.3) & 128
× 37 \\
\protect\hyperlink{S21}{S21}
& Plated designs:
every metric ranked
by correlation with
kcat/KM among the 16
designs that have a
value (complete
Table 9) & Activity
ranking (3.3) & 106
× 29 \\
\protect\hyperlink{S22}{S22}
& Plated designs:
ordering among the
16 active designs,
structure-space
metrics and
comparators &
Activity ranking
(3.3) & 81 × 23 \\
\protect\hyperlink{S23}{S23}
& Plated designs:
enrichment at every
keep fraction,
structure-space
metrics & Activity
ranking (3.3) & 816
× 27 \\
\protect\hyperlink{S24}{S24}
& Plated designs:
enrichment for
prediction-based
metrics & Activity
ranking (3.3) & 232
× 27 \\
\protect\hyperlink{S25}{S25}
& Sanity floors by
class of metric &
Detection ability
(3.4) & 53 × 8 \\
\protect\hyperlink{S26}{S26}
& Control matching:
the loosest tier
needed & Detection
ability (3.4) & 81 ×
6 \\
\protect\hyperlink{S27}{S27}
& Control matching:
balance of residues
changed in the two
arms & Detection
ability (3.4) & 15 ×
9 \\
\protect\hyperlink{S28}{S28}
& Catalytic lesion:
every metric scored
(Table 10 plus
supplementary
comparators and
per-step flags) &
Detection ability
(3.4) & 43 × 73 \\
\protect\hyperlink{S29}{S29}
& Structure-space
specificity by
metric family and
lesion step, with
the panel statistic
and sensitivity sets
& Detection ability
(3.4) & 20 × 27 \\
\protect\hyperlink{S30}{S30}
& Re-predicted
catalytic lesion:
the prediction-based
metrics under three
successive controls
and the
pre-specified rule &
Detection ability
(3.4) & 6 × 23 \\
\protect\hyperlink{S31}{S31}
& Sensitivity of the
specificity-ratio
counts to motif
size, packing
matching and count
gating & Detection
ability (3.4) & 20 ×
7 \\
\protect\hyperlink{S32}{S32}
& Re-predicted
ladder: seed-noise
context & Detection
ability (3.4) & 56 ×
12 \\
\protect\hyperlink{S33}{S33}
& Specificity ratio
of every eligible
metric at every step
& Detection ability
(3.4) & 256 × 29 \\
\protect\hyperlink{S34}{S34}
& Specificity ratio
of whole-protein and
sequence-only
comparators at every
step & Detection
ability (3.4) & 308
× 29 \\
\protect\hyperlink{S35}{S35}
& The specific cells
that survive the
controls & Detection
ability (3.4) & 97 ×
10 \\
\protect\hyperlink{S36}{S36}
& Rank statistic R
per lesion step on
the 143 main enzymes
& Detection ability
(3.4) & 313 × 14 \\
\protect\hyperlink{S37}{S37}
& Re-predicted
ladder: distance of
the substituted and
control residues to
the ligand &
Detection ability
(3.4) & 312 × 34 \\
\protect\hyperlink{S38}{S38}
& Re-predicted
ladder: candidates
for a new
distance-matched
control & Detection
ability (3.4) & 291
× 17 \\
\protect\hyperlink{S39}{S39}
& Re-predicted
ladder: specificity
ratio in strata of
the distance to the
ligand & Detection
ability (3.4) & 203
× 27 \\
\protect\hyperlink{S40}{S40}
& Re-predicted
ladder:
caliper-matched
subsets and the rule
labels & Detection
ability (3.4) & 348
× 25 \\
\protect\hyperlink{S41}{S41}
& Re-predicted
ladder:
regression-adjusted
ratios & Detection
ability (3.4) & 116
× 28 \\
\protect\hyperlink{S42}{S42}
& Panel equivalence
statistic: earlier
whole panel against
the main-set metrics
& Detection ability
(3.4) & 6 × 10 \\
\protect\hyperlink{S43}{S43}
& Re-predicted
ladder: specificity
ratio per step for
all metrics &
Detection ability
(3.4) & 231 × 16 \\
\protect\hyperlink{S44}{S44}
& Second-shell
lesion: every metric
scored (Table 11
plus supplementary
comparators and
per-dose flags) &
Detection ability
(3.4) & 37 × 59 \\
\protect\hyperlink{S45}{S45}
& Drift of the
current Chai-1
engine against the
earlier predictions
& Substrate
discrimination (3.2)
& 176 × 5 \\
\protect\hyperlink{S46}{S46}
& Effective
dimensionality of
the panel &
Detection ability
(3.4) & 4 × 10 \\
\protect\hyperlink{S47}{S47}
& De novo designs:
burial equivalence
of the created-site
control & Detection
ability (3.4) & 5 ×
11 \\
\protect\hyperlink{S48}{S48}
& Experiments that
could not be made or
were retired &
Detection ability
(3.4) & 9 × 4 \\
\protect\hyperlink{S49}{S49}
& PLACER ensemble
metrics on the
isosteric step (not
interpretable) &
Detection ability
(3.4) & 50 × 12 \\
\end{longtable}

}

\section{5
Supplementary
tables}

\hypertarget{S1}{%
\subsubsection{Table
S1. Main metric set:
the role of every
distinct metric, and
why}\label{S1}}

\begin{shaded}\small

\textbf{File.}
\texttt{supplementary/S1\_main\_metric\_set.csv}
· 191 rows × 16
columns · public

\textbf{Supports.}
Methods, Table 1 and
`Which metrics are
reported'; the main
set of 33 metrics
used in Tables 5 to
11. \textbf{Cited
in.} Methods 2.1;
Discussion.

\textbf{A row is}
one distinct metric
(a canonical key;
literal aliases are
merged), 191 in all.

\end{shaded}

The audit that fixes
which metrics the
paper reports. For
each distinct metric
it gives the
applicability status
in the detection
test of the
catalytic lesion, in
the
dead-against-active
tests and in the
activity-ranking
tests, and the
resulting role: main
(33, valid and
tested in the
deliberate-change
tests and in the
tests on real
proteins), reference
(6, comparators
chosen before any
result was seen),
excluded\_bookkeeping
(63 counters and
constants) or one of
the supplement roles
(supp\_global 44,
supp\_placer 25,
supp\_excluded 10,
supp\_aggregate 4,
supp\_partial 3,
supp\_real\_protein\_tests\_only
2, and supp\_derived
1, the Chai-1
combined score, a
rescaled ipTM). The
reason for each role
is in the last
column. The status
is decided from what
a metric reads,
whether it is
defined and the
design of the
experiment, never
from a result.

\textbf{Key
columns.}

{\small

\begin{longtable}[]{@{}
  >{\raggedright\arraybackslash}p{(\columnwidth - 2\tabcolsep) * \real{0.2909}}
  >{\raggedright\arraybackslash}p{(\columnwidth - 2\tabcolsep) * \real{0.7091}}@{}}
\toprule\noalign{}
\begin{minipage}[b]{\linewidth}\raggedright
Column
\end{minipage} &
\begin{minipage}[b]{\linewidth}\raggedright
Meaning
\end{minipage} \\
\midrule\noalign{}
\endhead
\bottomrule\noalign{}
\endlastfoot
\texttt{canonical\_key}
/ \texttt{members} /
\texttt{rep\_metric}
& the distinct
metric, the names
merged into it and
the representative
one \\
\texttt{mode} /
\texttt{input\_class}
/ \texttt{subfamily}
/
\texttt{aggregator}
& structure-space or
prediction-based,
what the metric
reads, its family
and how it is
summarised \\
\texttt{status\_detection\_test}
/
\texttt{fraction\_detection\_test}
& applicability
status in the
detection test of
the catalytic lesion
and the share of
units on which the
metric is defined \\
\texttt{status\_dead\_vs\_active\_tests}
/
\texttt{fraction\_dead\_vs\_active\_tests}
/
\texttt{partial\_dead\_vs\_active\_tests}
& the same for the
dead-against-active
tests \\
\texttt{status\_ranking\_tests}
/
\texttt{ranking\_tests\_ok}
& status in each
activity-ranking
test (BglB variants,
plated designs) and
the tests in which
the metric is OK \\
\texttt{role} &
main, reference,
excluded\_bookkeeping,
supp\_global,
supp\_placer,
supp\_excluded,
supp\_aggregate,
supp\_partial,
supp\_real\_protein\_tests\_only
or supp\_derived \\
\texttt{reason} &
why the metric has
that role \\
\end{longtable}

}

\textbf{Excerpt}
(first 3 rows where
role = main).

{\footnotesize

\begin{longtable}[]{@{}
  >{\raggedright\arraybackslash}p{(\columnwidth - 8\tabcolsep) * \real{0.2000}}
  >{\raggedright\arraybackslash}p{(\columnwidth - 8\tabcolsep) * \real{0.2000}}
  >{\raggedright\arraybackslash}p{(\columnwidth - 8\tabcolsep) * \real{0.2000}}
  >{\raggedright\arraybackslash}p{(\columnwidth - 8\tabcolsep) * \real{0.2000}}
  >{\raggedright\arraybackslash}p{(\columnwidth - 8\tabcolsep) * \real{0.2000}}@{}}
\toprule\noalign{}
\begin{minipage}[b]{\linewidth}\raggedright
\texttt{canonical\_key}
\end{minipage} &
\begin{minipage}[b]{\linewidth}\raggedright
\texttt{input\_class}
\end{minipage} &
\begin{minipage}[b]{\linewidth}\raggedright
\texttt{status\_detection\_test}
\end{minipage} &
\begin{minipage}[b]{\linewidth}\raggedright
\texttt{status\_dead\_vs\_active\_tests}
\end{minipage} &
\begin{minipage}[b]{\linewidth}\raggedright
\texttt{role}
\end{minipage} \\
\midrule\noalign{}
\endhead
\bottomrule\noalign{}
\endlastfoot
\texttt{ame\_crystal\_clash\_max}
&
\texttt{pred\_interface}
&
\texttt{OK-PROVISIONAL}
& \texttt{OK} &
\texttt{main} \\
\texttt{ame\_crystal\_rmsd}
& \texttt{ame} &
\texttt{OK-KNOCKON}
& \texttt{OK} &
\texttt{main} \\
\texttt{catalytic\_frac\_buried}
&
\texttt{struct\_site}
& \texttt{OK} &
\texttt{OK} &
\texttt{main} \\
\end{longtable}

}

\textbf{Notes.}

\begin{itemize}
\tightlist
\item
  The status
  vocabulary is in
  the Conventions.
\item
  NOT-RUN in
  status\_ranking\_tests
  marks the PLACER
  metrics, which
  were not computed
  in the BglB and
  plated-design
  analyses.
\end{itemize}

\hypertarget{S2}{%
\subsubsection{Table
S2. Metric
contracts: what each
named quantity
reads, over which
residues, and where
it is
valid}\label{S2}}

\begin{shaded}\small

\textbf{File.}
\texttt{supplementary/S2\_metric\_contracts.csv}
· 216 rows × 19
columns · public

\textbf{Supports.}
Methods, Table 1 and
`Which metrics are
reported'.
\textbf{Cited in.}
Methods 2.1.

\textbf{A row is}
one named quantity
(216 in all; an
alias is a separate
row linked by
alias\_of and
canonical\_key).

\end{shaded}

The input contract
of every quantity in
the panel: what it
reads (sequence
only, whole-protein
structure,
site-scoped
structure, predictor
output, active-site
accuracy, PLACER
output, deposit
metadata or
bookkeeping), over
which residues it is
scored in the
catalytic arm and in
the two kinds of
control arm, which
inputs it needs,
which kinds of
metric it is valid
for and which
changes it cannot
respond to.

\textbf{Key
columns.}

{\small

\begin{longtable}[]{@{}
  >{\raggedright\arraybackslash}p{(\columnwidth - 2\tabcolsep) * \real{0.2909}}
  >{\raggedright\arraybackslash}p{(\columnwidth - 2\tabcolsep) * \real{0.7091}}@{}}
\toprule\noalign{}
\begin{minipage}[b]{\linewidth}\raggedright
Column
\end{minipage} &
\begin{minipage}[b]{\linewidth}\raggedright
Meaning
\end{minipage} \\
\midrule\noalign{}
\endhead
\bottomrule\noalign{}
\endlastfoot
\texttt{metric\_name}
/
\texttt{canonical\_key}
/ \texttt{alias\_of}
/
\texttt{alias\_kind}
& the quantity, the
distinct metric
after aliases are
merged, and the kind
of alias \\
\texttt{aggregator}
/
\texttt{metric\_family}
/ \texttt{subfamily}
& how it is
summarised, and its
family and
subfamily \\
\texttt{input\_class}
& sequence-only,
whole-protein
structure,
site-scoped
structure,
prediction,
active-site accuracy
(AME), PLACER,
deposit metadata or
bookkeeping \\
\texttt{scope\_catalytic\_arm}
/
\texttt{scope\_control\_arm\_catalytic\_lesion}
/
\texttt{scope\_control\_arm\_second\_shell}
& the residues the
metric is scored
over in the
catalytic arm, in
the control arm of
the catalytic lesion
and in the control
arm of the
second-shell
lesion \\
\texttt{needs\_cat\_residue}
/
\texttt{needs\_ligand}
/
\texttt{needs\_reference}
/
\texttt{needs\_predictor}
& the inputs the
metric needs \\
\texttt{valid\_modes}
& structure-space,
prediction-based,
PLACER, any or
deposit metadata \\
\texttt{invariant\_under}
& changes the metric
cannot register (for
example sequence
composition under a
rotamer change) \\
\texttt{ame\_flavour}
/ \texttt{note} &
for active-site
accuracy, the
reference it is
measured against;
remarks on the
metric, including
where the metric as
computed differs
from its earlier
description \\
\end{longtable}

}

\textbf{Excerpt}
(first 3 rows where
input\_class =
struct\_site).

{\footnotesize

\begin{longtable}[]{@{}
  >{\raggedright\arraybackslash}p{(\columnwidth - 6\tabcolsep) * \real{0.2500}}
  >{\raggedright\arraybackslash}p{(\columnwidth - 6\tabcolsep) * \real{0.2500}}
  >{\raggedright\arraybackslash}p{(\columnwidth - 6\tabcolsep) * \real{0.2500}}
  >{\raggedright\arraybackslash}p{(\columnwidth - 6\tabcolsep) * \real{0.2500}}@{}}
\toprule\noalign{}
\begin{minipage}[b]{\linewidth}\raggedright
\texttt{metric\_name}
\end{minipage} &
\begin{minipage}[b]{\linewidth}\raggedright
\texttt{input\_class}
\end{minipage} &
\begin{minipage}[b]{\linewidth}\raggedright
\texttt{scope\_catalytic\_arm}
\end{minipage} &
\begin{minipage}[b]{\linewidth}\raggedright
\texttt{valid\_modes}
\end{minipage} \\
\midrule\noalign{}
\endhead
\bottomrule\noalign{}
\endlastfoot
\texttt{catalytic\_frac\_buried}
&
\texttt{struct\_site}
&
\texttt{scored\_residue\_set}
&
\texttt{structure-space} \\
\texttt{catalytic\_ligand\_min\_dist}
&
\texttt{struct\_site}
&
\texttt{scored\_residue\_set+ligand\_within\_…}
&
\texttt{structure-space} \\
\texttt{catalytic\_pairwise\_dist\_max}
&
\texttt{struct\_site}
&
\texttt{scored\_residue\_set}
&
\texttt{structure-space} \\
\end{longtable}

}

\hypertarget{S3}{%
\subsubsection{Table
S3. Applicability of
every metric in
every test (long
format)}\label{S3}}

\begin{shaded}\small

\textbf{File.}
\texttt{supplementary/S3\_applicability.csv}
· 6,480 rows × 11
columns · public

\textbf{Supports.}
Methods, `Which
metrics are
reported'; every
Results subsection.
\textbf{Cited in.}
Methods 2.1.

\textbf{A row is}
one named quantity
in one test (216
quantities x 30
tests = 6,480 rows).

\end{shaded}

Whether each
quantity can be
tested in each of
the 30 tests, and if
not why. Status
counts: OK 468,
GLOBAL 339, REF 15,
OK-KNOCKON 6,
OK-PROVISIONAL 10
(testable); X-ARM
2,296, X-BOOKKEEP
1,920, X-NOINPUT
576, X-WITHDRAWN
381, X-UNDEF-COV
203, X-NOCTRL 118,
X-CONFOUND 10,
NOT-RUN 79, DUP 59
(not testable, with
the reason). The
reason and the
coverage (defined\_k
of defined\_N units)
are given. Decided
from inputs,
definedness and
design only, never
from results.

\textbf{Key
columns.}

{\small

\begin{longtable}[]{@{}
  >{\raggedright\arraybackslash}p{(\columnwidth - 2\tabcolsep) * \real{0.2909}}
  >{\raggedright\arraybackslash}p{(\columnwidth - 2\tabcolsep) * \real{0.7091}}@{}}
\toprule\noalign{}
\begin{minipage}[b]{\linewidth}\raggedright
Column
\end{minipage} &
\begin{minipage}[b]{\linewidth}\raggedright
Meaning
\end{minipage} \\
\midrule\noalign{}
\endhead
\bottomrule\noalign{}
\endlastfoot
\texttt{metric\_name}
/
\texttt{canonical\_key}
/
\texttt{input\_class}
& the quantity, the
distinct metric and
its class \\
\texttt{test} /
\texttt{test\_class}
& the test (its name
and class as in
\protect\hyperlink{S6}{S6})
and the class alone,
for example `dead
against active,
prediction-based' \\
\texttt{status} /
\texttt{reason} &
the applicability
status and the
reason in words \\
\texttt{defined\_k}
/
\texttt{defined\_N}
/ \texttt{fraction}
/ \texttt{partial} &
number of units on
which the metric is
defined, of how
many, their ratio,
and whether coverage
is partial \\
\end{longtable}

}

\textbf{Excerpt}
(first 3 rows where
status = OK).

{\footnotesize

\begin{longtable}[]{@{}
  >{\raggedright\arraybackslash}p{(\columnwidth - 6\tabcolsep) * \real{0.2500}}
  >{\raggedright\arraybackslash}p{(\columnwidth - 6\tabcolsep) * \real{0.2500}}
  >{\raggedright\arraybackslash}p{(\columnwidth - 6\tabcolsep) * \real{0.2500}}
  >{\raggedright\arraybackslash}p{(\columnwidth - 6\tabcolsep) * \real{0.2500}}@{}}
\toprule\noalign{}
\begin{minipage}[b]{\linewidth}\raggedright
\texttt{metric\_name}
\end{minipage} &
\begin{minipage}[b]{\linewidth}\raggedright
\texttt{test}
\end{minipage} &
\begin{minipage}[b]{\linewidth}\raggedright
\texttt{status}
\end{minipage} &
\begin{minipage}[b]{\linewidth}\raggedright
\texttt{reason}
\end{minipage} \\
\midrule\noalign{}
\endhead
\bottomrule\noalign{}
\endlastfoot
\texttt{ame\_crystal\_clash\_max}
&
\texttt{BglB\ single-point\ variants\ (predi…}
& \texttt{OK} &
\texttt{eligible\ and\ scored} \\
\texttt{ame\_crystal\_clash\_max}
&
\texttt{zymogen-mature\ pairs\ re-predicted…}
& \texttt{OK} &
\texttt{eligible\ and\ scored} \\
\texttt{ame\_crystal\_clash\_max}
&
\texttt{substrate\ swap,\ natural\ enzymes\ {[}…}
& \texttt{OK} &
\texttt{eligible\ and\ scored} \\
\end{longtable}

}

\textbf{Notes.}

\begin{itemize}
\tightlist
\item
  Released as a CSV
  only: with 6,480
  rows it is not
  printed.
  \protect\hyperlink{S25}{S25}
  is its summary by
  class of metric.
\item
  Of the 79 cells
  with status
  NOT-RUN, 50 are
  PLACER metrics
  that were not
  computed in the
  BglB and
  plated-design
  analyses. The
  other 29 are
  prediction-based
  metrics in the
  ladder with
  ligand-distance-matched
  controls: this
  analysis was
  carried out
  (results in
  \protect\hyperlink{S40}{S40}
  and
  \protect\hyperlink{S41}{S41})
  but its cells keep
  the label NOT-RUN.
\end{itemize}

\hypertarget{S4}{%
\subsubsection{Table
S4. Applicability by
class of metric and
test}\label{S4}}

\begin{shaded}\small

\textbf{File.}
\texttt{supplementary/S4\_applicability\_by\_class.csv}
· 340 rows × 5
columns · public

\textbf{Supports.}
Methods, `Which
metrics are
reported' (summary
of
\protect\hyperlink{S3}{S3}).
\textbf{Cited in.}
Methods 2.1.

\textbf{A row is}
one class of metric
in one test with one
status (340 rows).

\end{shaded}

A summary of
\protect\hyperlink{S3}{S3}:
for each input class
and test, how many
metrics receive each
status. The quickest
way to see why a
whole class of
metrics is absent
from a result.

\textbf{Key
columns.}

{\small

\begin{longtable}[]{@{}
  >{\raggedright\arraybackslash}p{(\columnwidth - 2\tabcolsep) * \real{0.2909}}
  >{\raggedright\arraybackslash}p{(\columnwidth - 2\tabcolsep) * \real{0.7091}}@{}}
\toprule\noalign{}
\begin{minipage}[b]{\linewidth}\raggedright
Column
\end{minipage} &
\begin{minipage}[b]{\linewidth}\raggedright
Meaning
\end{minipage} \\
\midrule\noalign{}
\endhead
\bottomrule\noalign{}
\endlastfoot
\texttt{input\_class}
& class of metric
(ame, bookkeeping,
pred\_global,
pred\_interface,
seq\_only,
struct\_site,
struct\_whole,
\ldots) \\
\texttt{test} /
\texttt{test\_class}
& the test and its
class \\
\texttt{status} &
applicability
status \\
\texttt{n\_metrics}
& number of metrics
of that class with
that status in that
test \\
\end{longtable}

}

\textbf{Excerpt}
(first 3 rows where
status = OK).

{\footnotesize

\begin{longtable}[]{@{}
  >{\raggedright\arraybackslash}p{(\columnwidth - 6\tabcolsep) * \real{0.2500}}
  >{\raggedright\arraybackslash}p{(\columnwidth - 6\tabcolsep) * \real{0.2500}}
  >{\raggedright\arraybackslash}p{(\columnwidth - 6\tabcolsep) * \real{0.2500}}
  >{\raggedright\arraybackslash}p{(\columnwidth - 6\tabcolsep) * \real{0.2500}}@{}}
\toprule\noalign{}
\begin{minipage}[b]{\linewidth}\raggedright
\texttt{input\_class}
\end{minipage} &
\begin{minipage}[b]{\linewidth}\raggedright
\texttt{test\_class}
\end{minipage} &
\begin{minipage}[b]{\linewidth}\raggedright
\texttt{status}
\end{minipage} &
\begin{minipage}[b]{\linewidth}\raggedright
\texttt{n\_metrics}
\end{minipage} \\
\midrule\noalign{}
\endhead
\bottomrule\noalign{}
\endlastfoot
\texttt{ame} &
\texttt{activity\ ranking,\ prediction-based}
& \texttt{OK} &
\texttt{5} \\
\texttt{ame} &
\texttt{dead\ against\ active,\ prediction-b…}
& \texttt{OK} &
\texttt{6} \\
\texttt{ame} &
\texttt{substrate\ swap}
& \texttt{OK} &
\texttt{6} \\
\end{longtable}

}

\hypertarget{S5}{%
\subsubsection{Table
S5. Counts restated
on the eligible
metrics
(denominators before
and after the
audit)}\label{S5}}

\begin{shaded}\small

\textbf{File.}
\texttt{supplementary/S5\_denominator\_restatement.csv}
· 58 rows × 14
columns · public

\textbf{Supports.}
Discussion (`Test a
metric only where it
is defined');
Methods, `Which
metrics are
reported'.
\textbf{Cited in.}
Methods 2.1.

\textbf{A row is}
one headline count
of the earlier
whole-panel analysis
(58 in all).

\end{shaded}

The counts of
responding and
specific metrics of
the earlier analysis
(original\_k of
original\_N, over
every metric with a
testable cell)
restated on the
eligible metrics
only (eligible\_k of
eligible\_N, also as
distinct metrics and
for the
whole-protein
comparators), with
the breakdown of
what was excluded
and why. It shows
how much of the
earlier denominators
were metrics that
could not respond.

\textbf{Key
columns.}

{\small

\begin{longtable}[]{@{}
  >{\raggedright\arraybackslash}p{(\columnwidth - 2\tabcolsep) * \real{0.2909}}
  >{\raggedright\arraybackslash}p{(\columnwidth - 2\tabcolsep) * \real{0.7091}}@{}}
\toprule\noalign{}
\begin{minipage}[b]{\linewidth}\raggedright
Column
\end{minipage} &
\begin{minipage}[b]{\linewidth}\raggedright
Meaning
\end{minipage} \\
\midrule\noalign{}
\endhead
\bottomrule\noalign{}
\endlastfoot
\texttt{analysis} /
\texttt{level} /
\texttt{statistic} &
the earlier analysis
(data set and
lesion), the lesion
step and the
statistic
(`responds' or
`specific') \\
\texttt{original\_k}
/
\texttt{original\_N}
& count and
denominator of the
earlier analysis \\
\texttt{eligible\_k}
/
\texttt{eligible\_N}
/
\texttt{eligible\_distinct\_k}
/
\texttt{eligible\_distinct\_N}
& the same on
eligible metrics,
and as distinct
metrics \\
\texttt{global\_k} /
\texttt{global\_N} &
the count among
whole-protein
comparators \\
\texttt{excluded\_breakdown}
/
\texttt{recompute\_kind}
/ \texttt{note} &
which statuses were
excluded and how
many, how the
recount was done,
and a note \\
\end{longtable}

}

\textbf{Excerpt}
(first 3 rows).

{\footnotesize

\begin{longtable}[]{@{}
  >{\raggedright\arraybackslash}p{(\columnwidth - 12\tabcolsep) * \real{0.1429}}
  >{\raggedright\arraybackslash}p{(\columnwidth - 12\tabcolsep) * \real{0.1429}}
  >{\raggedright\arraybackslash}p{(\columnwidth - 12\tabcolsep) * \real{0.1429}}
  >{\raggedright\arraybackslash}p{(\columnwidth - 12\tabcolsep) * \real{0.1429}}
  >{\raggedright\arraybackslash}p{(\columnwidth - 12\tabcolsep) * \real{0.1429}}
  >{\raggedright\arraybackslash}p{(\columnwidth - 12\tabcolsep) * \real{0.1429}}
  >{\raggedright\arraybackslash}p{(\columnwidth - 12\tabcolsep) * \real{0.1429}}@{}}
\toprule\noalign{}
\begin{minipage}[b]{\linewidth}\raggedright
\texttt{analysis}
\end{minipage} &
\begin{minipage}[b]{\linewidth}\raggedright
\texttt{level}
\end{minipage} &
\begin{minipage}[b]{\linewidth}\raggedright
\texttt{statistic}
\end{minipage} &
\begin{minipage}[b]{\linewidth}\raggedright
\texttt{original\_k}
\end{minipage} &
\begin{minipage}[b]{\linewidth}\raggedright
\texttt{original\_N}
\end{minipage} &
\begin{minipage}[b]{\linewidth}\raggedright
\texttt{eligible\_k}
\end{minipage} &
\begin{minipage}[b]{\linewidth}\raggedright
\texttt{eligible\_N}
\end{minipage} \\
\midrule\noalign{}
\endhead
\bottomrule\noalign{}
\endlastfoot
\texttt{main\ set,\ packing-matched\ control…}
& \texttt{ala} &
\texttt{responds} &
\texttt{55} &
\texttt{102} &
\texttt{25} &
\texttt{37} \\
\texttt{main\ set,\ packing-matched\ control…}
& \texttt{ala} &
\texttt{specific} &
\texttt{11} &
\texttt{102} &
\texttt{10} &
\texttt{37} \\
\texttt{main\ set,\ packing-matched\ control…}
& \texttt{gly} &
\texttt{responds} &
\texttt{58} &
\texttt{98} &
\texttt{25} &
\texttt{34} \\
\end{longtable}

}

\hypertarget{S6}{%
\subsubsection{Table
S6. Test inventory:
the 30 tests, with
comparator, control
type, role and
caveats}\label{S6}}

\begin{shaded}\small

\textbf{File.}
\texttt{supplementary/S6\_test\_inventory.csv}
· 30 rows × 14
columns · public

\textbf{Supports.}
Methods 2.2 to 2.5;
the key to the test
names in
\protect\hyperlink{S3}{S3}
and
\protect\hyperlink{S4}{S4}.
\textbf{Cited in.}
Methods 2.1;
Discussion.

\textbf{A row is}
one test: an
experiment on a data
set with one kind of
metric (30 in all).

\end{shaded}

Every test with its
name and class, the
kind of metric
(structure-space or
prediction-based),
the unit and the
number of units, the
resampling unit, the
changes applied, the
comparator and
control type, the
role (primary,
sensitivity, sanity,
retired,
not\_measurable or
dup) and the
caveats.
\protect\hyperlink{S3}{S3}
and
\protect\hyperlink{S4}{S4}
use the test names
of this table.

\textbf{Key
columns.}

{\small

\begin{longtable}[]{@{}
  >{\raggedright\arraybackslash}p{(\columnwidth - 2\tabcolsep) * \real{0.2909}}
  >{\raggedright\arraybackslash}p{(\columnwidth - 2\tabcolsep) * \real{0.7091}}@{}}
\toprule\noalign{}
\begin{minipage}[b]{\linewidth}\raggedright
Column
\end{minipage} &
\begin{minipage}[b]{\linewidth}\raggedright
Meaning
\end{minipage} \\
\midrule\noalign{}
\endhead
\bottomrule\noalign{}
\endlastfoot
\texttt{test} /
\texttt{reader\_name}
/
\texttt{test\_class}
/ \texttt{mode} &
the test (name with
its class), the name
alone, the class,
and structure-space
or
prediction-based \\
\texttt{unit} /
\texttt{n\_units} /
\texttt{n\_clusters}
/
\texttt{cluster\_def}
& what is counted,
how many, and the
resampling unit \\
\texttt{levels} /
\texttt{comparator}
/
\texttt{control\_type}
& what is changed or
contrasted, what the
metric is compared
with and how the
control is built \\
\texttt{claim\_type}
/ \texttt{role} &
the kind of claim
and the role:
primary,
sensitivity, sanity,
retired,
not\_measurable or
dup \\
\texttt{caveats} &
limits of the
test \\
\end{longtable}

}

\textbf{Excerpt}
(first 3 rows).

{\footnotesize

\begin{longtable}[]{@{}
  >{\raggedright\arraybackslash}p{(\columnwidth - 8\tabcolsep) * \real{0.2000}}
  >{\raggedright\arraybackslash}p{(\columnwidth - 8\tabcolsep) * \real{0.2000}}
  >{\raggedright\arraybackslash}p{(\columnwidth - 8\tabcolsep) * \real{0.2000}}
  >{\raggedright\arraybackslash}p{(\columnwidth - 8\tabcolsep) * \real{0.2000}}
  >{\raggedright\arraybackslash}p{(\columnwidth - 8\tabcolsep) * \real{0.2000}}@{}}
\toprule\noalign{}
\begin{minipage}[b]{\linewidth}\raggedright
\texttt{reader\_name}
\end{minipage} &
\begin{minipage}[b]{\linewidth}\raggedright
\texttt{test\_class}
\end{minipage} &
\begin{minipage}[b]{\linewidth}\raggedright
\texttt{mode}
\end{minipage} &
\begin{minipage}[b]{\linewidth}\raggedright
\texttt{claim\_type}
\end{minipage} &
\begin{minipage}[b]{\linewidth}\raggedright
\texttt{role}
\end{minipage} \\
\midrule\noalign{}
\endhead
\bottomrule\noalign{}
\endlastfoot
\texttt{main\ set,\ packing-matched\ controls}
&
\texttt{catalytic\ lesion,\ structure-space}
&
\texttt{structure-space}
&
\texttt{specificity}
&
\texttt{primary} \\
\texttt{sanity\ floor:\ all\ catalytic\ resid…}
&
\texttt{detection\ floor\ (duplicate\ of\ the…}
&
\texttt{structure-space}
& \texttt{detection}
& \texttt{dup} \\
\texttt{sanity\ floor:\ permuted\ sequence\ o…}
&
\texttt{sanity\ floor}
&
\texttt{structure-space}
& \texttt{detection}
& \texttt{sanity} \\
\end{longtable}

}

\hypertarget{S7}{%
\subsubsection{Table
S7. Wider set of 49
zymogen-mature
pairs: all 136
scored metrics, with
the same-state
null}\label{S7}}

\begin{shaded}\small

\textbf{File.}
\texttt{supplementary/S7\_zymogen\_all\_metrics.csv}
· 136 rows × 25
columns · public

\textbf{Supports.}
Results 3.1.1 (wider
set); Methods 2.2.1.
\textbf{Cited in.}
Methods 2.2.1;
Results 3.1.1;
Discussion.

\textbf{A row is}
one scored metric
(136 in all).

\end{shaded}

The analysis on the
49 pairs for which
structure-space
metrics were scored
(before the
evaluation set of 21
pairs that share a
ligand was fixed):
the direction-free
AUROC of dead
against active with
interval, the
directed AUROC, the
number of pairs, and
the AUROC of the
same metric on 96
same-state null
pairs, which shows
how much separation
arises without any
dead/active
difference.

\textbf{Key
columns.}

{\small

\begin{longtable}[]{@{}
  >{\raggedright\arraybackslash}p{(\columnwidth - 2\tabcolsep) * \real{0.2909}}
  >{\raggedright\arraybackslash}p{(\columnwidth - 2\tabcolsep) * \real{0.7091}}@{}}
\toprule\noalign{}
\begin{minipage}[b]{\linewidth}\raggedright
Column
\end{minipage} &
\begin{minipage}[b]{\linewidth}\raggedright
Meaning
\end{minipage} \\
\midrule\noalign{}
\endhead
\bottomrule\noalign{}
\endlastfoot
\texttt{metric\_name}
/
\texttt{canonical\_key}
/
\texttt{input\_class}
/
\texttt{in\_main\_set}
& the metric, its
class and whether it
is a main-set
metric \\
\texttt{n\_pairs} &
pairs with a value
(49 at most) \\
\texttt{auroc} /
\texttt{auroc\_lo} /
\texttt{auroc\_hi} /
\texttt{auroc\_directed}
& direction-free
AUROC with 95\%
interval over pairs,
and the directed
value \\
\texttt{auroc\_vs\_same\_state\_null}
/
\texttt{n\_zymogen\_pairs}
/
\texttt{n\_same\_state\_null\_pairs}
& AUROC of zymogens
against the
same-state null
pairs and the
numbers in each
group \\
\texttt{applicability\_status}
& applicability
status \\
\end{longtable}

}

The shared block of
12 provenance
columns (see
Conventions) is also
present and is not
repeated above.

\textbf{Excerpt}
(first 3 rows).

{\footnotesize

\begin{longtable}[]{@{}
  >{\raggedright\arraybackslash}p{(\columnwidth - 10\tabcolsep) * \real{0.1667}}
  >{\raggedright\arraybackslash}p{(\columnwidth - 10\tabcolsep) * \real{0.1667}}
  >{\raggedright\arraybackslash}p{(\columnwidth - 10\tabcolsep) * \real{0.1667}}
  >{\raggedright\arraybackslash}p{(\columnwidth - 10\tabcolsep) * \real{0.1667}}
  >{\raggedright\arraybackslash}p{(\columnwidth - 10\tabcolsep) * \real{0.1667}}
  >{\raggedright\arraybackslash}p{(\columnwidth - 10\tabcolsep) * \real{0.1667}}@{}}
\toprule\noalign{}
\begin{minipage}[b]{\linewidth}\raggedright
\texttt{metric\_name}
\end{minipage} &
\begin{minipage}[b]{\linewidth}\raggedright
\texttt{in\_main\_set}
\end{minipage} &
\begin{minipage}[b]{\linewidth}\raggedright
\texttt{n\_pairs}
\end{minipage} &
\begin{minipage}[b]{\linewidth}\raggedright
\texttt{auroc}
\end{minipage} &
\begin{minipage}[b]{\linewidth}\raggedright
\texttt{auroc\_lo}
\end{minipage} &
\begin{minipage}[b]{\linewidth}\raggedright
\texttt{auroc\_hi}
\end{minipage} \\
\midrule\noalign{}
\endhead
\bottomrule\noalign{}
\endlastfoot
\texttt{catalytic\_frac\_buried}
& \texttt{1} &
\texttt{49} &
\texttt{0.5321} &
\texttt{0.5019} &
\texttt{0.6112} \\
\texttt{catalytic\_hydration\_count}
& \texttt{0} &
\texttt{48} &
\texttt{0.5686} &
\texttt{0.505} &
\texttt{0.6604} \\
\texttt{catalytic\_ligand\_min\_dist}
& \texttt{0} &
\texttt{25} &
\texttt{0.5264} &
\texttt{0.5024} &
\texttt{0.6464} \\
\end{longtable}

}

\textbf{Notes.}

\begin{itemize}
\tightlist
\item
  Structure-space
  and PLACER metrics
  only; the 4
  prediction-based
  metrics were
  scored on the 28
  re-predicted pairs
  (\protect\hyperlink{S8}{S8}).
\item
  Results are
  descriptive of a
  wider set; the
  evaluation set is
  the 21 pairs of
  \protect\hyperlink{S36}{S36}.
\end{itemize}

\hypertarget{S8}{%
\subsubsection{Table
S8. Zymogen pairs
re-predicted with
Chai-1, restated
with the definedness
gate}\label{S8}}

\begin{shaded}\small

\textbf{File.}
\texttt{supplementary/S8\_zymogen\_re\_predicted\_restated.csv}
· 33 rows × 22
columns · public

\textbf{Supports.}
Results 3.1.1
(wider-set check for
the prediction-based
metrics).
\textbf{Cited in.}
Methods 2.2.1;
Results 3.1.1.

\textbf{A row is}
one prediction-based
metric (33 in all).

\end{shaded}

The 28
zymogen-mature pairs
re-predicted with
Chai-1: the AUROC of
each metric with
interval, the number
of pairs reported
against the number
on which the metric
is defined, and a
note where the
earlier value was
not interpretable.
Active-site accuracy
is relabelled as
measured against
each form's own
deposited structure.

\textbf{Key
columns.}

{\small

\begin{longtable}[]{@{}
  >{\raggedright\arraybackslash}p{(\columnwidth - 2\tabcolsep) * \real{0.2909}}
  >{\raggedright\arraybackslash}p{(\columnwidth - 2\tabcolsep) * \real{0.7091}}@{}}
\toprule\noalign{}
\begin{minipage}[b]{\linewidth}\raggedright
Column
\end{minipage} &
\begin{minipage}[b]{\linewidth}\raggedright
Meaning
\end{minipage} \\
\midrule\noalign{}
\endhead
\bottomrule\noalign{}
\endlastfoot
\texttt{metric\_name}
/
\texttt{metric\_reported\_as}
& the metric and the
name it carried in
the earlier
report \\
\texttt{input\_class}
& class of metric \\
\texttt{n\_pairs\_reported}
/
\texttt{n\_pairs\_defined\_restated}
& pairs in the
earlier report and
pairs on which the
metric is defined \\
\texttt{auroc} /
\texttt{auroc\_lo} /
\texttt{auroc\_hi} &
direction-free AUROC
with 95\%
interval \\
\texttt{applicability\_status}
/
\texttt{restatement\_note}
& applicability
status and what was
restated \\
\end{longtable}

}

The shared block of
12 provenance
columns (see
Conventions) is also
present and is not
repeated above.

\textbf{Excerpt}
(first 3 rows where
applicability\_status
= OK).

{\footnotesize

\begin{longtable}[]{@{}
  >{\raggedright\arraybackslash}p{(\columnwidth - 8\tabcolsep) * \real{0.2000}}
  >{\raggedright\arraybackslash}p{(\columnwidth - 8\tabcolsep) * \real{0.2000}}
  >{\raggedright\arraybackslash}p{(\columnwidth - 8\tabcolsep) * \real{0.2000}}
  >{\raggedright\arraybackslash}p{(\columnwidth - 8\tabcolsep) * \real{0.2000}}
  >{\raggedright\arraybackslash}p{(\columnwidth - 8\tabcolsep) * \real{0.2000}}@{}}
\toprule\noalign{}
\begin{minipage}[b]{\linewidth}\raggedright
\texttt{metric\_name}
\end{minipage} &
\begin{minipage}[b]{\linewidth}\raggedright
\texttt{n\_pairs\_reported}
\end{minipage} &
\begin{minipage}[b]{\linewidth}\raggedright
\texttt{n\_pairs\_defined\_restated}
\end{minipage} &
\begin{minipage}[b]{\linewidth}\raggedright
\texttt{auroc}
\end{minipage} &
\begin{minipage}[b]{\linewidth}\raggedright
\texttt{applicability\_status}
\end{minipage} \\
\midrule\noalign{}
\endhead
\bottomrule\noalign{}
\endlastfoot
\texttt{ame\_crystal\_clash\_max}
& \texttt{21} &
\texttt{21} &
\texttt{0.5351} &
\texttt{OK} \\
\texttt{ame\_crystal\_pass}
& \texttt{28} &
\texttt{28} &
\texttt{0.5045} &
\texttt{OK} \\
\texttt{ame\_crystal\_rmsd\_max}
& \texttt{28} &
\texttt{28} &
\texttt{0.5332} &
\texttt{OK} \\
\end{longtable}

}

\textbf{Notes.}

\begin{itemize}
\tightlist
\item
  28 pairs include
  the 7 apo pairs
  that are not in
  the 21-pair
  evaluation set.
\end{itemize}

\hypertarget{S9}{%
\subsubsection{Table
S9. Definitions of
the pair
sets}\label{S9}}

\begin{shaded}\small

\textbf{File.}
\texttt{supplementary/S9\_pair\_set\_definitions.csv}
· 6 rows × 4 columns
· public

\textbf{Supports.}
Methods 2.2.1;
Results 3.1.1.
\textbf{Cited in.}
Methods 2.2.1;
Results 3.1.1.

\textbf{A row is}
one set of
zymogen-mature pairs
(6 in all).

\end{shaded}

The six sets used in
3.1 with their size
and definition: 55
candidate rows, 49
scored pairs, 45
with PLACER, 28
re-predicted pairs,
the 21-pair
evaluation set
(pairs that share a
ligand) and the 96
same-state null
pairs, and which
table uses each.

\textbf{Key
columns.}

{\small

\begin{longtable}[]{@{}
  >{\raggedright\arraybackslash}p{(\columnwidth - 2\tabcolsep) * \real{0.2909}}
  >{\raggedright\arraybackslash}p{(\columnwidth - 2\tabcolsep) * \real{0.7091}}@{}}
\toprule\noalign{}
\begin{minipage}[b]{\linewidth}\raggedright
Column
\end{minipage} &
\begin{minipage}[b]{\linewidth}\raggedright
Meaning
\end{minipage} \\
\midrule\noalign{}
\endhead
\bottomrule\noalign{}
\endlastfoot
\texttt{pair\_set} /
\texttt{n} & the set
and its size \\
\texttt{definition}
& how the set is
defined \\
\texttt{used\_for} &
the table that uses
it \\
\end{longtable}

}

\textbf{Excerpt}
(first 3 rows).

{\footnotesize

\begin{longtable}[]{@{}
  >{\raggedright\arraybackslash}p{(\columnwidth - 4\tabcolsep) * \real{0.3333}}
  >{\raggedright\arraybackslash}p{(\columnwidth - 4\tabcolsep) * \real{0.3333}}
  >{\raggedright\arraybackslash}p{(\columnwidth - 4\tabcolsep) * \real{0.3333}}@{}}
\toprule\noalign{}
\begin{minipage}[b]{\linewidth}\raggedright
\texttt{pair\_set}
\end{minipage} &
\begin{minipage}[b]{\linewidth}\raggedright
\texttt{n}
\end{minipage} &
\begin{minipage}[b]{\linewidth}\raggedright
\texttt{used\_for}
\end{minipage} \\
\midrule\noalign{}
\endhead
\bottomrule\noalign{}
\endlastfoot
\texttt{zymogen\ candidate\ rows}
& \texttt{55} &
\texttt{candidate\ list} \\
\texttt{scored\ pairs\ (wider\ set)}
& \texttt{49} &
\texttt{S7\ (49\ =\ max\ n\_pairs)} \\
\texttt{pairs\ with\ PLACER}
& \texttt{45} &
\texttt{S49} \\
\end{longtable}

}

\hypertarget{S10}{%
\subsubsection{Table
S10.
Substrate-trapping
mutants checked
against the
literature (tier
retired)}\label{S10}}

\begin{shaded}\small

\textbf{File.}
\texttt{supplementary/S10\_trapping\_verification.csv}
· 6 rows × 4 columns
· public

\textbf{Supports.}
Results 3.1.1
(retired tier);
Methods 2.2.1.
\textbf{Cited in.}
Methods 2.2.1;
Results 3.1.1.

\textbf{A row is}
one verdict category
(6 rows).

\end{shaded}

The 30
trapping-mutant
pairs verified
against the primary
literature: 15
confirmed reduced, 4
confirmed inactive,
4 active and 7 not
stated; 26 of the 30
(86.7\%, 95\%
interval 70.3 to
94.7\%) are not
confirmed dead.
Every deposited
construct carries an
engineered
substitution; the
failure is the
inference from
substitution to dead
enzyme. The tier is
retired.

\textbf{Key
columns.}

{\small

\begin{longtable}[]{@{}
  >{\raggedright\arraybackslash}p{(\columnwidth - 2\tabcolsep) * \real{0.2909}}
  >{\raggedright\arraybackslash}p{(\columnwidth - 2\tabcolsep) * \real{0.7091}}@{}}
\toprule\noalign{}
\begin{minipage}[b]{\linewidth}\raggedright
Column
\end{minipage} &
\begin{minipage}[b]{\linewidth}\raggedright
Meaning
\end{minipage} \\
\midrule\noalign{}
\endhead
\bottomrule\noalign{}
\endlastfoot
\texttt{verdict} &
confirmed\_reduced,
confirmed\_inactive,
active, not\_stated,
and the summary
rows \\
\texttt{n\_pairs} /
\texttt{denominator}
& number of pairs
and the total
(30) \\
\texttt{note} &
interval of the
mislabelling rate \\
\end{longtable}

}

\textbf{Excerpt}
(first 3 rows).

{\footnotesize

\begin{longtable}[]{@{}
  >{\raggedright\arraybackslash}p{(\columnwidth - 6\tabcolsep) * \real{0.2500}}
  >{\raggedright\arraybackslash}p{(\columnwidth - 6\tabcolsep) * \real{0.2500}}
  >{\raggedright\arraybackslash}p{(\columnwidth - 6\tabcolsep) * \real{0.2500}}
  >{\raggedright\arraybackslash}p{(\columnwidth - 6\tabcolsep) * \real{0.2500}}@{}}
\toprule\noalign{}
\begin{minipage}[b]{\linewidth}\raggedright
\texttt{verdict}
\end{minipage} &
\begin{minipage}[b]{\linewidth}\raggedright
\texttt{n\_pairs}
\end{minipage} &
\begin{minipage}[b]{\linewidth}\raggedright
\texttt{denominator}
\end{minipage} &
\begin{minipage}[b]{\linewidth}\raggedright
\texttt{note}
\end{minipage} \\
\midrule\noalign{}
\endhead
\bottomrule\noalign{}
\endlastfoot
\texttt{confirmed\_reduced}
& \texttt{15} &
\texttt{30} &
\texttt{} \\
\texttt{not\_stated}
& \texttt{7} &
\texttt{30} &
\texttt{} \\
\texttt{confirmed\_inactive}
& \texttt{4} &
\texttt{30} &
\texttt{} \\
\end{longtable}

}

\hypertarget{S11}{%
\subsubsection{Table
S11. Zymogen-mature
evaluation set: all
158 metrics ranked
by AUROC (complete
Table
5)}\label{S11}}

\begin{shaded}\small

\textbf{File.}
\texttt{supplementary/S11\_zymogen\_21pair\_all\_metrics.csv}
· 158 rows × 27
columns · public

\textbf{Supports.}
Results 3.1.1; the
complete version of
Table 5.
\textbf{Cited in.}
Methods 2.2.1;
Results 3.1.1.

\textbf{A row is}
one metric scored on
the 21 pairs (158 in
all).

\end{shaded}

Every metric scored
on the 21 pairs that
share a ligand:
structure-space,
PLACER and
prediction-based
metrics, including
whole-protein and
sequence-only
comparators, ranked
by direction-free
AUROC of dead
against active. The
best-of-158
threshold (95th
percentile of the
best AUROC among 158
unrelated metrics in
label-swap
simulations) is
marked: seven
metrics exceed it,
all summaries of
whole-prediction
confidence (pLDDT
and per-chain pTM),
which are
comparators and not
main-set metrics.

\textbf{Key
columns.}

{\small

\begin{longtable}[]{@{}
  >{\raggedright\arraybackslash}p{(\columnwidth - 2\tabcolsep) * \real{0.2909}}
  >{\raggedright\arraybackslash}p{(\columnwidth - 2\tabcolsep) * \real{0.7091}}@{}}
\toprule\noalign{}
\begin{minipage}[b]{\linewidth}\raggedright
Column
\end{minipage} &
\begin{minipage}[b]{\linewidth}\raggedright
Meaning
\end{minipage} \\
\midrule\noalign{}
\endhead
\bottomrule\noalign{}
\endlastfoot
\texttt{rank} /
\texttt{rank\_in\_table\_5}
& rank among all 158
and rank among the
39 items of Table
5 \\
\texttt{metric\_name}
/
\texttt{canonical\_key}
/ \texttt{source} /
\texttt{input\_class}
& the metric, the
distinct metric,
structure-space /
PLACER /
prediction-based,
and class \\
\texttt{role} &
main, reference or a
supplement role (43
rows carry main
because aliases of
the 33 main-set
metrics are separate
rows) \\
\texttt{n\_pairs} /
\texttt{auroc} /
\texttt{auroc\_ci\_lo}
/
\texttt{auroc\_ci\_hi}
& pairs (21) and
AUROC with 95\%
interval over
pairs \\
\texttt{higher\_in}
/
\texttt{auroc\_directed}
& the form (dead or
active) in which the
metric is larger,
and the directed
AUROC \\
\texttt{above\_best\_of\_158\_threshold}
/
\texttt{applicability\_status}
& 1 if above the
threshold;
applicability
status \\
\end{longtable}

}

The shared block of
12 provenance
columns (see
Conventions) is also
present and is not
repeated above.

\textbf{Excerpt}
(first 3 rows).

{\footnotesize

\begin{longtable}[]{@{}
  >{\raggedright\arraybackslash}p{(\columnwidth - 10\tabcolsep) * \real{0.1667}}
  >{\raggedright\arraybackslash}p{(\columnwidth - 10\tabcolsep) * \real{0.1667}}
  >{\raggedright\arraybackslash}p{(\columnwidth - 10\tabcolsep) * \real{0.1667}}
  >{\raggedright\arraybackslash}p{(\columnwidth - 10\tabcolsep) * \real{0.1667}}
  >{\raggedright\arraybackslash}p{(\columnwidth - 10\tabcolsep) * \real{0.1667}}
  >{\raggedright\arraybackslash}p{(\columnwidth - 10\tabcolsep) * \real{0.1667}}@{}}
\toprule\noalign{}
\begin{minipage}[b]{\linewidth}\raggedright
\texttt{rank}
\end{minipage} &
\begin{minipage}[b]{\linewidth}\raggedright
\texttt{metric\_name}
\end{minipage} &
\begin{minipage}[b]{\linewidth}\raggedright
\texttt{role}
\end{minipage} &
\begin{minipage}[b]{\linewidth}\raggedright
\texttt{auroc}
\end{minipage} &
\begin{minipage}[b]{\linewidth}\raggedright
\texttt{auroc\_ci\_lo}
\end{minipage} &
\begin{minipage}[b]{\linewidth}\raggedright
\texttt{auroc\_ci\_hi}
\end{minipage} \\
\midrule\noalign{}
\endhead
\bottomrule\noalign{}
\endlastfoot
\texttt{1} &
\texttt{chai\_plddt\_min}
&
\texttt{supp\_global}
& \texttt{0.8821} &
\texttt{0.7483} &
\texttt{0.9932} \\
\texttt{2} &
\texttt{chai\_plddt\_best\_model}
&
\texttt{supp\_global}
& \texttt{0.7914} &
\texttt{0.6757} &
\texttt{0.9025} \\
\texttt{3} &
\texttt{chai\_plddt\_mean}
& \texttt{reference}
& \texttt{0.7914} &
\texttt{0.6757} &
\texttt{0.9025} \\
\end{longtable}

}

\hypertarget{S12}{%
\subsubsection{Table
S12. Reference rows
in full, across the
activity
tests}\label{S12}}

\begin{shaded}\small

\textbf{File.}
\texttt{supplementary/S12\_reference\_rows.csv}
· 35 rows × 9
columns · D2DCure
aggregate (licence
of the source data
unstated; aggregates
only, no per-variant
rows)

\textbf{Supports.}
Results 3.1 and 3.3;
Tables 5, 6, 8 and
9; Methods 2.2 and
2.4. \textbf{Cited
in.} Methods 2.2.1,
2.4.1; Results
3.1.1, 3.1.2, 3.3.1.

\textbf{A row is}
one reference item
for one measure (35
in all).

\end{shaded}

The rows that need
no catalytic
information,
reported in full
with intervals:
crystallographic
resolution, sequence
length, net charge
and tyrosine
fraction (AUROC on
the 49 zymogen
pairs) and pLDDT and
pTM (AUROC on the 28
re-predicted pairs);
the five declared
BglB baselines
(side-chain volume
change, negative
distance to the
catalytic residues,
burial, BLOSUM62 and
the Rosetta score of
the BglB data;
signed AUROC on 432
variants); and
expected hits per
96-well plate at
each keep fraction
for tyrosine
fraction, net charge
and sequence length
on the plates.

\textbf{Key
columns.}

{\small

\begin{longtable}[]{@{}
  >{\raggedright\arraybackslash}p{(\columnwidth - 2\tabcolsep) * \real{0.2909}}
  >{\raggedright\arraybackslash}p{(\columnwidth - 2\tabcolsep) * \real{0.7091}}@{}}
\toprule\noalign{}
\begin{minipage}[b]{\linewidth}\raggedright
Column
\end{minipage} &
\begin{minipage}[b]{\linewidth}\raggedright
Meaning
\end{minipage} \\
\midrule\noalign{}
\endhead
\bottomrule\noalign{}
\endlastfoot
\texttt{reference\_item}
/ \texttt{metric} &
the reference row
(and the end kept
for plate rows) \\
\texttt{measure} &
AUROC
(direction-free),
AUROC (signed,
pre-directed) or
expected hits per
96-well plate \\
\texttt{keep\_fraction}
& fraction of
designs kept (plate
rows only) \\
\texttt{estimate} /
\texttt{ci\_lo} /
\texttt{ci\_hi} /
\texttt{n} & the
value, its interval
and the number of
units \\
\end{longtable}

}

\textbf{Excerpt}
(first 3 rows where
measure = AUROC
(signed,
pre-directed)).

{\footnotesize

\begin{longtable}[]{@{}
  >{\raggedright\arraybackslash}p{(\columnwidth - 10\tabcolsep) * \real{0.1667}}
  >{\raggedright\arraybackslash}p{(\columnwidth - 10\tabcolsep) * \real{0.1667}}
  >{\raggedright\arraybackslash}p{(\columnwidth - 10\tabcolsep) * \real{0.1667}}
  >{\raggedright\arraybackslash}p{(\columnwidth - 10\tabcolsep) * \real{0.1667}}
  >{\raggedright\arraybackslash}p{(\columnwidth - 10\tabcolsep) * \real{0.1667}}
  >{\raggedright\arraybackslash}p{(\columnwidth - 10\tabcolsep) * \real{0.1667}}@{}}
\toprule\noalign{}
\begin{minipage}[b]{\linewidth}\raggedright
\texttt{reference\_item}
\end{minipage} &
\begin{minipage}[b]{\linewidth}\raggedright
\texttt{measure}
\end{minipage} &
\begin{minipage}[b]{\linewidth}\raggedright
\texttt{estimate}
\end{minipage} &
\begin{minipage}[b]{\linewidth}\raggedright
\texttt{ci\_lo}
\end{minipage} &
\begin{minipage}[b]{\linewidth}\raggedright
\texttt{ci\_hi}
\end{minipage} &
\begin{minipage}[b]{\linewidth}\raggedright
\texttt{n}
\end{minipage} \\
\midrule\noalign{}
\endhead
\bottomrule\noalign{}
\endlastfoot
\texttt{abs\_dvol\_A3}
&
\texttt{AUROC\ (signed,\ pre-directed)}
& \texttt{0.5513} &
\texttt{0.4728} &
\texttt{0.6278} &
\texttt{432} \\
\texttt{neg\_d\_cat}
&
\texttt{AUROC\ (signed,\ pre-directed)}
& \texttt{0.7633} &
\texttt{0.6872} &
\texttt{0.8243} &
\texttt{432} \\
\texttt{neg\_rel\_sasa}
&
\texttt{AUROC\ (signed,\ pre-directed)}
& \texttt{0.6831} &
\texttt{0.6058} &
\texttt{0.748} &
\texttt{432} \\
\end{longtable}

}

\textbf{Notes.}

\begin{itemize}
\tightlist
\item
  Contains
  aggregates of
  D2DCure data
  (licence
  unstated); no
  per-variant rows.
\end{itemize}

\hypertarget{S13}{%
\subsubsection{Table
S13. Plated designs:
all 128 metrics
ranked by the AUROC
of active against no
active designs
(complete Table
6)}\label{S13}}

\begin{shaded}\small

\textbf{File.}
\texttt{supplementary/S13\_plates\_all\_metrics\_auroc.csv}
· 128 rows × 29
columns · public

\textbf{Supports.}
Results 3.1.2; the
complete version of
Table 6.
\textbf{Cited in.}
Methods 2.2.2;
Results 3.1.2.

\textbf{A row is}
one metric scored on
the 192 designs (128
in all).

\end{shaded}

Every metric scored
on the 192 plated
designs, ranked by
the direction-free
AUROC of the 16
active designs
against the 176 no
active designs, with
95\% interval over
the 136 backbone
clusters. The
best-of-128
threshold (permuting
labels among
designs) is marked:
ten metrics exceed
it. They are two
main-set metrics
(the intra-residue
repulsion, 0.80, and
the repulsion of the
catalytic residues,
0.78), an alias of
the latter, and
sequence-composition
and whole-protein
comparators.

\textbf{Key
columns.}

{\small

\begin{longtable}[]{@{}
  >{\raggedright\arraybackslash}p{(\columnwidth - 2\tabcolsep) * \real{0.2909}}
  >{\raggedright\arraybackslash}p{(\columnwidth - 2\tabcolsep) * \real{0.7091}}@{}}
\toprule\noalign{}
\begin{minipage}[b]{\linewidth}\raggedright
Column
\end{minipage} &
\begin{minipage}[b]{\linewidth}\raggedright
Meaning
\end{minipage} \\
\midrule\noalign{}
\endhead
\bottomrule\noalign{}
\endlastfoot
\texttt{rank} /
\texttt{rank\_in\_table\_6}
/
\texttt{in\_table\_6}
& rank among 128,
rank among the 36
items of Table 6 and
membership of Table
6 \\
\texttt{metric\_name}
/
\texttt{canonical\_key}
/ \texttt{source} /
\texttt{input\_class}
/ \texttt{role} &
the metric, its
source, class and
role \\
\texttt{n\_designs}
/ \texttt{n\_active}
& 192 and 16 \\
\texttt{auroc} /
\texttt{auroc\_ci\_lo}
/
\texttt{auroc\_ci\_hi}
/
\texttt{auroc\_directed}
/
\texttt{higher\_in}
& AUROC, interval,
directed AUROC and
the group in which
the metric is
larger \\
\texttt{above\_best\_of\_all\_metrics\_threshold}
/
\texttt{applicability\_status}
& 1 if above the
best-of-128
threshold;
applicability
status \\
\end{longtable}

}

The shared block of
12 provenance
columns (see
Conventions) is also
present and is not
repeated above.

\textbf{Excerpt}
(first 3 rows).

{\footnotesize

\begin{longtable}[]{@{}
  >{\raggedright\arraybackslash}p{(\columnwidth - 10\tabcolsep) * \real{0.1667}}
  >{\raggedright\arraybackslash}p{(\columnwidth - 10\tabcolsep) * \real{0.1667}}
  >{\raggedright\arraybackslash}p{(\columnwidth - 10\tabcolsep) * \real{0.1667}}
  >{\raggedright\arraybackslash}p{(\columnwidth - 10\tabcolsep) * \real{0.1667}}
  >{\raggedright\arraybackslash}p{(\columnwidth - 10\tabcolsep) * \real{0.1667}}
  >{\raggedright\arraybackslash}p{(\columnwidth - 10\tabcolsep) * \real{0.1667}}@{}}
\toprule\noalign{}
\begin{minipage}[b]{\linewidth}\raggedright
\texttt{rank}
\end{minipage} &
\begin{minipage}[b]{\linewidth}\raggedright
\texttt{metric\_name}
\end{minipage} &
\begin{minipage}[b]{\linewidth}\raggedright
\texttt{role}
\end{minipage} &
\begin{minipage}[b]{\linewidth}\raggedright
\texttt{auroc}
\end{minipage} &
\begin{minipage}[b]{\linewidth}\raggedright
\texttt{auroc\_ci\_lo}
\end{minipage} &
\begin{minipage}[b]{\linewidth}\raggedright
\texttt{auroc\_ci\_hi}
\end{minipage} \\
\midrule\noalign{}
\endhead
\bottomrule\noalign{}
\endlastfoot
\texttt{1} &
\texttt{rosetta\_cat\_fa\_intra\_rep}
& \texttt{main} &
\texttt{0.8018} &
\texttt{0.6385} &
\texttt{0.8929} \\
\texttt{2} &
\texttt{frac\_G} &
\texttt{supp\_global}
& \texttt{0.7926} &
\texttt{0.6157} &
\texttt{0.8698} \\
\texttt{3} &
\texttt{rosetta\_ref}
&
\texttt{excluded\_bookkeeping}
& \texttt{0.7876} &
\texttt{0.6708} &
\texttt{0.8914} \\
\end{longtable}

}

\textbf{Notes.}

\begin{itemize}
\tightlist
\item
  Designs are not
  exchangeable (10
  of the 16 active
  designs lie in two
  backbone
  clusters), so the
  threshold
  understates the
  chance range.
\end{itemize}

\hypertarget{S14}{%
\subsubsection{Table
S14. Substrate swap
on all systems: all
33 prediction-based
quantities, natural
enzymes and de novo
designs}\label{S14}}

\begin{shaded}\small

\textbf{File.}
\texttt{supplementary/S14\_substrate\_swap\_per\_metric.csv}
· 242 rows × 17
columns · public

\textbf{Supports.}
Results 3.2; the
all-system analysis
behind Table 7.
\textbf{Cited in.}
Methods 2.3; Results
3.2.

\textbf{A row is}
one prediction-based
quantity under one
wrong-ligand
condition, for one
source (natural
enzymes or de novo
designs).

\end{shaded}

The substrate swap
before the 55-enzyme
evaluation set was
fixed: all 59
natural enzymes
(cognate value
averaged over five
seeds) and all 30 de
novo designs, for
every one of the 33
prediction-based
quantities that were
scored, including
the Chai-1 combined
score (a rescaled
ipTM) that Table 7
leaves out. For each
quantity and each
wrong ligand (same
enzyme class,
different class,
decoy; apo for the
quantities that are
defined without a
ligand) it gives the
mean under the
cognate and under
the other condition,
the paired change,
the change in units
of seed noise, and
the direction-free
AUROC with its
interval.

\textbf{Key
columns.}

{\small

\begin{longtable}[]{@{}
  >{\raggedright\arraybackslash}p{(\columnwidth - 2\tabcolsep) * \real{0.2909}}
  >{\raggedright\arraybackslash}p{(\columnwidth - 2\tabcolsep) * \real{0.7091}}@{}}
\toprule\noalign{}
\begin{minipage}[b]{\linewidth}\raggedright
Column
\end{minipage} &
\begin{minipage}[b]{\linewidth}\raggedright
Meaning
\end{minipage} \\
\midrule\noalign{}
\endhead
\bottomrule\noalign{}
\endlastfoot
\texttt{source} &
natural or denovo \\
\texttt{metric\_name}
/
\texttt{canonical\_key}
/
\texttt{input\_class}
& the quantity, the
distinct metric it
belongs to and its
class \\
\texttt{ame\_flavour}
& crystal =
active-site accuracy
against the
deposited structure;
self\_design =
against the design's
own model (never
pooled) \\
\texttt{in\_main\_set}
& 1 for the four
main-set metrics
(each appears
several times: mean,
maximum and spread
over the five
predicted models are
separate rows) \\
\texttt{condition} &
same\_ec, diff\_ec,
decoy or apo \\
\texttt{n\_systems}
& systems with a
value (59 natural,
30 de novo) \\
\texttt{mean\_cognate}
/
\texttt{mean\_condition}
/
\texttt{mean\_delta}
& mean value under
the cognate ligand,
under the condition,
and the paired
difference \\
\texttt{median\_delta\_over\_sigma}
/
\texttt{frac\_beyond\_1sigma}
& median change in
units of the
seed-to-seed
standard deviation,
and the share of
systems that move by
more than one \\
\texttt{auroc} /
\texttt{auroc\_lo} /
\texttt{auroc\_hi} &
direction-free AUROC
of cognate against
the condition, 95\%
interval over
systems \\
\end{longtable}

}

\textbf{Excerpt}
(first 3 rows where
source = natural,
condition =
diff\_ec,
in\_main\_set = 1).

{\footnotesize

\begin{longtable}[]{@{}
  >{\raggedright\arraybackslash}p{(\columnwidth - 12\tabcolsep) * \real{0.1429}}
  >{\raggedright\arraybackslash}p{(\columnwidth - 12\tabcolsep) * \real{0.1429}}
  >{\raggedright\arraybackslash}p{(\columnwidth - 12\tabcolsep) * \real{0.1429}}
  >{\raggedright\arraybackslash}p{(\columnwidth - 12\tabcolsep) * \real{0.1429}}
  >{\raggedright\arraybackslash}p{(\columnwidth - 12\tabcolsep) * \real{0.1429}}
  >{\raggedright\arraybackslash}p{(\columnwidth - 12\tabcolsep) * \real{0.1429}}
  >{\raggedright\arraybackslash}p{(\columnwidth - 12\tabcolsep) * \real{0.1429}}@{}}
\toprule\noalign{}
\begin{minipage}[b]{\linewidth}\raggedright
\texttt{source}
\end{minipage} &
\begin{minipage}[b]{\linewidth}\raggedright
\texttt{metric\_name}
\end{minipage} &
\begin{minipage}[b]{\linewidth}\raggedright
\texttt{condition}
\end{minipage} &
\begin{minipage}[b]{\linewidth}\raggedright
\texttt{n\_systems}
\end{minipage} &
\begin{minipage}[b]{\linewidth}\raggedright
\texttt{mean\_cognate}
\end{minipage} &
\begin{minipage}[b]{\linewidth}\raggedright
\texttt{mean\_condition}
\end{minipage} &
\begin{minipage}[b]{\linewidth}\raggedright
\texttt{auroc}
\end{minipage} \\
\midrule\noalign{}
\endhead
\bottomrule\noalign{}
\endlastfoot
\texttt{natural} &
\texttt{ame\_crystal\_clash\_max}
& \texttt{diff\_ec}
& \texttt{58} &
\texttt{3.2826} &
\texttt{3.6103} &
\texttt{0.6564} \\
\texttt{natural} &
\texttt{ame\_crystal\_rmsd\_max}
& \texttt{diff\_ec}
& \texttt{59} &
\texttt{1.4469} &
\texttt{1.5572} &
\texttt{0.5395} \\
\texttt{natural} &
\texttt{ame\_crystal\_rmsd\_mean}
& \texttt{diff\_ec}
& \texttt{59} &
\texttt{1.2516} &
\texttt{1.2895} &
\texttt{0.5306} \\
\end{longtable}

}

\textbf{Notes.}

\begin{itemize}
\tightlist
\item
  n is 59 natural
  enzymes here and
  55 in Table 7,
  whose evaluation
  set drops four
  enzymes whose
  ligand could not
  be built.
\item
  Per-metric results
  of this table are
  descriptive; the
  ranking and the
  best-of-30
  threshold are in
  Table 7.
\end{itemize}

\hypertarget{S15}{%
\subsubsection{Table
S15. Substrate swap
stratified by ligand
size and
charge}\label{S15}}

\begin{shaded}\small

\textbf{File.}
\texttt{supplementary/S15\_substrate\_swap\_size\_strata.csv}
· 1,136 rows × 27
columns · public

\textbf{Supports.}
Results 3.2; Methods
2.3 (size check, Eq.
10). \textbf{Cited
in.} Methods 2.3;
Results 3.2.

\textbf{A row is}
one quantity,
source, wrong-ligand
condition and
stratum of the
difference in ligand
size or charge, for
one analysis
variant.

\end{shaded}

The check of whether
ligand size or
charge explains the
discrimination. For
each headline
quantity (ipTM,
per-chain pTM
minimum, the
combined score,
ligand clearance),
comparator (pTM,
pLDDT, AME RMSD) and
the two baselines
(ligand heavy-atom
count and formal
charge), the cells
(an enzyme and one
wrong ligand) are
stratified by the
difference in
heavy-atom count
(smaller, matched,
larger), by the
charge difference
(same, different)
and jointly. For
each stratum it
gives the paired win
fraction, the AUROC
of the absolute
value, the Spearman
correlation of the
size difference with
the change,
intervals over
systems and
clusters, and
whether the effect
survives.

\textbf{Key
columns.}

{\small

\begin{longtable}[]{@{}
  >{\raggedright\arraybackslash}p{(\columnwidth - 2\tabcolsep) * \real{0.2909}}
  >{\raggedright\arraybackslash}p{(\columnwidth - 2\tabcolsep) * \real{0.7091}}@{}}
\toprule\noalign{}
\begin{minipage}[b]{\linewidth}\raggedright
Column
\end{minipage} &
\begin{minipage}[b]{\linewidth}\raggedright
Meaning
\end{minipage} \\
\midrule\noalign{}
\endhead
\bottomrule\noalign{}
\endlastfoot
\texttt{source} /
\texttt{variant} &
natural or denovo;
primary or the
variant that
excludes ipTM = 0
cells
(excl\_iptm0) \\
\texttt{metric\_name}
/ \texttt{role} &
the quantity and
headline, comparator
or baseline \\
\texttt{stratum\_family}
/ \texttt{stratum} &
all, size (smaller,
matched, larger),
charge (q\_same,
q\_diff) or joint,
and the stratum \\
\texttt{condition} &
same\_ec, diff\_ec,
decoy or the three
pooled
(pooled\_3) \\
\texttt{k\_cells} /
\texttt{N\_cells} /
\texttt{k\_systems}
/
\texttt{N\_systems}
& cells and systems
in the stratum \\
\texttt{win\_frac} /
\texttt{win\_lo95} /
\texttt{win\_hi95} &
paired win fraction
(cognate higher than
wrong ligand) with
95\% interval \\
\texttt{auroc\_abs}
/
\texttt{spearman\_dn\_vs\_change}
/ \texttt{survives}
& AUROC of absolute
value, Spearman
correlation of size
difference with the
change, and 1 / 0 /
`n\textless10, not
estimated' \\
\end{longtable}

}

\textbf{Excerpt}
(first 3 rows where
variant = primary,
source = natural,
stratum\_family =
size).

{\footnotesize

\begin{longtable}[]{@{}
  >{\raggedright\arraybackslash}p{(\columnwidth - 10\tabcolsep) * \real{0.1667}}
  >{\raggedright\arraybackslash}p{(\columnwidth - 10\tabcolsep) * \real{0.1667}}
  >{\raggedright\arraybackslash}p{(\columnwidth - 10\tabcolsep) * \real{0.1667}}
  >{\raggedright\arraybackslash}p{(\columnwidth - 10\tabcolsep) * \real{0.1667}}
  >{\raggedright\arraybackslash}p{(\columnwidth - 10\tabcolsep) * \real{0.1667}}
  >{\raggedright\arraybackslash}p{(\columnwidth - 10\tabcolsep) * \real{0.1667}}@{}}
\toprule\noalign{}
\begin{minipage}[b]{\linewidth}\raggedright
\texttt{metric\_name}
\end{minipage} &
\begin{minipage}[b]{\linewidth}\raggedright
\texttt{stratum}
\end{minipage} &
\begin{minipage}[b]{\linewidth}\raggedright
\texttt{condition}
\end{minipage} &
\begin{minipage}[b]{\linewidth}\raggedright
\texttt{k\_cells}
\end{minipage} &
\begin{minipage}[b]{\linewidth}\raggedright
\texttt{N\_cells}
\end{minipage} &
\begin{minipage}[b]{\linewidth}\raggedright
\texttt{win\_frac}
\end{minipage} \\
\midrule\noalign{}
\endhead
\bottomrule\noalign{}
\endlastfoot
\texttt{chai\_iptm\_mean}
& \texttt{smaller} &
\texttt{same\_ec} &
\texttt{18} &
\texttt{177} &
\texttt{1.0} \\
\texttt{chai\_aggregate\_score\_mean}
& \texttt{smaller} &
\texttt{same\_ec} &
\texttt{18} &
\texttt{177} &
\texttt{1.0} \\
\texttt{chai\_per\_chain\_ptm\_min\_mean}
& \texttt{smaller} &
\texttt{same\_ec} &
\texttt{18} &
\texttt{177} &
\texttt{0.8333} \\
\end{longtable}

}

\textbf{Notes.}

\begin{itemize}
\tightlist
\item
  Released as a CSV
  only for the full
  1,136 rows; the
  excerpt shows the
  structure.
\end{itemize}

\hypertarget{S16}{%
\subsubsection{Table
S16. BglB: every
analysed quantity
ranked by rank
correlation with the
measured impairment
(complete Table
8)}\label{S16}}

\begin{shaded}\small

\textbf{File.}
\texttt{supplementary/S16\_bglb\_all\_quantities\_ranked.csv}
· 114 rows × 35
columns · D2DCure
aggregate (licence
of the source data
unstated; aggregates
only, no per-variant
rows)

\textbf{Supports.}
Results 3.3.1; the
complete version of
Table 8.
\textbf{Cited in.}
Methods 2.4, 2.4.1;
Results 3.3.1.

\textbf{A row is}
one quantity
analysed on the 432
BglB variants (114
in all).

\end{shaded}

Every quantity
analysed on the
432-variant
evaluation set: 84
structure-space
metrics, 25
prediction-based
metrics and the five
declared baselines,
ranked by the
Spearman correlation
of the size of its
change with the
measured impairment,
with 95\% interval
over the 175
positions. It
carries the
best-of-114
threshold, whether
the quantity beats
the distance
baseline, the
correlation of the
signed change, and
beside them the
AUROC of the binary
contrast (impaired
against
wild-type-like
variants) with its
interval, threshold
and
beats-the-distance
flag. No quantity
beats the distance
baseline; 34 exceed
the best-of-114
threshold.

\textbf{Key
columns.}

{\small

\begin{longtable}[]{@{}
  >{\raggedright\arraybackslash}p{(\columnwidth - 2\tabcolsep) * \real{0.2909}}
  >{\raggedright\arraybackslash}p{(\columnwidth - 2\tabcolsep) * \real{0.7091}}@{}}
\toprule\noalign{}
\begin{minipage}[b]{\linewidth}\raggedright
Column
\end{minipage} &
\begin{minipage}[b]{\linewidth}\raggedright
Meaning
\end{minipage} \\
\midrule\noalign{}
\endhead
\bottomrule\noalign{}
\endlastfoot
\texttt{rank} /
\texttt{rank\_in\_table\_8}
/
\texttt{in\_table\_8}
& rank among 114,
rank among the 41
items of Table 8 and
membership of Table
8 \\
\texttt{metric\_name}
/
\texttt{canonical\_key}
/ \texttt{source} /
\texttt{input\_class}
/ \texttt{role} &
the quantity, its
source
(structure-space,
prediction-based or
baseline), class and
role \\
\texttt{n\_variants}
& variants with a
value (431 or 432;
248 for the Rosetta
score of the BglB
data) \\
\texttt{rho} /
\texttt{rho\_ci\_lo}
/
\texttt{rho\_ci\_hi}
& Spearman
correlation and 95\%
interval over
positions \\
\texttt{above\_best\_of\_all\_quantities\_threshold}
/
\texttt{beats\_distance\_baseline}
/
\texttt{diff\_vs\_distance\_baseline\_ci\_lo}
& threshold flag,
beats-the-distance
flag and the lower
end of the
difference in
correlation \\
\texttt{rho\_signed\_delta}
& correlation of the
signed change \\
\texttt{auroc} /
\texttt{auroc\_ci\_lo}
/
\texttt{auroc\_ci\_hi}
& AUROC of the
binary contrast with
95\% interval; the
columns that follow
(threshold,
beats-the-distance
flag and difference
against the distance
baseline) are for
the AUROC \\
\end{longtable}

}

The shared block of
12 provenance
columns (see
Conventions) is also
present and is not
repeated above.

\textbf{Excerpt}
(first 3 rows).

{\footnotesize

\begin{longtable}[]{@{}
  >{\raggedright\arraybackslash}p{(\columnwidth - 10\tabcolsep) * \real{0.1667}}
  >{\raggedright\arraybackslash}p{(\columnwidth - 10\tabcolsep) * \real{0.1667}}
  >{\raggedright\arraybackslash}p{(\columnwidth - 10\tabcolsep) * \real{0.1667}}
  >{\raggedright\arraybackslash}p{(\columnwidth - 10\tabcolsep) * \real{0.1667}}
  >{\raggedright\arraybackslash}p{(\columnwidth - 10\tabcolsep) * \real{0.1667}}
  >{\raggedright\arraybackslash}p{(\columnwidth - 10\tabcolsep) * \real{0.1667}}@{}}
\toprule\noalign{}
\begin{minipage}[b]{\linewidth}\raggedright
\texttt{rank}
\end{minipage} &
\begin{minipage}[b]{\linewidth}\raggedright
\texttt{metric\_name}
\end{minipage} &
\begin{minipage}[b]{\linewidth}\raggedright
\texttt{role}
\end{minipage} &
\begin{minipage}[b]{\linewidth}\raggedright
\texttt{rho}
\end{minipage} &
\begin{minipage}[b]{\linewidth}\raggedright
\texttt{rho\_ci\_lo}
\end{minipage} &
\begin{minipage}[b]{\linewidth}\raggedright
\texttt{rho\_ci\_hi}
\end{minipage} \\
\midrule\noalign{}
\endhead
\bottomrule\noalign{}
\endlastfoot
\texttt{1} &
\texttt{propka\_catalytic\_pka\_range}
& \texttt{main} &
\texttt{0.48} &
\texttt{0.3474} &
\texttt{0.592} \\
\texttt{2} &
\texttt{propka\_catalytic\_pka\_min}
& \texttt{main} &
\texttt{0.4622} &
\texttt{0.3224} &
\texttt{0.5748} \\
\texttt{3} &
\texttt{propka\_catalytic\_shift\_absmax}
& \texttt{main} &
\texttt{0.4587} &
\texttt{0.319} &
\texttt{0.5782} \\
\end{longtable}

}

\textbf{Notes.}

\begin{itemize}
\tightlist
\item
  Aggregates of
  D2DCure data
  (licence
  unstated); no
  per-variant rows.
\end{itemize}

\hypertarget{S17}{%
\subsubsection{Table
S17. BglB: all 84
structure-space and
comparator
metrics}\label{S17}}

\begin{shaded}\small

\textbf{File.}
\texttt{supplementary/S17\_bglb\_all\_metrics.csv}
· 84 rows × 29
columns · D2DCure
aggregate (licence
of the source data
unstated; aggregates
only, no per-variant
rows)

\textbf{Supports.}
Results 3.3.1;
Methods 2.4.1.
\textbf{Cited in.}
Methods 2.4.1;
Results 3.3.1.

\textbf{A row is}
one metric analysed
on the 432 BglB
variants (84 in
all).

\end{shaded}

Spearman
correlation, signed,
direction-free and
magnitude AUROC with
intervals, the
paired difference
against the distance
baseline and whether
the metric beats all
baselines, for every
structure-space
metric and
comparator (36
site-scoped, 28
sequence-only, 15
whole-protein, 5
bookkeeping). None
beats all baselines.

\textbf{Key
columns.}

{\small

\begin{longtable}[]{@{}
  >{\raggedright\arraybackslash}p{(\columnwidth - 2\tabcolsep) * \real{0.2909}}
  >{\raggedright\arraybackslash}p{(\columnwidth - 2\tabcolsep) * \real{0.7091}}@{}}
\toprule\noalign{}
\begin{minipage}[b]{\linewidth}\raggedright
Column
\end{minipage} &
\begin{minipage}[b]{\linewidth}\raggedright
Meaning
\end{minipage} \\
\midrule\noalign{}
\endhead
\bottomrule\noalign{}
\endlastfoot
\texttt{metric} /
\texttt{canonical\_key}
/
\texttt{input\_class}
/
\texttt{in\_main\_set}
& the metric, its
class and main-set
membership \\
\texttt{rho} /
\texttt{rho\_lo} /
\texttt{rho\_hi} &
Spearman correlation
with 95\% interval
over positions \\
\texttt{auc\_signed}
/
\texttt{auc\_dirfree}
/ auc\_mag (+ \_lo,
\_hi) & AUROC of the
binary contrast
using the signed
change, the
direction-free form
and the magnitude,
with intervals \\
\texttt{diff\_vs\_neg\_d\_cat\_lo}
/
\texttt{beats\_all\_baselines}
& lower end of the
difference against
the distance
baseline; whether
the metric beats all
five baselines \\
\end{longtable}

}

The shared block of
12 provenance
columns (see
Conventions) is also
present and is not
repeated above.

\textbf{Excerpt}
(first 3 rows where
in\_main\_set = 1).

{\footnotesize

\begin{longtable}[]{@{}
  >{\raggedright\arraybackslash}p{(\columnwidth - 10\tabcolsep) * \real{0.1667}}
  >{\raggedright\arraybackslash}p{(\columnwidth - 10\tabcolsep) * \real{0.1667}}
  >{\raggedright\arraybackslash}p{(\columnwidth - 10\tabcolsep) * \real{0.1667}}
  >{\raggedright\arraybackslash}p{(\columnwidth - 10\tabcolsep) * \real{0.1667}}
  >{\raggedright\arraybackslash}p{(\columnwidth - 10\tabcolsep) * \real{0.1667}}
  >{\raggedright\arraybackslash}p{(\columnwidth - 10\tabcolsep) * \real{0.1667}}@{}}
\toprule\noalign{}
\begin{minipage}[b]{\linewidth}\raggedright
\texttt{metric}
\end{minipage} &
\begin{minipage}[b]{\linewidth}\raggedright
\texttt{in\_main\_set}
\end{minipage} &
\begin{minipage}[b]{\linewidth}\raggedright
\texttt{rho}
\end{minipage} &
\begin{minipage}[b]{\linewidth}\raggedright
\texttt{rho\_lo}
\end{minipage} &
\begin{minipage}[b]{\linewidth}\raggedright
\texttt{rho\_hi}
\end{minipage} &
\begin{minipage}[b]{\linewidth}\raggedright
\texttt{beats\_all\_baselines}
\end{minipage} \\
\midrule\noalign{}
\endhead
\bottomrule\noalign{}
\endlastfoot
\texttt{catalytic\_frac\_buried}
& \texttt{1} &
\texttt{-0.1909} &
\texttt{-0.3048} &
\texttt{-0.057} &
\texttt{0} \\
\texttt{catalytic\_pairwise\_dist\_max}
& \texttt{1} &
\texttt{0.0431} &
\texttt{-0.1841} &
\texttt{0.2709} &
\texttt{0} \\
\texttt{catalytic\_pairwise\_dist\_mean}
& \texttt{1} &
\texttt{0.09} &
\texttt{-0.074} &
\texttt{0.2468} &
\texttt{0} \\
\end{longtable}

}

\textbf{Notes.}

\begin{itemize}
\tightlist
\item
  Aggregates of
  D2DCure data
  (licence
  unstated); no
  per-variant rows.
\end{itemize}

\hypertarget{S18}{%
\subsubsection{Table
S18. BglB:
sensitivity
analyses}\label{S18}}

\begin{shaded}\small

\textbf{File.}
\texttt{supplementary/S18\_bglb\_sensitivities.csv}
· 11 rows × 22
columns · D2DCure
aggregate (licence
of the source data
unstated; aggregates
only, no per-variant
rows)

\textbf{Supports.}
Results 3.3.1;
Discussion.
\textbf{Cited in.}
Methods 2.4.1;
Results 3.3.1;
Discussion.

\textbf{A row is}
the main analysis
and one of ten
sensitivity analyses
(11 rows).

\end{shaded}

The main analysis
and ten sensitivity
analyses with their
sizes, the number of
metrics with an
interval above 0.5,
the number that beat
the distance
baseline and the
AUROC of the
distance baseline.
The analyses are:
the main analysis
(432 variants);
variants near to and
distal from the
catalytic residues
(144 and 288);
variants with a
measured stability
(292); each
institution left out
in turn (FIU,
MiraCosta, UC
Davis); other
thresholds for the
impaired label (1
and 3 standard
deviations, with 173
and 139 impaired
variants);
responding metrics
only (70 metrics
analysed); and
without the variants
at catalytic
positions (431
variants, one
position removed).
In every analysis no
metric beats the
distance baseline.

\textbf{Key
columns.}

{\small

\begin{longtable}[]{@{}
  >{\raggedright\arraybackslash}p{(\columnwidth - 2\tabcolsep) * \real{0.2909}}
  >{\raggedright\arraybackslash}p{(\columnwidth - 2\tabcolsep) * \real{0.7091}}@{}}
\toprule\noalign{}
\begin{minipage}[b]{\linewidth}\raggedright
Column
\end{minipage} &
\begin{minipage}[b]{\linewidth}\raggedright
Meaning
\end{minipage} \\
\midrule\noalign{}
\endhead
\bottomrule\noalign{}
\endlastfoot
\texttt{analysis} &
the analysis \\
\texttt{n\_variants}
/
\texttt{n\_positions}
/
\texttt{n\_impaired}
/ \texttt{n\_wtlike}
& size of the
analysis \\
\texttt{n\_metrics\_analysed}
/
\texttt{n\_auc\_mag\_ci\_above\_half}
& metrics analysed
and metrics whose
AUROC interval lies
above 0.5 \\
\texttt{n\_beats\_neg\_d\_cat}
/
\texttt{n\_beats\_all\_baselines}
& metrics that beat
the distance
baseline and all
baselines \\
\texttt{auc\_neg\_d\_cat\_baseline}
& AUROC of the
distance baseline \\
\end{longtable}

}

The shared block of
12 provenance
columns (see
Conventions) is also
present and is not
repeated above.

\textbf{Excerpt}
(first 3 rows).

{\footnotesize

\begin{longtable}[]{@{}
  >{\raggedright\arraybackslash}p{(\columnwidth - 10\tabcolsep) * \real{0.1667}}
  >{\raggedright\arraybackslash}p{(\columnwidth - 10\tabcolsep) * \real{0.1667}}
  >{\raggedright\arraybackslash}p{(\columnwidth - 10\tabcolsep) * \real{0.1667}}
  >{\raggedright\arraybackslash}p{(\columnwidth - 10\tabcolsep) * \real{0.1667}}
  >{\raggedright\arraybackslash}p{(\columnwidth - 10\tabcolsep) * \real{0.1667}}
  >{\raggedright\arraybackslash}p{(\columnwidth - 10\tabcolsep) * \real{0.1667}}@{}}
\toprule\noalign{}
\begin{minipage}[b]{\linewidth}\raggedright
\texttt{analysis}
\end{minipage} &
\begin{minipage}[b]{\linewidth}\raggedright
\texttt{n\_variants}
\end{minipage} &
\begin{minipage}[b]{\linewidth}\raggedright
\texttt{n\_positions}
\end{minipage} &
\begin{minipage}[b]{\linewidth}\raggedright
\texttt{n\_impaired}
\end{minipage} &
\begin{minipage}[b]{\linewidth}\raggedright
\texttt{n\_wtlike}
\end{minipage} &
\begin{minipage}[b]{\linewidth}\raggedright
\texttt{n\_beats\_neg\_d\_cat}
\end{minipage} \\
\midrule\noalign{}
\endhead
\bottomrule\noalign{}
\endlastfoot
\texttt{main\ analysis}
& \texttt{432} &
\texttt{175} &
\texttt{166} &
\texttt{254} &
\texttt{0} \\
\texttt{variants\ distal\ from\ the\ catalyti…}
& \texttt{288} &
\texttt{136} &
\texttt{71} &
\texttt{207} &
\texttt{0} \\
\texttt{variants\ near\ the\ catalytic\ resid…}
& \texttt{144} &
\texttt{39} &
\texttt{95} &
\texttt{47} &
\texttt{0} \\
\end{longtable}

}

\textbf{Notes.}

\begin{itemize}
\tightlist
\item
  Aggregates of
  D2DCure data
  (licence
  unstated); no
  per-variant rows.
\end{itemize}

\hypertarget{S19}{%
\subsubsection{Table
S19. BglB:
prediction-based
metrics (Chai-1 with
the assay
substrate)}\label{S19}}

\begin{shaded}\small

\textbf{File.}
\texttt{supplementary/S19\_bglb\_prediction\_based\_all\_metrics.csv}
· 28 rows × 29
columns · D2DCure
aggregate (licence
of the source data
unstated; aggregates
only, no per-variant
rows)

\textbf{Supports.}
Results 3.3.1;
Methods 2.4.1.
\textbf{Cited in.}
Methods 2.4.1;
Results 3.3.1.

\textbf{A row is}
one prediction-based
metric analysed on
the 432 BglB
variants (28 in
all).

\end{shaded}

The same columns as
\protect\hyperlink{S17}{S17}
for the metrics
computed from Chai-1
predictions of each
variant with the
assay substrate
(pNPG, no metal);
active-site accuracy
is measured against
the wild-type
crystal, which
carries a covalent
glucosyl
intermediate. Twelve
rows are main-set
metrics; the largest
correlation is 0.17.

\textbf{Key
columns.}

{\small

\begin{longtable}[]{@{}
  >{\raggedright\arraybackslash}p{(\columnwidth - 2\tabcolsep) * \real{0.2909}}
  >{\raggedright\arraybackslash}p{(\columnwidth - 2\tabcolsep) * \real{0.7091}}@{}}
\toprule\noalign{}
\begin{minipage}[b]{\linewidth}\raggedright
Column
\end{minipage} &
\begin{minipage}[b]{\linewidth}\raggedright
Meaning
\end{minipage} \\
\midrule\noalign{}
\endhead
\bottomrule\noalign{}
\endlastfoot
\texttt{metric} /
\texttt{canonical\_key}
/
\texttt{input\_class}
/
\texttt{in\_main\_set}
& the metric, its
class and main-set
membership \\
\texttt{rho} /
\texttt{rho\_lo} /
\texttt{rho\_hi} &
Spearman correlation
with 95\% interval
over positions \\
\texttt{auc\_signed}
/
\texttt{auc\_dirfree}
/ auc\_mag (+ \_lo,
\_hi) & AUROC of the
binary contrast \\
\texttt{beats\_all\_baselines}
& whether the metric
beats all five
baselines \\
\end{longtable}

}

The shared block of
12 provenance
columns (see
Conventions) is also
present and is not
repeated above.

\textbf{Excerpt}
(first 3 rows where
in\_main\_set = 1).

{\footnotesize

\begin{longtable}[]{@{}
  >{\raggedright\arraybackslash}p{(\columnwidth - 8\tabcolsep) * \real{0.2000}}
  >{\raggedright\arraybackslash}p{(\columnwidth - 8\tabcolsep) * \real{0.2000}}
  >{\raggedright\arraybackslash}p{(\columnwidth - 8\tabcolsep) * \real{0.2000}}
  >{\raggedright\arraybackslash}p{(\columnwidth - 8\tabcolsep) * \real{0.2000}}
  >{\raggedright\arraybackslash}p{(\columnwidth - 8\tabcolsep) * \real{0.2000}}@{}}
\toprule\noalign{}
\begin{minipage}[b]{\linewidth}\raggedright
\texttt{metric}
\end{minipage} &
\begin{minipage}[b]{\linewidth}\raggedright
\texttt{in\_main\_set}
\end{minipage} &
\begin{minipage}[b]{\linewidth}\raggedright
\texttt{rho}
\end{minipage} &
\begin{minipage}[b]{\linewidth}\raggedright
\texttt{rho\_lo}
\end{minipage} &
\begin{minipage}[b]{\linewidth}\raggedright
\texttt{rho\_hi}
\end{minipage} \\
\midrule\noalign{}
\endhead
\bottomrule\noalign{}
\endlastfoot
\texttt{ame\_crystal\_clash\_max}
& \texttt{1} &
\texttt{0.0828} &
\texttt{-0.0534} &
\texttt{0.2162} \\
\texttt{ame\_crystal\_rmsd\_max}
& \texttt{1} &
\texttt{0.0632} &
\texttt{-0.0622} &
\texttt{0.1849} \\
\texttt{ame\_crystal\_rmsd\_mean}
& \texttt{1} &
\texttt{0.1213} &
\texttt{-0.0079} &
\texttt{0.259} \\
\end{longtable}

}

\textbf{Notes.}

\begin{itemize}
\tightlist
\item
  Aggregates of
  D2DCure data
  (licence
  unstated); no
  per-variant rows.
\end{itemize}

\hypertarget{S20}{%
\subsubsection{Table
S20. Plated designs:
hits per 96-well
plate for every
metric (filter-style
analysis)}\label{S20}}

\begin{shaded}\small

\textbf{File.}
\texttt{supplementary/S20\_plates\_hits\_per\_plate\_ranked.csv}
· 128 rows × 37
columns · public

\textbf{Supports.}
Results 3.3.2;
Methods 2.4.2 (Eqs.
16 and 17).
\textbf{Cited in.}
Methods 2.4.2;
Results 3.3.2.

\textbf{A row is}
one metric scored on
the 192 designs (128
in all).

\end{shaded}

The filter-style
analysis, for every
metric: the expected
hits per 96-well
plate and the
actives among the 48
designs kept when
the quarter of
designs that the
metric ranks
highest, or lowest,
fills the plate (8.0
hits when
unfiltered), each
with a 90\% interval
over designs, ranked
by the better end.
30 metrics exceed
the random-filter
band (14.0 hits).

\textbf{Key
columns.}

{\small

\begin{longtable}[]{@{}
  >{\raggedright\arraybackslash}p{(\columnwidth - 2\tabcolsep) * \real{0.2909}}
  >{\raggedright\arraybackslash}p{(\columnwidth - 2\tabcolsep) * \real{0.7091}}@{}}
\toprule\noalign{}
\begin{minipage}[b]{\linewidth}\raggedright
Column
\end{minipage} &
\begin{minipage}[b]{\linewidth}\raggedright
Meaning
\end{minipage} \\
\midrule\noalign{}
\endhead
\bottomrule\noalign{}
\endlastfoot
\texttt{rank} /
\texttt{rank\_among\_ranked\_items}
/
\texttt{ranked\_item}
& rank among 128,
rank among the 36
items of Table 9 and
whether the metric
is one of them \\
\texttt{metric\_name}
/
\texttt{canonical\_key}
/ \texttt{source} /
\texttt{input\_class}
/ \texttt{role} &
the metric and its
class \\
\texttt{n\_designs}
/ \texttt{n\_active}
/
\texttt{n\_kept\_highest}
/
\texttt{n\_kept\_lowest}
& 192, 16 and the
number kept (fewer
than 48 when designs
tie at the cut) \\
\texttt{hits\_highest\_quarter}
/
\texttt{hits\_highest\_ci\_lo90}
/
\texttt{hits\_highest\_ci\_hi90}
/
\texttt{actives\_in\_highest\_quarter}
& hits per plate for
the highest quarter,
90\% interval and
actives kept \\
\texttt{hits\_lowest\_quarter}
/
\texttt{hits\_lowest\_ci\_lo90}
/
\texttt{hits\_lowest\_ci\_hi90}
/
\texttt{actives\_in\_lowest\_quarter}
& the same for the
lowest quarter \\
\texttt{better\_end}
/
\texttt{best\_hits\_per\_plate}
/
\texttt{above\_random\_band}
& the better end,
its hits per plate
and whether it
exceeds the random
band \\
\texttt{note} & for
example a tie at the
cut \\
\end{longtable}

}

The shared block of
12 provenance
columns (see
Conventions) is also
present and is not
repeated above.

\textbf{Excerpt}
(first 3 rows).

{\footnotesize

\begin{longtable}[]{@{}
  >{\raggedright\arraybackslash}p{(\columnwidth - 10\tabcolsep) * \real{0.1667}}
  >{\raggedright\arraybackslash}p{(\columnwidth - 10\tabcolsep) * \real{0.1667}}
  >{\raggedright\arraybackslash}p{(\columnwidth - 10\tabcolsep) * \real{0.1667}}
  >{\raggedright\arraybackslash}p{(\columnwidth - 10\tabcolsep) * \real{0.1667}}
  >{\raggedright\arraybackslash}p{(\columnwidth - 10\tabcolsep) * \real{0.1667}}
  >{\raggedright\arraybackslash}p{(\columnwidth - 10\tabcolsep) * \real{0.1667}}@{}}
\toprule\noalign{}
\begin{minipage}[b]{\linewidth}\raggedright
\texttt{rank}
\end{minipage} &
\begin{minipage}[b]{\linewidth}\raggedright
\texttt{metric\_name}
\end{minipage} &
\begin{minipage}[b]{\linewidth}\raggedright
\texttt{role}
\end{minipage} &
\begin{minipage}[b]{\linewidth}\raggedright
\texttt{better\_end}
\end{minipage} &
\begin{minipage}[b]{\linewidth}\raggedright
\texttt{best\_hits\_per\_plate}
\end{minipage} &
\begin{minipage}[b]{\linewidth}\raggedright
\texttt{above\_random\_band}
\end{minipage} \\
\midrule\noalign{}
\endhead
\bottomrule\noalign{}
\endlastfoot
\texttt{1} &
\texttt{frac\_Y} &
\texttt{reference} &
\texttt{highest} &
\texttt{24.0} &
\texttt{1} \\
\texttt{2} &
\texttt{frac\_G} &
\texttt{supp\_global}
& \texttt{lowest} &
\texttt{22.0} &
\texttt{1} \\
\texttt{3} &
\texttt{rosetta\_cat\_fa\_intra\_rep}
& \texttt{main} &
\texttt{lowest} &
\texttt{22.0} &
\texttt{1} \\
\end{longtable}

}

\textbf{Notes.}

\begin{itemize}
\tightlist
\item
  Descriptive only:
  no hypothesis
  test.
\end{itemize}

\hypertarget{S21}{%
\subsubsection{Table
S21. Plated designs:
every metric ranked
by correlation with
kcat/KM among the 16
designs that have a
value (complete
Table
9)}\label{S21}}

\begin{shaded}\small

\textbf{File.}
\texttt{supplementary/S21\_plates\_all\_metrics\_ranked.csv}
· 106 rows × 29
columns · public

\textbf{Supports.}
Results 3.3.2; the
complete version of
Table 9.
\textbf{Cited in.}
Methods 2.4.2;
Results 3.3.2.

\textbf{A row is}
one metric with a
rank correlation on
the 16 designs with
a reported kcat/KM
(106 in all).

\end{shaded}

Every metric that
has a correlation,
ranked by the size
of the Spearman
correlation with
log10 kcat/KM, with
a 90\% interval over
designs, whether the
interval excludes 0
(33 do), the
direction, and the
correlation with the
sequence length
partialled out and
whether it survives
that control.

\textbf{Key
columns.}

{\small

\begin{longtable}[]{@{}
  >{\raggedright\arraybackslash}p{(\columnwidth - 2\tabcolsep) * \real{0.2909}}
  >{\raggedright\arraybackslash}p{(\columnwidth - 2\tabcolsep) * \real{0.7091}}@{}}
\toprule\noalign{}
\begin{minipage}[b]{\linewidth}\raggedright
Column
\end{minipage} &
\begin{minipage}[b]{\linewidth}\raggedright
Meaning
\end{minipage} \\
\midrule\noalign{}
\endhead
\bottomrule\noalign{}
\endlastfoot
\texttt{rank} /
\texttt{rank\_in\_table\_9}
/
\texttt{in\_table\_9}
& rank among 106,
rank among the 36
items of Table 9 and
membership of Table
9 \\
\texttt{metric\_name}
/
\texttt{canonical\_key}
/ \texttt{source} /
\texttt{input\_class}
/ \texttt{role} &
the metric and its
class \\
\texttt{n\_designs}
/ \texttt{rho} /
\texttt{rho\_ci\_lo90}
/
\texttt{rho\_ci\_hi90}
& 16 designs,
Spearman correlation
and 90\% interval
over designs \\
\texttt{interval\_excludes\_zero}
/
\texttt{higher\_value\_means}
& whether the
interval excludes 0
and whether a larger
value means higher
or lower activity \\
\texttt{rho\_given\_length}
/
\texttt{survives\_length\_control}
& correlation with
the sequence length
partialled out of
both sides, and
whether it
survives \\
\end{longtable}

}

The shared block of
12 provenance
columns (see
Conventions) is also
present and is not
repeated above.

\textbf{Excerpt}
(first 3 rows).

{\footnotesize

\begin{longtable}[]{@{}
  >{\raggedright\arraybackslash}p{(\columnwidth - 10\tabcolsep) * \real{0.1667}}
  >{\raggedright\arraybackslash}p{(\columnwidth - 10\tabcolsep) * \real{0.1667}}
  >{\raggedright\arraybackslash}p{(\columnwidth - 10\tabcolsep) * \real{0.1667}}
  >{\raggedright\arraybackslash}p{(\columnwidth - 10\tabcolsep) * \real{0.1667}}
  >{\raggedright\arraybackslash}p{(\columnwidth - 10\tabcolsep) * \real{0.1667}}
  >{\raggedright\arraybackslash}p{(\columnwidth - 10\tabcolsep) * \real{0.1667}}@{}}
\toprule\noalign{}
\begin{minipage}[b]{\linewidth}\raggedright
\texttt{rank}
\end{minipage} &
\begin{minipage}[b]{\linewidth}\raggedright
\texttt{metric\_name}
\end{minipage} &
\begin{minipage}[b]{\linewidth}\raggedright
\texttt{rho}
\end{minipage} &
\begin{minipage}[b]{\linewidth}\raggedright
\texttt{rho\_ci\_lo90}
\end{minipage} &
\begin{minipage}[b]{\linewidth}\raggedright
\texttt{rho\_ci\_hi90}
\end{minipage} &
\begin{minipage}[b]{\linewidth}\raggedright
\texttt{higher\_value\_means}
\end{minipage} \\
\midrule\noalign{}
\endhead
\bottomrule\noalign{}
\endlastfoot
\texttt{1} &
\texttt{frac\_I} &
\texttt{-0.7821} &
\texttt{-0.9494} &
\texttt{-0.5042} &
\texttt{lower\ activity} \\
\texttt{2} &
\texttt{frac\_A} &
\texttt{0.776} &
\texttt{0.4627} &
\texttt{0.9348} &
\texttt{higher\ activity} \\
\texttt{3} &
\texttt{rosetta\_omega}
& \texttt{-0.7756} &
\texttt{-0.8904} &
\texttt{-0.5} &
\texttt{lower\ activity} \\
\end{longtable}

}

\textbf{Notes.}

\begin{itemize}
\tightlist
\item
  Descriptive only;
  about 10\% of
  unrelated metrics
  are expected to
  have an interval
  that excludes 0.
\end{itemize}

\hypertarget{S22}{%
\subsubsection{Table
S22. Plated designs:
ordering among the
16 active designs,
structure-space
metrics and
comparators}\label{S22}}

\begin{shaded}\small

\textbf{File.}
\texttt{supplementary/S22\_plates\_ordering.csv}
· 81 rows × 23
columns · public

\textbf{Supports.}
Results 3.3.2;
Methods 2.4.2 (Eq.
15). \textbf{Cited
in.} Methods 2.4.2;
Results 3.3.2.

\textbf{A row is}
one metric of the
structure-space
panel or a
comparator (81 in
all).

\end{shaded}

The Spearman
correlation with log
kcat/KM among the 16
active designs with
and without the
sequence length
partialled out, with
90\% intervals and
whether the interval
excludes 0 (24 do).

\textbf{Key
columns.}

{\small

\begin{longtable}[]{@{}
  >{\raggedright\arraybackslash}p{(\columnwidth - 2\tabcolsep) * \real{0.2909}}
  >{\raggedright\arraybackslash}p{(\columnwidth - 2\tabcolsep) * \real{0.7091}}@{}}
\toprule\noalign{}
\begin{minipage}[b]{\linewidth}\raggedright
Column
\end{minipage} &
\begin{minipage}[b]{\linewidth}\raggedright
Meaning
\end{minipage} \\
\midrule\noalign{}
\endhead
\bottomrule\noalign{}
\endlastfoot
\texttt{metric\_name}
/
\texttt{canonical\_key}
/
\texttt{input\_class}
/
\texttt{in\_main\_set}
& the metric and its
class \\
\texttt{spearman\_rho}
/ \texttt{rho\_lo90}
/ \texttt{rho\_hi90}
/
\texttt{excludes\_zero}
& correlation, 90\%
interval and whether
it excludes 0 \\
\texttt{rho\_given\_length}
/
\texttt{survives\_length\_control}
& correlation with
length partialled
out and whether it
survives \\
\end{longtable}

}

The shared block of
12 provenance
columns (see
Conventions) is also
present and is not
repeated above.

\textbf{Excerpt}
(first 3 rows where
in\_main\_set = 1).

{\footnotesize

\begin{longtable}[]{@{}
  >{\raggedright\arraybackslash}p{(\columnwidth - 10\tabcolsep) * \real{0.1667}}
  >{\raggedright\arraybackslash}p{(\columnwidth - 10\tabcolsep) * \real{0.1667}}
  >{\raggedright\arraybackslash}p{(\columnwidth - 10\tabcolsep) * \real{0.1667}}
  >{\raggedright\arraybackslash}p{(\columnwidth - 10\tabcolsep) * \real{0.1667}}
  >{\raggedright\arraybackslash}p{(\columnwidth - 10\tabcolsep) * \real{0.1667}}
  >{\raggedright\arraybackslash}p{(\columnwidth - 10\tabcolsep) * \real{0.1667}}@{}}
\toprule\noalign{}
\begin{minipage}[b]{\linewidth}\raggedright
\texttt{metric\_name}
\end{minipage} &
\begin{minipage}[b]{\linewidth}\raggedright
\texttt{in\_main\_set}
\end{minipage} &
\begin{minipage}[b]{\linewidth}\raggedright
\texttt{spearman\_rho}
\end{minipage} &
\begin{minipage}[b]{\linewidth}\raggedright
\texttt{rho\_lo90}
\end{minipage} &
\begin{minipage}[b]{\linewidth}\raggedright
\texttt{rho\_hi90}
\end{minipage} &
\begin{minipage}[b]{\linewidth}\raggedright
\texttt{rho\_given\_length}
\end{minipage} \\
\midrule\noalign{}
\endhead
\bottomrule\noalign{}
\endlastfoot
\texttt{catalytic\_polar\_contacts}
& \texttt{1} &
\texttt{0.6711} &
\texttt{0.4115} &
\texttt{0.8501} &
\texttt{0.5183} \\
\texttt{rosetta\_catalytic\_vdw\_net}
& \texttt{1} &
\texttt{-0.6299} &
\texttt{-0.8777} &
\texttt{-0.2414} &
\texttt{-0.4743} \\
\texttt{rosetta\_cat\_fa\_sol}
& \texttt{1} &
\texttt{0.5651} &
\texttt{0.1543} &
\texttt{0.8487} &
\texttt{0.538} \\
\end{longtable}

}

\hypertarget{S23}{%
\subsubsection{Table
S23. Plated designs:
enrichment at every
keep fraction,
structure-space
metrics}\label{S23}}

\begin{shaded}\small

\textbf{File.}
\texttt{supplementary/S23\_plates\_enrichment.csv}
· 816 rows × 27
columns · public

\textbf{Supports.}
Results 3.3.2;
Methods 2.4.2.
\textbf{Cited in.}
Methods 2.4.2;
Results 3.3.2.

\textbf{A row is}
one metric, one
direction (keep the
highest or the
lowest values) and
one keep fraction
(816 rows).

\end{shaded}

Enrichment on the
192 plated designs
for the 102 metrics
of the
structure-space
panel and its
comparators, in both
directions and at
keep fractions of 5,
10, 25 and 50\%: the
number kept, the
actives kept and the
expected hits per
96-well plate with a
90\% interval
against the
unfiltered 8.0.

\textbf{Key
columns.}

{\small

\begin{longtable}[]{@{}
  >{\raggedright\arraybackslash}p{(\columnwidth - 2\tabcolsep) * \real{0.2909}}
  >{\raggedright\arraybackslash}p{(\columnwidth - 2\tabcolsep) * \real{0.7091}}@{}}
\toprule\noalign{}
\begin{minipage}[b]{\linewidth}\raggedright
Column
\end{minipage} &
\begin{minipage}[b]{\linewidth}\raggedright
Meaning
\end{minipage} \\
\midrule\noalign{}
\endhead
\bottomrule\noalign{}
\endlastfoot
\texttt{metric} /
\texttt{canonical\_key}
/
\texttt{input\_class}
/
\texttt{in\_main\_set}
& the metric \\
\texttt{direction} /
\texttt{keep\_fraction}
/ \texttt{n\_kept} &
keep the highest or
the lowest values;
the fraction and the
number of designs
kept \\
\texttt{n\_active\_kept}
/
\texttt{n\_active\_total}
& actives among
those kept, and
16 \\
\texttt{expected\_hits\_per\_plate}
/
\texttt{hits\_lo90}
/
\texttt{hits\_hi90}
& expected hits per
plate with 90\%
interval \\
\texttt{baseline\_hits\_per\_plate}
/
\texttt{enrichment\_vs\_baseline}
& 8.0 and the ratio
to it \\
\end{longtable}

}

The shared block of
12 provenance
columns (see
Conventions) is also
present and is not
repeated above.

\textbf{Excerpt}
(first 3 rows where
in\_main\_set = 1,
keep\_fraction =
0.25).

{\footnotesize

\begin{longtable}[]{@{}
  >{\raggedright\arraybackslash}p{(\columnwidth - 10\tabcolsep) * \real{0.1667}}
  >{\raggedright\arraybackslash}p{(\columnwidth - 10\tabcolsep) * \real{0.1667}}
  >{\raggedright\arraybackslash}p{(\columnwidth - 10\tabcolsep) * \real{0.1667}}
  >{\raggedright\arraybackslash}p{(\columnwidth - 10\tabcolsep) * \real{0.1667}}
  >{\raggedright\arraybackslash}p{(\columnwidth - 10\tabcolsep) * \real{0.1667}}
  >{\raggedright\arraybackslash}p{(\columnwidth - 10\tabcolsep) * \real{0.1667}}@{}}
\toprule\noalign{}
\begin{minipage}[b]{\linewidth}\raggedright
\texttt{metric}
\end{minipage} &
\begin{minipage}[b]{\linewidth}\raggedright
\texttt{direction}
\end{minipage} &
\begin{minipage}[b]{\linewidth}\raggedright
\texttt{keep\_fraction}
\end{minipage} &
\begin{minipage}[b]{\linewidth}\raggedright
\texttt{n\_kept}
\end{minipage} &
\begin{minipage}[b]{\linewidth}\raggedright
\texttt{n\_active\_kept}
\end{minipage} &
\begin{minipage}[b]{\linewidth}\raggedright
\texttt{expected\_hits\_per\_plate}
\end{minipage} \\
\midrule\noalign{}
\endhead
\bottomrule\noalign{}
\endlastfoot
\texttt{catalytic\_frac\_buried}
& \texttt{low} &
\texttt{0.25} &
\texttt{48} &
\texttt{2} &
\texttt{4.0} \\
\texttt{catalytic\_frac\_buried}
& \texttt{high} &
\texttt{0.25} &
\texttt{48} &
\texttt{1} &
\texttt{2.0} \\
\texttt{catalytic\_pairwise\_dist\_max}
& \texttt{low} &
\texttt{0.25} &
\texttt{48} &
\texttt{3} &
\texttt{6.0} \\
\end{longtable}

}

\textbf{Notes.}

\begin{itemize}
\tightlist
\item
  Descriptive only.
  Released as a CSV
  only for the full
  816 rows.
\end{itemize}

\hypertarget{S24}{%
\subsubsection{Table
S24. Plated designs:
enrichment for
prediction-based
metrics}\label{S24}}

\begin{shaded}\small

\textbf{File.}
\texttt{supplementary/S24\_plates\_prediction\_based\_enrichment.csv}
· 232 rows × 27
columns · public

\textbf{Supports.}
Results 3.3.2;
Methods 2.4.2.
\textbf{Cited in.}
Methods 2.4.2;
Results 3.3.2.

\textbf{A row is}
one prediction-based
metric, one
direction and one
keep fraction (232
rows).

\end{shaded}

The same as
\protect\hyperlink{S23}{S23}
for the 29
prediction-based
metrics (Chai-1 with
the transition-state
analogue and zinc;
active-site accuracy
against the design's
own model).

\textbf{Key
columns.}

{\small

\begin{longtable}[]{@{}
  >{\raggedright\arraybackslash}p{(\columnwidth - 2\tabcolsep) * \real{0.2909}}
  >{\raggedright\arraybackslash}p{(\columnwidth - 2\tabcolsep) * \real{0.7091}}@{}}
\toprule\noalign{}
\begin{minipage}[b]{\linewidth}\raggedright
Column
\end{minipage} &
\begin{minipage}[b]{\linewidth}\raggedright
Meaning
\end{minipage} \\
\midrule\noalign{}
\endhead
\bottomrule\noalign{}
\endlastfoot
\texttt{metric} /
\texttt{canonical\_key}
/
\texttt{input\_class}
/
\texttt{in\_main\_set}
& the metric \\
\texttt{direction} /
\texttt{keep\_fraction}
/ \texttt{n\_kept} /
\texttt{n\_active\_kept}
& as in
\protect\hyperlink{S23}{S23} \\
\texttt{expected\_hits\_per\_plate}
/
\texttt{hits\_lo90}
/
\texttt{hits\_hi90}
/
\texttt{enrichment\_vs\_baseline}
& expected hits per
plate, 90\% interval
and ratio to 8.0 \\
\end{longtable}

}

The shared block of
12 provenance
columns (see
Conventions) is also
present and is not
repeated above.

\textbf{Excerpt}
(first 3 rows where
keep\_fraction =
0.25, direction =
high).

{\footnotesize

\begin{longtable}[]{@{}
  >{\raggedright\arraybackslash}p{(\columnwidth - 8\tabcolsep) * \real{0.2000}}
  >{\raggedright\arraybackslash}p{(\columnwidth - 8\tabcolsep) * \real{0.2000}}
  >{\raggedright\arraybackslash}p{(\columnwidth - 8\tabcolsep) * \real{0.2000}}
  >{\raggedright\arraybackslash}p{(\columnwidth - 8\tabcolsep) * \real{0.2000}}
  >{\raggedright\arraybackslash}p{(\columnwidth - 8\tabcolsep) * \real{0.2000}}@{}}
\toprule\noalign{}
\begin{minipage}[b]{\linewidth}\raggedright
\texttt{metric}
\end{minipage} &
\begin{minipage}[b]{\linewidth}\raggedright
\texttt{direction}
\end{minipage} &
\begin{minipage}[b]{\linewidth}\raggedright
\texttt{keep\_fraction}
\end{minipage} &
\begin{minipage}[b]{\linewidth}\raggedright
\texttt{n\_active\_kept}
\end{minipage} &
\begin{minipage}[b]{\linewidth}\raggedright
\texttt{expected\_hits\_per\_plate}
\end{minipage} \\
\midrule\noalign{}
\endhead
\bottomrule\noalign{}
\endlastfoot
\texttt{chai\_plddt\_mean}
& \texttt{high} &
\texttt{0.25} &
\texttt{8} &
\texttt{16.0} \\
\texttt{chai\_ptm\_mean}
& \texttt{high} &
\texttt{0.25} &
\texttt{9} &
\texttt{18.0} \\
\texttt{frac\_Y} &
\texttt{high} &
\texttt{0.25} &
\texttt{12} &
\texttt{24.0} \\
\end{longtable}

}

\textbf{Notes.}

\begin{itemize}
\tightlist
\item
  Descriptive only
  (16 active
  designs).
\end{itemize}

\hypertarget{S25}{%
\subsubsection{Table
S25. Sanity floors
by class of
metric}\label{S25}}

\begin{shaded}\small

\textbf{File.}
\texttt{supplementary/S25\_sanity\_floor\_by\_input\_class.csv}
· 53 rows × 8
columns · public

\textbf{Supports.}
Methods 2.5
(detection floor);
Results 3.4.1.
\textbf{Cited in.}
Methods 2.5; Results
3.4.1.

\textbf{A row is}
one sanity-floor
condition on one
data set and one
class of metric (53
rows).

\end{shaded}

The sanity-floor
conditions (all
catalytic residues
replaced by Gly,
scrambled sequence
on the native
backbone, unrelated
protein, ligand
removed on the pilot
set) for the
143-enzyme set, the
shortened-motif set,
the pilot set and
the de novo set: how
many metrics of each
class were scored
and how many
responded. One row
reconciles the
detection floor with
the Gly step: the
responses are
identical on 98 of
98 metrics (the same
change scored
twice).

\textbf{Key
columns.}

{\small

\begin{longtable}[]{@{}
  >{\raggedright\arraybackslash}p{(\columnwidth - 2\tabcolsep) * \real{0.2909}}
  >{\raggedright\arraybackslash}p{(\columnwidth - 2\tabcolsep) * \real{0.7091}}@{}}
\toprule\noalign{}
\begin{minipage}[b]{\linewidth}\raggedright
Column
\end{minipage} &
\begin{minipage}[b]{\linewidth}\raggedright
Meaning
\end{minipage} \\
\midrule\noalign{}
\endhead
\bottomrule\noalign{}
\endlastfoot
\texttt{dataset} /
\texttt{condition} /
\texttt{input\_class}
& the data set, the
condition and the
class of metric \\
\texttt{n\_metrics}
/ \texttt{n\_scored}
/
\texttt{n\_responds}
/
\texttt{n\_unscored}
& metrics in the
class, scored,
responding and not
scored \\
\texttt{note} &
remarks, including
the
reconciliation \\
\end{longtable}

}

\textbf{Excerpt}
(first 3 rows where
dataset = main set
(143), condition =
cat\_to\_gly).

{\footnotesize

\begin{longtable}[]{@{}
  >{\raggedright\arraybackslash}p{(\columnwidth - 10\tabcolsep) * \real{0.1667}}
  >{\raggedright\arraybackslash}p{(\columnwidth - 10\tabcolsep) * \real{0.1667}}
  >{\raggedright\arraybackslash}p{(\columnwidth - 10\tabcolsep) * \real{0.1667}}
  >{\raggedright\arraybackslash}p{(\columnwidth - 10\tabcolsep) * \real{0.1667}}
  >{\raggedright\arraybackslash}p{(\columnwidth - 10\tabcolsep) * \real{0.1667}}
  >{\raggedright\arraybackslash}p{(\columnwidth - 10\tabcolsep) * \real{0.1667}}@{}}
\toprule\noalign{}
\begin{minipage}[b]{\linewidth}\raggedright
\texttt{dataset}
\end{minipage} &
\begin{minipage}[b]{\linewidth}\raggedright
\texttt{condition}
\end{minipage} &
\begin{minipage}[b]{\linewidth}\raggedright
\texttt{input\_class}
\end{minipage} &
\begin{minipage}[b]{\linewidth}\raggedright
\texttt{n\_metrics}
\end{minipage} &
\begin{minipage}[b]{\linewidth}\raggedright
\texttt{n\_scored}
\end{minipage} &
\begin{minipage}[b]{\linewidth}\raggedright
\texttt{n\_responds}
\end{minipage} \\
\midrule\noalign{}
\endhead
\bottomrule\noalign{}
\endlastfoot
\texttt{main\ set\ (143)}
&
\texttt{cat\_to\_gly}
&
\texttt{bookkeeping}
& \texttt{20} &
\texttt{20} &
\texttt{8} \\
\texttt{main\ set\ (143)}
&
\texttt{cat\_to\_gly}
& \texttt{seq\_only}
& \texttt{29} &
\texttt{29} &
\texttt{10} \\
\texttt{main\ set\ (143)}
&
\texttt{cat\_to\_gly}
&
\texttt{struct\_site}
& \texttt{34} &
\texttt{26} &
\texttt{25} \\
\end{longtable}

}

\hypertarget{S26}{%
\subsubsection{Table
S26. Control
matching: the
loosest tier
needed}\label{S26}}

\begin{shaded}\small

\textbf{File.}
\texttt{supplementary/S26\_control\_match\_tiers.csv}
· 81 rows × 6
columns · public

\textbf{Supports.}
Methods 2.5; Results
3.4.1. \textbf{Cited
in.} Methods 2.5;
Results 3.4.1;
Discussion.

\textbf{A row is}
one analysis, lesion
test and matching
tier (81 rows).

\end{shaded}

For each analysis
and lesion test, how
many (enzyme, step)
entries needed each
loosest matching
tier: 0 = exact
match, tiers 3 and
above drop packing
matching, tier 5
also relaxes burial
(main set only), -1
= no control found;
and the share of the
lesion test.

\textbf{Key
columns.}

{\small

\begin{longtable}[]{@{}
  >{\raggedright\arraybackslash}p{(\columnwidth - 2\tabcolsep) * \real{0.2909}}
  >{\raggedright\arraybackslash}p{(\columnwidth - 2\tabcolsep) * \real{0.7091}}@{}}
\toprule\noalign{}
\begin{minipage}[b]{\linewidth}\raggedright
Column
\end{minipage} &
\begin{minipage}[b]{\linewidth}\raggedright
Meaning
\end{minipage} \\
\midrule\noalign{}
\endhead
\bottomrule\noalign{}
\endlastfoot
\texttt{analysis\_set}
/
\texttt{lesion\_test}
& the analysis (data
set and controls)
and the lesion test
(catalytic lesion,
second-shell lesion
or the retired
deformation
experiment) \\
\texttt{worst\_tier}
/
\texttt{n\_entries}
/
\texttt{share\_of\_lesion\_test}
& loosest tier used,
number of entries
and their share of
the lesion test \\
\texttt{note} &
definition of the
tiers \\
\end{longtable}

}

\textbf{Excerpt}
(first 3 rows where
analysis\_set = main
set, packing-matched
controls,
lesion\_test =
catalytic lesion).

{\footnotesize

\begin{longtable}[]{@{}
  >{\raggedright\arraybackslash}p{(\columnwidth - 8\tabcolsep) * \real{0.2000}}
  >{\raggedright\arraybackslash}p{(\columnwidth - 8\tabcolsep) * \real{0.2000}}
  >{\raggedright\arraybackslash}p{(\columnwidth - 8\tabcolsep) * \real{0.2000}}
  >{\raggedright\arraybackslash}p{(\columnwidth - 8\tabcolsep) * \real{0.2000}}
  >{\raggedright\arraybackslash}p{(\columnwidth - 8\tabcolsep) * \real{0.2000}}@{}}
\toprule\noalign{}
\begin{minipage}[b]{\linewidth}\raggedright
\texttt{analysis\_set}
\end{minipage} &
\begin{minipage}[b]{\linewidth}\raggedright
\texttt{lesion\_test}
\end{minipage} &
\begin{minipage}[b]{\linewidth}\raggedright
\texttt{worst\_tier}
\end{minipage} &
\begin{minipage}[b]{\linewidth}\raggedright
\texttt{n\_entries}
\end{minipage} &
\begin{minipage}[b]{\linewidth}\raggedright
\texttt{share\_of\_lesion\_test}
\end{minipage} \\
\midrule\noalign{}
\endhead
\bottomrule\noalign{}
\endlastfoot
\texttt{main\ set,\ packing-matched\ controls}
&
\texttt{catalytic\ lesion}
& \texttt{-1.0} &
\texttt{8} &
\texttt{0.0035} \\
\texttt{main\ set,\ packing-matched\ controls}
&
\texttt{catalytic\ lesion}
& \texttt{0.0} &
\texttt{226} &
\texttt{0.0997} \\
\texttt{main\ set,\ packing-matched\ controls}
&
\texttt{catalytic\ lesion}
& \texttt{2.0} &
\texttt{496} &
\texttt{0.2189} \\
\end{longtable}

}

\hypertarget{S27}{%
\subsubsection{Table
S27. Control
matching: balance of
residues changed in
the two
arms}\label{S27}}

\begin{shaded}\small

\textbf{File.}
\texttt{supplementary/S27\_control\_arm\_balance.csv}
· 15 rows × 9
columns · public

\textbf{Supports.}
Methods 2.5; Results
3.4.1. \textbf{Cited
in.} Methods 2.5;
Results 3.4.1;
Discussion.

\textbf{A row is}
one analysis and
lesion test (15
rows).

\end{shaded}

The share of
(enzyme, step) pairs
in which the control
arm changed fewer or
more residues than
the catalytic arm,
the median and mean
ratio, and the
number with no
control. The
comparison of
medians used
elsewhere cannot see
this imbalance.

\textbf{Key
columns.}

{\small

\begin{longtable}[]{@{}
  >{\raggedright\arraybackslash}p{(\columnwidth - 2\tabcolsep) * \real{0.2909}}
  >{\raggedright\arraybackslash}p{(\columnwidth - 2\tabcolsep) * \real{0.7091}}@{}}
\toprule\noalign{}
\begin{minipage}[b]{\linewidth}\raggedright
Column
\end{minipage} &
\begin{minipage}[b]{\linewidth}\raggedright
Meaning
\end{minipage} \\
\midrule\noalign{}
\endhead
\bottomrule\noalign{}
\endlastfoot
\texttt{analysis\_set}
/
\texttt{lesion\_test}
/ \texttt{pairs} &
the analysis, the
lesion test and the
number of pairs \\
\texttt{control\_fewer\_share}
/
\texttt{control\_more\_share}
& shares with fewer
or more residues
changed in the
control arm \\
\texttt{median\_ratio}
/
\texttt{mean\_ratio}
/
\texttt{catalytic\_only\_no\_control}
& ratio of residues
changed and pairs
with no control \\
\end{longtable}

}

\textbf{Excerpt}
(first 3 rows).

{\footnotesize

\begin{longtable}[]{@{}
  >{\raggedright\arraybackslash}p{(\columnwidth - 10\tabcolsep) * \real{0.1667}}
  >{\raggedright\arraybackslash}p{(\columnwidth - 10\tabcolsep) * \real{0.1667}}
  >{\raggedright\arraybackslash}p{(\columnwidth - 10\tabcolsep) * \real{0.1667}}
  >{\raggedright\arraybackslash}p{(\columnwidth - 10\tabcolsep) * \real{0.1667}}
  >{\raggedright\arraybackslash}p{(\columnwidth - 10\tabcolsep) * \real{0.1667}}
  >{\raggedright\arraybackslash}p{(\columnwidth - 10\tabcolsep) * \real{0.1667}}@{}}
\toprule\noalign{}
\begin{minipage}[b]{\linewidth}\raggedright
\texttt{analysis\_set}
\end{minipage} &
\begin{minipage}[b]{\linewidth}\raggedright
\texttt{lesion\_test}
\end{minipage} &
\begin{minipage}[b]{\linewidth}\raggedright
\texttt{pairs}
\end{minipage} &
\begin{minipage}[b]{\linewidth}\raggedright
\texttt{control\_fewer\_share}
\end{minipage} &
\begin{minipage}[b]{\linewidth}\raggedright
\texttt{control\_more\_share}
\end{minipage} &
\begin{minipage}[b]{\linewidth}\raggedright
\texttt{median\_ratio}
\end{minipage} \\
\midrule\noalign{}
\endhead
\bottomrule\noalign{}
\endlastfoot
\texttt{pilot\ set}
&
\texttt{catalytic\ lesion}
& \texttt{212} &
\texttt{0.1981} &
\texttt{0.0} &
\texttt{1.0} \\
\texttt{pilot\ set}
&
\texttt{second-shell\ lesion}
& \texttt{159} &
\texttt{0.0} &
\texttt{0.0} &
\texttt{1.0} \\
\texttt{pilot\ set}
&
\texttt{deformation}
& \texttt{205} &
\texttt{0.0} &
\texttt{0.0} &
\texttt{1.0} \\
\end{longtable}

}

\hypertarget{S28}{%
\subsubsection{Table
S28. Catalytic
lesion: every metric
scored (Table 10
plus supplementary
comparators and
per-step
flags)}\label{S28}}

\begin{shaded}\small

\textbf{File.}
\texttt{supplementary/S28\_catalytic\_lesion\_all\_metrics\_ranked.csv}
· 43 rows × 73
columns · public

\textbf{Supports.}
Results 3.4.1; the
complete version of
Table 10.
\textbf{Cited in.}
Methods 2.5.1;
Results 3.4.1;
Discussion.

\textbf{A row is}
one metric scored on
the catalytic-lesion
ladder (43 in all).

\end{shaded}

Table 10 with the 8
supplementary
comparators and
every flag that
decides whether a
step counts. Metrics
are grouped by the
lesion step at which
they are first
detected and ranked
by the mean of the
rank statistic R
over the counted
steps. For each step
(isosteric 15.3,
non-isosteric 36.4,
Ala 53.2, Gly 80.5
A3) it gives R with
interval, the number
of informative
proteins and the
flags that the
metric responds,
that the control
response is above
its noise and that
the metric is
specific; the number
of steps counted;
the first step
detected and the
first step at which
the metric is
specific (the onset
of specificity); the
mean R of the 143
main enzymes alone;
and, for the
prediction-based
metrics, the
distance-matched R.
Structure-space
metrics: 195 pooled
natural enzymes (143
main + 52 pilot);
prediction-based: 59
enzymes, 286 pairs.

\textbf{Key
columns.}

{\small

\begin{longtable}[]{@{}
  >{\raggedright\arraybackslash}p{(\columnwidth - 2\tabcolsep) * \real{0.2909}}
  >{\raggedright\arraybackslash}p{(\columnwidth - 2\tabcolsep) * \real{0.7091}}@{}}
\toprule\noalign{}
\begin{minipage}[b]{\linewidth}\raggedright
Column
\end{minipage} &
\begin{minipage}[b]{\linewidth}\raggedright
Meaning
\end{minipage} \\
\midrule\noalign{}
\endhead
\bottomrule\noalign{}
\endlastfoot
\texttt{detection\_group}
/
\texttt{rank\_in\_group}
/ \texttt{rank} /
\texttt{rank\_in\_table\_10}
& group by first
detected step, rank
in the group,
overall rank and
rank among the 35
items of Table 10 \\
\texttt{metric\_name}
/
\texttt{canonical\_key}
/ \texttt{source} /
\texttt{input\_class}
/ \texttt{role} /
\texttt{in\_ranked\_set}
& the metric, class,
role and whether it
is one of the 35
ranked items \\
\texttt{n\_steps} /
\texttt{n\_steps\_interpretable}
& steps with a value
and steps counted \\
\texttt{r\_mean} /
\texttt{r\_lo} /
\texttt{r\_hi} &
mean R over counted
steps with 95\%
interval over enzyme
sub-subclasses \\
r\_ / \_lo /
\emph{hi / n} & R,
interval and
informative proteins
for isosteric,
non\_isosteric, ala
and gly \\
responds\_ /
control\_above\_noise\_
/ specific\_ & the
per-step flags \\
\texttt{first\_step\_detected}
/
\texttt{first\_step\_specific}
& onsets of
detection and of
specificity \\
\texttt{r\_matched}
/
\texttt{n\_pairs\_matched}
/
\texttt{n\_systems\_matched}
& distance-matched R
and its size
(prediction-based
metrics) \\
\texttt{r\_mean\_primary\_only}
/
\texttt{above\_best\_of\_all\_quantities\_threshold}
& mean R of the 143
main enzymes and the
best-of-N flag \\
\end{longtable}

}

The shared block of
12 provenance
columns (see
Conventions) is also
present and is not
repeated above.

\textbf{Excerpt}
(first 3 rows).

{\footnotesize

\begin{longtable}[]{@{}
  >{\raggedright\arraybackslash}p{(\columnwidth - 10\tabcolsep) * \real{0.1667}}
  >{\raggedright\arraybackslash}p{(\columnwidth - 10\tabcolsep) * \real{0.1667}}
  >{\raggedright\arraybackslash}p{(\columnwidth - 10\tabcolsep) * \real{0.1667}}
  >{\raggedright\arraybackslash}p{(\columnwidth - 10\tabcolsep) * \real{0.1667}}
  >{\raggedright\arraybackslash}p{(\columnwidth - 10\tabcolsep) * \real{0.1667}}
  >{\raggedright\arraybackslash}p{(\columnwidth - 10\tabcolsep) * \real{0.1667}}@{}}
\toprule\noalign{}
\begin{minipage}[b]{\linewidth}\raggedright
\texttt{detection\_group}
\end{minipage} &
\begin{minipage}[b]{\linewidth}\raggedright
\texttt{rank\_in\_group}
\end{minipage} &
\begin{minipage}[b]{\linewidth}\raggedright
\texttt{metric\_name}
\end{minipage} &
\begin{minipage}[b]{\linewidth}\raggedright
\texttt{r\_mean}
\end{minipage} &
\begin{minipage}[b]{\linewidth}\raggedright
\texttt{r\_lo}
\end{minipage} &
\begin{minipage}[b]{\linewidth}\raggedright
\texttt{r\_hi}
\end{minipage} \\
\midrule\noalign{}
\endhead
\bottomrule\noalign{}
\endlastfoot
\texttt{isosteric} &
\texttt{1} &
\texttt{chai\_per\_chain\_ptm\_min\_mean}
& \texttt{0.7524} &
\texttt{0.6768} &
\texttt{0.8238} \\
\texttt{isosteric} &
\texttt{2} &
\texttt{propka\_catalytic\_pka\_range}
& \texttt{0.6834} &
\texttt{0.5759} &
\texttt{0.7806} \\
\texttt{isosteric} &
\texttt{3} &
\texttt{chai\_iptm\_mean}
& \texttt{0.6732} &
\texttt{0.5991} &
\texttt{0.7463} \\
\end{longtable}

}

\textbf{Notes.}

\begin{itemize}
\tightlist
\item
  Descriptive of a
  pooled set.
  \protect\hyperlink{S29}{S29}
  and
  \protect\hyperlink{S30}{S30}
  hold the
  specificity-ratio
  analysis.
\end{itemize}

\hypertarget{S29}{%
\subsubsection{Table
S29. Structure-space
specificity by
metric family and
lesion step, with
the panel statistic
and sensitivity
sets}\label{S29}}

\begin{shaded}\small

\textbf{File.}
\texttt{supplementary/S29\_specificity\_by\_family\_and\_step.csv}
· 20 rows × 27
columns · public

\textbf{Supports.}
Results 3.4.1 and
3.4.2; Discussion.
\textbf{Cited in.}
Methods 2.5.1;
Results 3.4.1,
3.4.2; Discussion.

\textbf{A row is}
one family of
metrics at one
lesion step or dose
(20 rows).

\end{shaded}

The number of
eligible, responding
and specific
structure-space
metrics (as distinct
metrics) by family
(site-scoped Rosetta
energy, catalytic
geometry, catalytic
pKa, all main-set
metrics) at each
step of the
catalytic lesion and
each second-shell
dose, the
geometric-mean
specificity ratio of
the main-set panel
(isosteric, 1.075
{[}0.958, 1.178{]}),
and the counts in
four sensitivity
sets (shortened
motif, no packing
matching or count
gate, pilot set, de
novo designs). At
the isosteric step
23 of 29 metrics
respond and none is
specific;
specificity appears
for 5 metrics at the
non-isosteric step
and for 8 at the Ala
and Gly steps.

\textbf{Key
columns.}

{\small

\begin{longtable}[]{@{}
  >{\raggedright\arraybackslash}p{(\columnwidth - 2\tabcolsep) * \real{0.2909}}
  >{\raggedright\arraybackslash}p{(\columnwidth - 2\tabcolsep) * \real{0.7091}}@{}}
\toprule\noalign{}
\begin{minipage}[b]{\linewidth}\raggedright
Column
\end{minipage} &
\begin{minipage}[b]{\linewidth}\raggedright
Meaning
\end{minipage} \\
\midrule\noalign{}
\endhead
\bottomrule\noalign{}
\endlastfoot
\texttt{family} /
\texttt{lesion\_test}
/ \texttt{level} /
\texttt{dose\_value}
/
\texttt{dose\_unit}
& family, lesion
test (catalytic
lesion or
second-shell
lesion), step or
dose and its size \\
\texttt{n\_eligible\_keys}
/
\texttt{n\_responds\_keys}
/
\texttt{n\_specific\_keys}
/ \texttt{n\_cells}
& counts as distinct
metrics, and
cells \\
\texttt{sr\_geomean\_ci90}
& geometric-mean
specificity ratio of
the panel with 90\%
interval \\
\texttt{sens\_shortened\_motif}
/
\texttt{sens\_no\_packing\_match}
/
\texttt{sens\_pilot\_set}
/
\texttt{sens\_de\_novo\_set}
& specific /
eligible in the
shortened-motif,
no-packing-match,
pilot and de novo
sets \\
\end{longtable}

}

The shared block of
12 provenance
columns (see
Conventions) is also
present and is not
repeated above.

\textbf{Excerpt}
(first 3 rows where
lesion\_test =
catalytic lesion).

{\footnotesize

\begin{longtable}[]{@{}
  >{\raggedright\arraybackslash}p{(\columnwidth - 10\tabcolsep) * \real{0.1667}}
  >{\raggedright\arraybackslash}p{(\columnwidth - 10\tabcolsep) * \real{0.1667}}
  >{\raggedright\arraybackslash}p{(\columnwidth - 10\tabcolsep) * \real{0.1667}}
  >{\raggedright\arraybackslash}p{(\columnwidth - 10\tabcolsep) * \real{0.1667}}
  >{\raggedright\arraybackslash}p{(\columnwidth - 10\tabcolsep) * \real{0.1667}}
  >{\raggedright\arraybackslash}p{(\columnwidth - 10\tabcolsep) * \real{0.1667}}@{}}
\toprule\noalign{}
\begin{minipage}[b]{\linewidth}\raggedright
\texttt{family}
\end{minipage} &
\begin{minipage}[b]{\linewidth}\raggedright
\texttt{level}
\end{minipage} &
\begin{minipage}[b]{\linewidth}\raggedright
\texttt{n\_eligible\_keys}
\end{minipage} &
\begin{minipage}[b]{\linewidth}\raggedright
\texttt{n\_responds\_keys}
\end{minipage} &
\begin{minipage}[b]{\linewidth}\raggedright
\texttt{n\_specific\_keys}
\end{minipage} &
\begin{minipage}[b]{\linewidth}\raggedright
\texttt{sens\_shortened\_motif}
\end{minipage} \\
\midrule\noalign{}
\endhead
\bottomrule\noalign{}
\endlastfoot
\texttt{Rosetta\ energy\ (site-scoped)}
& \texttt{isosteric}
& \texttt{15} &
\texttt{13} &
\texttt{0} &
\texttt{0\ of\ 15} \\
\texttt{Catalytic\ geometry\ (site-scoped)}
& \texttt{isosteric}
& \texttt{8} &
\texttt{5} &
\texttt{0} &
\texttt{0\ of\ 8} \\
\texttt{Catalytic\ pKa}
& \texttt{isosteric}
& \texttt{6} &
\texttt{5} &
\texttt{0} &
\texttt{} \\
\end{longtable}

}

\hypertarget{S30}{%
\subsubsection{Table
S30. Re-predicted
catalytic lesion:
the prediction-based
metrics under three
successive controls
and the
pre-specified
rule}\label{S30}}

\begin{shaded}\small

\textbf{File.}
\texttt{supplementary/S30\_predictor\_ladder\_three\_controls.csv}
· 6 rows × 23
columns · public

\textbf{Supports.}
Results 3.4.1;
Methods 2.5.1 (Eq.
24). \textbf{Cited
in.} Methods 2.5.1;
Results 3.4.1.

\textbf{A row is}
one prediction-based
metric or reference
row (6 rows).

\end{shaded}

Specificity of the
prediction-based
metrics on the
natural arm of the
re-predicted ladder
(59 enzymes, 286
pairs): the
specificity ratio on
the isosteric step
with the original
control, in the
distance-matched
subset (68 pairs, 21
systems) and after
regression
adjustment on all
steps (with the
unadjusted ratio on
the same pairs), and
the label of the
rule specified
before the analysis.
The ratio of the
interface terms
falls once the
control is matched
or adjusted on the
distance to the
ligand (ipTM: 1.96
with the original
control, 0.47 in the
distance-matched
subset and 1.18
after adjustment);
no interface term
meets the rule.

\textbf{Key
columns.}

{\small

\begin{longtable}[]{@{}
  >{\raggedright\arraybackslash}p{(\columnwidth - 2\tabcolsep) * \real{0.2909}}
  >{\raggedright\arraybackslash}p{(\columnwidth - 2\tabcolsep) * \real{0.7091}}@{}}
\toprule\noalign{}
\begin{minipage}[b]{\linewidth}\raggedright
Column
\end{minipage} &
\begin{minipage}[b]{\linewidth}\raggedright
Meaning
\end{minipage} \\
\midrule\noalign{}
\endhead
\bottomrule\noalign{}
\endlastfoot
\texttt{metric} /
\texttt{role} & the
metric and main or
reference \\
\texttt{n\_pairs\_isosteric}
/
\texttt{n\_systems\_isosteric}
& pairs and systems
at the isosteric
step \\
\texttt{sr\_mean\_isosteric\_original\_control}
& specificity ratio
with the original
control \\
\texttt{distance\_matched\_k\_of\_N}
/
\texttt{sr\_distance\_matched}
& size of the
distance-matched
subset and the ratio
in it \\
\texttt{unadjusted\_gmr\_all\_steps}
/
\texttt{adjusted\_gmr\_all\_steps}
& geometric-mean
ratio over all
steps, unadjusted
and adjusted for
distance and
burial \\
\texttt{rule\_label}
& the label from the
pre-specified rule
(lower 90\% bound
above 1.2 with at
least 30 pairs in 15
systems) \\
\end{longtable}

}

The shared block of
12 provenance
columns (see
Conventions) is also
present and is not
repeated above.

\textbf{Excerpt}
(first 3 rows).

{\footnotesize

\begin{longtable}[]{@{}
  >{\raggedright\arraybackslash}p{(\columnwidth - 8\tabcolsep) * \real{0.2000}}
  >{\raggedright\arraybackslash}p{(\columnwidth - 8\tabcolsep) * \real{0.2000}}
  >{\raggedright\arraybackslash}p{(\columnwidth - 8\tabcolsep) * \real{0.2000}}
  >{\raggedright\arraybackslash}p{(\columnwidth - 8\tabcolsep) * \real{0.2000}}
  >{\raggedright\arraybackslash}p{(\columnwidth - 8\tabcolsep) * \real{0.2000}}@{}}
\toprule\noalign{}
\begin{minipage}[b]{\linewidth}\raggedright
\texttt{metric}
\end{minipage} &
\begin{minipage}[b]{\linewidth}\raggedright
\texttt{sr\_mean\_isosteric\_original\_control}
\end{minipage} &
\begin{minipage}[b]{\linewidth}\raggedright
\texttt{sr\_distance\_matched}
\end{minipage} &
\begin{minipage}[b]{\linewidth}\raggedright
\texttt{adjusted\_gmr\_all\_steps}
\end{minipage} &
\begin{minipage}[b]{\linewidth}\raggedright
\texttt{rule\_label}
\end{minipage} \\
\midrule\noalign{}
\endhead
\bottomrule\noalign{}
\endlastfoot
\texttt{ame\_crystal\_clash\_max}
&
\texttt{1.84\ {[}1.31,\ 2.69{]}}
&
\texttt{1.05\ {[}0.79,\ 1.39{]}}
&
\texttt{1.18\ {[}0.95,\ 1.46{]}}
&
\texttt{testable;\ not\ specific\ (lower\ 90\%…} \\
\texttt{ame\_crystal\_rmsd\_min}
&
\texttt{2.67\ {[}1.02,\ 4.74{]}}
&
\texttt{0.98\ {[}0.68,\ 1.21{]}}
&
\texttt{1.90\ {[}1.40,\ 2.53{]}}
&
\texttt{knockon\_displacement\_no\_specifici…} \\
\texttt{chai\_iptm\_mean}
&
\texttt{1.96\ {[}1.44,\ 2.71{]}}
&
\texttt{0.47\ {[}0.24,\ 1.04{]}}
&
\texttt{1.18\ {[}0.85,\ 1.60{]}}
&
\texttt{testable;\ not\ specific\ (lower\ 90\%…} \\
\end{longtable}

}

\hypertarget{S31}{%
\subsubsection{Table
S31. Sensitivity of
the
specificity-ratio
counts to motif
size, packing
matching and count
gating}\label{S31}}

\begin{shaded}\small

\textbf{File.}
\texttt{supplementary/S31\_motif\_packing\_sensitivity.csv}
· 20 rows × 7
columns · public

\textbf{Supports.}
Results 3.4.1;
Discussion
(limitations).
\textbf{Cited in.}
Methods 2.5.1;
Results 3.4.1;
Discussion.

\textbf{A row is}
one analysis and one
lesion step (20
rows).

\end{shaded}

Counts of eligible,
responding and
specific metrics at
each step for five
analyses: the pilot
set, the main set
without packing
matching or count
gating, the main set
with packing-matched
controls (primary),
the main set with
the motif shortened
to at most 3
residues, and the 30
de novo designs with
a created-site
control.

\textbf{Key
columns.}

{\small

\begin{longtable}[]{@{}
  >{\raggedright\arraybackslash}p{(\columnwidth - 2\tabcolsep) * \real{0.2909}}
  >{\raggedright\arraybackslash}p{(\columnwidth - 2\tabcolsep) * \real{0.7091}}@{}}
\toprule\noalign{}
\begin{minipage}[b]{\linewidth}\raggedright
Column
\end{minipage} &
\begin{minipage}[b]{\linewidth}\raggedright
Meaning
\end{minipage} \\
\midrule\noalign{}
\endhead
\bottomrule\noalign{}
\endlastfoot
\texttt{reader\_name}
/
\texttt{motif\_and\_controls}
& the analysis and
how its motif and
controls are
defined \\
\texttt{level} &
lesion step \\
\texttt{n\_eligible\_keys}
/
\texttt{n\_responds\_keys}
/
\texttt{n\_specific\_keys}
/ \texttt{n\_cells}
& counts as distinct
metrics, and
cells \\
\end{longtable}

}

\textbf{Excerpt}
(first 3 rows).

{\footnotesize

\begin{longtable}[]{@{}
  >{\raggedright\arraybackslash}p{(\columnwidth - 8\tabcolsep) * \real{0.2000}}
  >{\raggedright\arraybackslash}p{(\columnwidth - 8\tabcolsep) * \real{0.2000}}
  >{\raggedright\arraybackslash}p{(\columnwidth - 8\tabcolsep) * \real{0.2000}}
  >{\raggedright\arraybackslash}p{(\columnwidth - 8\tabcolsep) * \real{0.2000}}
  >{\raggedright\arraybackslash}p{(\columnwidth - 8\tabcolsep) * \real{0.2000}}@{}}
\toprule\noalign{}
\begin{minipage}[b]{\linewidth}\raggedright
\texttt{reader\_name}
\end{minipage} &
\begin{minipage}[b]{\linewidth}\raggedright
\texttt{level}
\end{minipage} &
\begin{minipage}[b]{\linewidth}\raggedright
\texttt{n\_eligible\_keys}
\end{minipage} &
\begin{minipage}[b]{\linewidth}\raggedright
\texttt{n\_responds\_keys}
\end{minipage} &
\begin{minipage}[b]{\linewidth}\raggedright
\texttt{n\_specific\_keys}
\end{minipage} \\
\midrule\noalign{}
\endhead
\bottomrule\noalign{}
\endlastfoot
\texttt{pilot\ set}
& \texttt{isosteric}
& \texttt{23} &
\texttt{12} &
\texttt{0} \\
\texttt{pilot\ set}
&
\texttt{non\_isosteric}
& \texttt{23} &
\texttt{12} &
\texttt{0} \\
\texttt{pilot\ set}
& \texttt{ala} &
\texttt{23} &
\texttt{14} &
\texttt{0} \\
\end{longtable}

}

\hypertarget{S32}{%
\subsubsection{Table
S32. Re-predicted
ladder: seed-noise
context}\label{S32}}

\begin{shaded}\small

\textbf{File.}
\texttt{supplementary/S32\_replicate\_noise\_context.csv}
· 56 rows × 12
columns · public

\textbf{Supports.}
Methods 2.5.1;
Results 3.4.1.
\textbf{Cited in.}
Methods 2.5.1;
Results 3.4.1.

\textbf{A row is}
one prediction-based
metric in one arm
and subset (56
rows).

\end{shaded}

Noise context from
the second-seed
replicate on 30
natural pairs: the
median absolute
change in each arm
against the seed
noise of the metric,
for the nine
prediction-based
metrics, on all
pairs and on the
original matched
subset.

\textbf{Key
columns.}

{\small

\begin{longtable}[]{@{}
  >{\raggedright\arraybackslash}p{(\columnwidth - 2\tabcolsep) * \real{0.2909}}
  >{\raggedright\arraybackslash}p{(\columnwidth - 2\tabcolsep) * \real{0.7091}}@{}}
\toprule\noalign{}
\begin{minipage}[b]{\linewidth}\raggedright
Column
\end{minipage} &
\begin{minipage}[b]{\linewidth}\raggedright
Meaning
\end{minipage} \\
\midrule\noalign{}
\endhead
\bottomrule\noalign{}
\endlastfoot
\texttt{source} /
\texttt{metric\_name}
/ \texttt{subset} /
\texttt{arm} &
natural or de novo,
the metric,
all\_pairs or
M0\_matched, and
catalytic or
control \\
\texttt{k\_pairs} /
\texttt{N\_pairs} /
\texttt{k\_systems}
/
\texttt{N\_systems}
& pairs and
systems \\
\texttt{median\_abs\_delta}
/
\texttt{median\_sigma\_native}
/
\texttt{frac\_abs\_delta\_gt\_sigma\_native}
& median absolute
change, the metric's
seed noise and the
share of pairs whose
change exceeds it \\
\end{longtable}

}

\textbf{Excerpt}
(first 3 rows where
source = natural,
subset =
all\_pairs).

{\footnotesize

\begin{longtable}[]{@{}
  >{\raggedright\arraybackslash}p{(\columnwidth - 10\tabcolsep) * \real{0.1667}}
  >{\raggedright\arraybackslash}p{(\columnwidth - 10\tabcolsep) * \real{0.1667}}
  >{\raggedright\arraybackslash}p{(\columnwidth - 10\tabcolsep) * \real{0.1667}}
  >{\raggedright\arraybackslash}p{(\columnwidth - 10\tabcolsep) * \real{0.1667}}
  >{\raggedright\arraybackslash}p{(\columnwidth - 10\tabcolsep) * \real{0.1667}}
  >{\raggedright\arraybackslash}p{(\columnwidth - 10\tabcolsep) * \real{0.1667}}@{}}
\toprule\noalign{}
\begin{minipage}[b]{\linewidth}\raggedright
\texttt{metric\_name}
\end{minipage} &
\begin{minipage}[b]{\linewidth}\raggedright
\texttt{subset}
\end{minipage} &
\begin{minipage}[b]{\linewidth}\raggedright
\texttt{arm}
\end{minipage} &
\begin{minipage}[b]{\linewidth}\raggedright
\texttt{median\_abs\_delta}
\end{minipage} &
\begin{minipage}[b]{\linewidth}\raggedright
\texttt{median\_sigma\_native}
\end{minipage} &
\begin{minipage}[b]{\linewidth}\raggedright
\texttt{frac\_abs\_delta\_gt\_sigma\_native}
\end{minipage} \\
\midrule\noalign{}
\endhead
\bottomrule\noalign{}
\endlastfoot
\texttt{chai\_iptm\_mean}
&
\texttt{all\_pairs}
& \texttt{catalytic}
& \texttt{0.00666} &
\texttt{0.00154} &
\texttt{0.8292} \\
\texttt{chai\_iptm\_mean}
&
\texttt{all\_pairs}
& \texttt{control} &
\texttt{0.00254} &
\texttt{0.00154} &
\texttt{0.6797} \\
\texttt{chai\_aggregate\_score\_mean}
&
\texttt{all\_pairs}
& \texttt{catalytic}
& \texttt{0.00556} &
\texttt{0.00127} &
\texttt{0.8392} \\
\end{longtable}

}

\hypertarget{S33}{%
\subsubsection{Table
S33. Specificity
ratio of every
eligible metric at
every
step}\label{S33}}

\begin{shaded}\small

\textbf{File.}
\texttt{supplementary/S33\_eligible\_metrics.csv}
· 256 rows × 29
columns · public

\textbf{Supports.}
Methods 2.5.1 (Eq.
22); Results 3.4.1
and 3.4.2.
\textbf{Cited in.}
Methods 2.5.1;
Results 3.4.1,
3.4.2.

\textbf{A row is}
one eligible
(site-scoped) metric
at one lesion step
or dose (256 rows).

\end{shaded}

The median change at
the catalytic lesion
and its interval,
the specificity
ratio with 95\%
interval, whether
the metric responds
and whether it is
specific, and the
verdict: specific,
non\_specific, blind
(does not respond),
invariant (does not
change) or below the
floor of five units.
Steps of the
catalytic lesion are
isosteric,
non\_isosteric, ala,
gly; second-shell
doses are
second\_shell\_1,
\_2, \_4.

\textbf{Key
columns.}

{\small

\begin{longtable}[]{@{}
  >{\raggedright\arraybackslash}p{(\columnwidth - 2\tabcolsep) * \real{0.2909}}
  >{\raggedright\arraybackslash}p{(\columnwidth - 2\tabcolsep) * \real{0.7091}}@{}}
\toprule\noalign{}
\begin{minipage}[b]{\linewidth}\raggedright
Column
\end{minipage} &
\begin{minipage}[b]{\linewidth}\raggedright
Meaning
\end{minipage} \\
\midrule\noalign{}
\endhead
\bottomrule\noalign{}
\endlastfoot
\texttt{canonical\_key}
/
\texttt{metric\_name}
/
\texttt{input\_class}
/
\texttt{in\_main\_set}
& the metric
(aliases collapse in
canonical\_key) \\
\texttt{lesion\_test}
/ \texttt{level} /
\texttt{n\_systems}
& catalytic lesion
or second-shell
lesion, the step or
dose and the number
of enzymes \\
\texttt{delta\_median}
/
\texttt{delta\_ci\_lo}
/
\texttt{delta\_ci\_hi}
& median change at
the catalytic lesion
with interval \\
\texttt{sr\_median}
/
\texttt{sr\_ci\_lo}
/
\texttt{sr\_ci\_hi}
& specificity ratio
with 95\%
interval \\
\texttt{responds} /
\texttt{specific} /
\texttt{verdict} /
\texttt{applicability\_status}
& the flags, the
verdict and the
applicability
status \\
\end{longtable}

}

The shared block of
12 provenance
columns (see
Conventions) is also
present and is not
repeated above.

\textbf{Excerpt}
(first 3 rows where
in\_main\_set = 1,
level = ala).

{\footnotesize

\begin{longtable}[]{@{}
  >{\raggedright\arraybackslash}p{(\columnwidth - 10\tabcolsep) * \real{0.1667}}
  >{\raggedright\arraybackslash}p{(\columnwidth - 10\tabcolsep) * \real{0.1667}}
  >{\raggedright\arraybackslash}p{(\columnwidth - 10\tabcolsep) * \real{0.1667}}
  >{\raggedright\arraybackslash}p{(\columnwidth - 10\tabcolsep) * \real{0.1667}}
  >{\raggedright\arraybackslash}p{(\columnwidth - 10\tabcolsep) * \real{0.1667}}
  >{\raggedright\arraybackslash}p{(\columnwidth - 10\tabcolsep) * \real{0.1667}}@{}}
\toprule\noalign{}
\begin{minipage}[b]{\linewidth}\raggedright
\texttt{canonical\_key}
\end{minipage} &
\begin{minipage}[b]{\linewidth}\raggedright
\texttt{lesion\_test}
\end{minipage} &
\begin{minipage}[b]{\linewidth}\raggedright
\texttt{level}
\end{minipage} &
\begin{minipage}[b]{\linewidth}\raggedright
\texttt{n\_systems}
\end{minipage} &
\begin{minipage}[b]{\linewidth}\raggedright
\texttt{sr\_median}
\end{minipage} &
\begin{minipage}[b]{\linewidth}\raggedright
\texttt{verdict}
\end{minipage} \\
\midrule\noalign{}
\endhead
\bottomrule\noalign{}
\endlastfoot
\texttt{catalytic\_frac\_buried}
&
\texttt{catalytic\ lesion}
& \texttt{ala} &
\texttt{143} &
\texttt{1.9091} &
\texttt{specific} \\
\texttt{catalytic\_pairwise\_dist\_max}
&
\texttt{catalytic\ lesion}
& \texttt{ala} &
\texttt{124} &
\texttt{0.0} &
\texttt{blind} \\
\texttt{catalytic\_pairwise\_dist\_mean}
&
\texttt{catalytic\ lesion}
& \texttt{ala} &
\texttt{124} &
\texttt{0.7262} &
\texttt{non\_specific} \\
\end{longtable}

}

\hypertarget{S34}{%
\subsubsection{Table
S34. Specificity
ratio of
whole-protein and
sequence-only
comparators at every
step}\label{S34}}

\begin{shaded}\small

\textbf{File.}
\texttt{supplementary/S34\_global\_comparators.csv}
· 308 rows × 29
columns · public

\textbf{Supports.}
Methods 2.5.1;
Results 3.4.
\textbf{Cited in.}
Methods 2.5.1;
Results 3.4.1,
3.4.2.

\textbf{A row is}
one comparator
metric at one lesion
step or dose (308
rows).

\end{shaded}

The same cells as
\protect\hyperlink{S33}{S33}
for the
whole-protein and
sequence-only
metrics, which are
not in the main set.
They show that
composition metrics
have a ratio of 1 by
construction and
that whole-protein
energies move with
the lesion.

\textbf{Key
columns.}

{\small

\begin{longtable}[]{@{}
  >{\raggedright\arraybackslash}p{(\columnwidth - 2\tabcolsep) * \real{0.2909}}
  >{\raggedright\arraybackslash}p{(\columnwidth - 2\tabcolsep) * \real{0.7091}}@{}}
\toprule\noalign{}
\begin{minipage}[b]{\linewidth}\raggedright
Column
\end{minipage} &
\begin{minipage}[b]{\linewidth}\raggedright
Meaning
\end{minipage} \\
\midrule\noalign{}
\endhead
\bottomrule\noalign{}
\endlastfoot
\texttt{canonical\_key}
/
\texttt{metric\_name}
/
\texttt{input\_class}
& the comparator \\
\texttt{lesion\_test}
/ \texttt{level} /
\texttt{n\_systems}
& as in
\protect\hyperlink{S33}{S33} \\
delta\_\emph{,
sr\_}, responds,
specific, verdict &
as in
\protect\hyperlink{S33}{S33} \\
\end{longtable}

}

The shared block of
12 provenance
columns (see
Conventions) is also
present and is not
repeated above.

\textbf{Excerpt}
(first 3 rows where
level = ala).

{\footnotesize

\begin{longtable}[]{@{}
  >{\raggedright\arraybackslash}p{(\columnwidth - 10\tabcolsep) * \real{0.1667}}
  >{\raggedright\arraybackslash}p{(\columnwidth - 10\tabcolsep) * \real{0.1667}}
  >{\raggedright\arraybackslash}p{(\columnwidth - 10\tabcolsep) * \real{0.1667}}
  >{\raggedright\arraybackslash}p{(\columnwidth - 10\tabcolsep) * \real{0.1667}}
  >{\raggedright\arraybackslash}p{(\columnwidth - 10\tabcolsep) * \real{0.1667}}
  >{\raggedright\arraybackslash}p{(\columnwidth - 10\tabcolsep) * \real{0.1667}}@{}}
\toprule\noalign{}
\begin{minipage}[b]{\linewidth}\raggedright
\texttt{canonical\_key}
\end{minipage} &
\begin{minipage}[b]{\linewidth}\raggedright
\texttt{lesion\_test}
\end{minipage} &
\begin{minipage}[b]{\linewidth}\raggedright
\texttt{level}
\end{minipage} &
\begin{minipage}[b]{\linewidth}\raggedright
\texttt{n\_systems}
\end{minipage} &
\begin{minipage}[b]{\linewidth}\raggedright
\texttt{sr\_median}
\end{minipage} &
\begin{minipage}[b]{\linewidth}\raggedright
\texttt{verdict}
\end{minipage} \\
\midrule\noalign{}
\endhead
\bottomrule\noalign{}
\endlastfoot
\texttt{frac\_A} &
\texttt{catalytic\ lesion}
& \texttt{ala} &
\texttt{143} &
\texttt{1.1744} &
\texttt{non\_specific} \\
\texttt{frac\_C} &
\texttt{catalytic\ lesion}
& \texttt{ala} &
\texttt{143} &
\texttt{nan} &
\texttt{blind} \\
\texttt{frac\_D} &
\texttt{catalytic\ lesion}
& \texttt{ala} &
\texttt{143} &
\texttt{0.9842} &
\texttt{non\_specific} \\
\end{longtable}

}

\hypertarget{S35}{%
\subsubsection{Table
S35. The specific
cells that survive
the
controls}\label{S35}}

\begin{shaded}\small

\textbf{File.}
\texttt{supplementary/S35\_survivors\_by\_status.csv}
· 97 rows × 10
columns · public

\textbf{Supports.}
Methods 2.5.1;
Results 3.4.1;
Discussion.
\textbf{Cited in.}
Methods 2.5.1;
Results 3.4.1.

\textbf{A row is}
one metric at one
step in one analysis
(97 rows: 74
natural, 23 de
novo).

\end{shaded}

The cells that are
specific under the
control of their
analysis, with the
specificity ratio,
the population
(natural or de novo)
and the
applicability status
of the metric. Of
the 97, 69 are of
metrics with status
OK, 27 of
whole-protein
comparators (GLOBAL)
and 1 of a metric
with too little
coverage.

\textbf{Key
columns.}

{\small

\begin{longtable}[]{@{}
  >{\raggedright\arraybackslash}p{(\columnwidth - 2\tabcolsep) * \real{0.2909}}
  >{\raggedright\arraybackslash}p{(\columnwidth - 2\tabcolsep) * \real{0.7091}}@{}}
\toprule\noalign{}
\begin{minipage}[b]{\linewidth}\raggedright
Column
\end{minipage} &
\begin{minipage}[b]{\linewidth}\raggedright
Meaning
\end{minipage} \\
\midrule\noalign{}
\endhead
\bottomrule\noalign{}
\endlastfoot
\texttt{analysis\_set}
/
\texttt{lesion\_test}
/ \texttt{level} &
the analysis, the
lesion test and the
step \\
\texttt{metric\_name}
/
\texttt{canonical\_key}
/
\texttt{input\_class}
/
\texttt{in\_main\_set}
& the metric \\
\texttt{sr\_text} &
specificity ratio
with interval, as
text \\
\texttt{applicability\_status}
/
\texttt{population}
& applicability
status and natural
or de novo \\
\end{longtable}

}

\textbf{Excerpt}
(first 3 rows where
analysis\_set = main
set, packing-matched
controls,
in\_main\_set = 1).

{\footnotesize

\begin{longtable}[]{@{}
  >{\raggedright\arraybackslash}p{(\columnwidth - 10\tabcolsep) * \real{0.1667}}
  >{\raggedright\arraybackslash}p{(\columnwidth - 10\tabcolsep) * \real{0.1667}}
  >{\raggedright\arraybackslash}p{(\columnwidth - 10\tabcolsep) * \real{0.1667}}
  >{\raggedright\arraybackslash}p{(\columnwidth - 10\tabcolsep) * \real{0.1667}}
  >{\raggedright\arraybackslash}p{(\columnwidth - 10\tabcolsep) * \real{0.1667}}
  >{\raggedright\arraybackslash}p{(\columnwidth - 10\tabcolsep) * \real{0.1667}}@{}}
\toprule\noalign{}
\begin{minipage}[b]{\linewidth}\raggedright
\texttt{analysis\_set}
\end{minipage} &
\begin{minipage}[b]{\linewidth}\raggedright
\texttt{lesion\_test}
\end{minipage} &
\begin{minipage}[b]{\linewidth}\raggedright
\texttt{level}
\end{minipage} &
\begin{minipage}[b]{\linewidth}\raggedright
\texttt{metric\_name}
\end{minipage} &
\begin{minipage}[b]{\linewidth}\raggedright
\texttt{sr\_text}
\end{minipage} &
\begin{minipage}[b]{\linewidth}\raggedright
\texttt{population}
\end{minipage} \\
\midrule\noalign{}
\endhead
\bottomrule\noalign{}
\endlastfoot
\texttt{main\ set,\ packing-matched\ controls}
&
\texttt{catalytic\ lesion}
& \texttt{ala} &
\texttt{catalytic\_frac\_buried}
&
\texttt{1.91\ {[}1.4,\ 2.72{]}}
&
\texttt{natural} \\
\texttt{main\ set,\ packing-matched\ controls}
&
\texttt{catalytic\ lesion}
& \texttt{ala} &
\texttt{catalytic\_radius\_of\_gyration}
&
\texttt{1.76\ {[}1.43,\ 2.19{]}}
&
\texttt{natural} \\
\texttt{main\ set,\ packing-matched\ controls}
&
\texttt{catalytic\ lesion}
& \texttt{ala} &
\texttt{catalytic\_rel\_sasa\_mean}
&
\texttt{2.4\ {[}1.93,\ 3.03{]}}
&
\texttt{natural} \\
\end{longtable}

}

\hypertarget{S36}{%
\subsubsection{Table
S36. Rank statistic
R per lesion step on
the 143 main
enzymes}\label{S36}}

\begin{shaded}\small

\textbf{File.}
\texttt{supplementary/S36\_ranking\_statistic\_by\_lesion\_step.csv}
· 313 rows × 14
columns · public

\textbf{Supports.}
Methods 2.5.1 (Eq.
19). \textbf{Cited
in.} Methods 2.5.1;
Results 3.4.1.

\textbf{A row is}
one metric at one
step of the
catalytic lesion
(313 rows).

\end{shaded}

The rank statistic R
(probability that a
metric changes more
at the catalytic
lesion than at the
matched control) for
the 143 main enzymes
alone, with 90\%
intervals, for 83
distinct metrics
including
comparators. This is
the per-step form of
R on the main
enzymes; Table 10
and
\protect\hyperlink{S28}{S28}
report the pooled
analysis on 195
enzymes with 95\%
intervals and the
counted-step rule,
so the values for
the same metric
differ.

\textbf{Key
columns.}

{\small

\begin{longtable}[]{@{}
  >{\raggedright\arraybackslash}p{(\columnwidth - 2\tabcolsep) * \real{0.2909}}
  >{\raggedright\arraybackslash}p{(\columnwidth - 2\tabcolsep) * \real{0.7091}}@{}}
\toprule\noalign{}
\begin{minipage}[b]{\linewidth}\raggedright
Column
\end{minipage} &
\begin{minipage}[b]{\linewidth}\raggedright
Meaning
\end{minipage} \\
\midrule\noalign{}
\endhead
\bottomrule\noalign{}
\endlastfoot
\texttt{metric\_name}
/
\texttt{canonical\_key}
/
\texttt{input\_class}
/
\texttt{in\_main\_set}
& the metric \\
\texttt{lesion\_step}
/
\texttt{volume\_A3}
& lesion step and
its median
side-chain volume
change \\
\texttt{R} /
\texttt{R\_lo90} /
\texttt{R\_hi90} /
\texttt{n\_systems}
/
\texttt{n\_clusters}
& R with 90\%
interval, enzymes
and enzyme
sub-subclasses \\
\texttt{responds} /
\texttt{interpretable}
& whether the metric
responds and whether
R is
interpretable \\
\texttt{applicability\_status}
& applicability
status \\
\end{longtable}

}

\textbf{Excerpt}
(first 3 rows where
in\_main\_set = 1,
lesion\_step = ala).

{\footnotesize

\begin{longtable}[]{@{}
  >{\raggedright\arraybackslash}p{(\columnwidth - 10\tabcolsep) * \real{0.1667}}
  >{\raggedright\arraybackslash}p{(\columnwidth - 10\tabcolsep) * \real{0.1667}}
  >{\raggedright\arraybackslash}p{(\columnwidth - 10\tabcolsep) * \real{0.1667}}
  >{\raggedright\arraybackslash}p{(\columnwidth - 10\tabcolsep) * \real{0.1667}}
  >{\raggedright\arraybackslash}p{(\columnwidth - 10\tabcolsep) * \real{0.1667}}
  >{\raggedright\arraybackslash}p{(\columnwidth - 10\tabcolsep) * \real{0.1667}}@{}}
\toprule\noalign{}
\begin{minipage}[b]{\linewidth}\raggedright
\texttt{metric\_name}
\end{minipage} &
\begin{minipage}[b]{\linewidth}\raggedright
\texttt{lesion\_step}
\end{minipage} &
\begin{minipage}[b]{\linewidth}\raggedright
\texttt{R}
\end{minipage} &
\begin{minipage}[b]{\linewidth}\raggedright
\texttt{R\_lo90}
\end{minipage} &
\begin{minipage}[b]{\linewidth}\raggedright
\texttt{R\_hi90}
\end{minipage} &
\begin{minipage}[b]{\linewidth}\raggedright
\texttt{n\_systems}
\end{minipage} \\
\midrule\noalign{}
\endhead
\bottomrule\noalign{}
\endlastfoot
\texttt{catalytic\_frac\_buried}
& \texttt{ala} &
\texttt{0.7557} &
\texttt{0.686} &
\texttt{0.8189} &
\texttt{131} \\
\texttt{catalytic\_pairwise\_dist\_max}
& \texttt{ala} &
\texttt{0.203} &
\texttt{0.1522} &
\texttt{0.2574} &
\texttt{133} \\
\texttt{catalytic\_pairwise\_dist\_mean}
& \texttt{ala} &
\texttt{0.3262} &
\texttt{0.2553} &
\texttt{0.4} &
\texttt{141} \\
\end{longtable}

}

\hypertarget{S37}{%
\subsubsection{Table
S37. Re-predicted
ladder: distance of
the substituted and
control residues to
the
ligand}\label{S37}}

\begin{shaded}\small

\textbf{File.}
\texttt{supplementary/S37\_ladder\_pair\_distances.csv}
· 312 rows × 34
columns · public

\textbf{Supports.}
Methods 2.5.1;
Results 3.4.1.
\textbf{Cited in.}
Methods 2.5.1;
Results 3.4.1.

\textbf{A row is}
one
catalytic/control
pair of the ladder
(312 pairs: 286
natural, 26 de
novo).

\end{shaded}

For every pair the
residue types, the
burial of each
residue, the
distance of the
substituted and of
the control residue
to the ligand (in
the reference
structure and in the
prediction), whether
the control is
adjacent to a
catalytic residue,
and whether the pair
falls in each of the
calipers used in
\protect\hyperlink{S40}{S40}.

\textbf{Key
columns.}

{\small

\begin{longtable}[]{@{}
  >{\raggedright\arraybackslash}p{(\columnwidth - 2\tabcolsep) * \real{0.2909}}
  >{\raggedright\arraybackslash}p{(\columnwidth - 2\tabcolsep) * \real{0.7091}}@{}}
\toprule\noalign{}
\begin{minipage}[b]{\linewidth}\raggedright
Column
\end{minipage} &
\begin{minipage}[b]{\linewidth}\raggedright
Meaning
\end{minipage} \\
\midrule\noalign{}
\endhead
\bottomrule\noalign{}
\endlastfoot
\texttt{source} /
\texttt{base\_system}
/ \texttt{level} &
natural or denovo,
the system and the
lesion step \\
\texttt{wt} /
\texttt{new\_aa} /
\texttt{cat\_position}
/
\texttt{ctl\_position}
& residue
substituted and its
position in the
catalytic and
control arms \\
\texttt{cat\_rel\_sasa}
/
\texttt{ctl\_rel\_sasa}
/
\texttt{abs\_diff\_rel\_sasa}
& relative
accessible area of
the two residues and
the difference \\
\texttt{d\_lig\_cat\_ref}
/
\texttt{d\_lig\_ctl\_ref}
/
\texttt{abs\_diff\_d\_lig\_ref}
& distance to the
ligand in the
reference structure
(A) \\
\texttt{d\_lig\_cat\_chai}
/
\texttt{d\_lig\_ctl\_chai}
& distance in the
Chai-1 prediction \\
\texttt{control\_adjacent\_to\_catalytic}
/
\texttt{ligand\_status}
& adjacency to a
catalytic residue;
whether the ligand
could be located \\
in\_caliper\_* &
membership of the
primary caliper and
of sensitivity
calipers \\
\end{longtable}

}

\textbf{Excerpt}
(first 3 rows where
source = natural,
level = isosteric,
ligand\_status =
ok).

{\footnotesize

\begin{longtable}[]{@{}
  >{\raggedright\arraybackslash}p{(\columnwidth - 12\tabcolsep) * \real{0.1429}}
  >{\raggedright\arraybackslash}p{(\columnwidth - 12\tabcolsep) * \real{0.1429}}
  >{\raggedright\arraybackslash}p{(\columnwidth - 12\tabcolsep) * \real{0.1429}}
  >{\raggedright\arraybackslash}p{(\columnwidth - 12\tabcolsep) * \real{0.1429}}
  >{\raggedright\arraybackslash}p{(\columnwidth - 12\tabcolsep) * \real{0.1429}}
  >{\raggedright\arraybackslash}p{(\columnwidth - 12\tabcolsep) * \real{0.1429}}
  >{\raggedright\arraybackslash}p{(\columnwidth - 12\tabcolsep) * \real{0.1429}}@{}}
\toprule\noalign{}
\begin{minipage}[b]{\linewidth}\raggedright
\texttt{base\_system}
\end{minipage} &
\begin{minipage}[b]{\linewidth}\raggedright
\texttt{level}
\end{minipage} &
\begin{minipage}[b]{\linewidth}\raggedright
\texttt{wt}
\end{minipage} &
\begin{minipage}[b]{\linewidth}\raggedright
\texttt{new\_aa}
\end{minipage} &
\begin{minipage}[b]{\linewidth}\raggedright
\texttt{d\_lig\_cat\_ref}
\end{minipage} &
\begin{minipage}[b]{\linewidth}\raggedright
\texttt{d\_lig\_ctl\_ref}
\end{minipage} &
\begin{minipage}[b]{\linewidth}\raggedright
\texttt{abs\_diff\_rel\_sasa}
\end{minipage} \\
\midrule\noalign{}
\endhead
\bottomrule\noalign{}
\endlastfoot
\texttt{M0024\_1nzy\_a2}
& \texttt{isosteric}
& \texttt{F} &
\texttt{L} &
\texttt{3.433} &
\texttt{24.224} &
\texttt{0.0031} \\
\texttt{M0024\_1nzy\_a2}
& \texttt{isosteric}
& \texttt{F} &
\texttt{Y} &
\texttt{3.433} &
\texttt{10.504} &
\texttt{0.0008} \\
\texttt{M0050\_1dbt\_a1}
& \texttt{isosteric}
& \texttt{D} &
\texttt{N} &
\texttt{3.703} &
\texttt{17.378} &
\texttt{0.0002} \\
\end{longtable}

}

\textbf{Notes.}

\begin{itemize}
\tightlist
\item
  Released as a CSV
  only for the full
  312 rows.
\end{itemize}

\hypertarget{S38}{%
\subsubsection{Table
S38. Re-predicted
ladder: candidates
for a new
distance-matched
control}\label{S38}}

\begin{shaded}\small

\textbf{File.}
\texttt{supplementary/S38\_new\_distance\_matched\_controls.csv}
· 291 rows × 17
columns · public

\textbf{Supports.}
Methods 2.5.1;
Results 3.4.1.
\textbf{Cited in.}
Methods 2.5.1;
Results 3.4.1.

\textbf{A row is}
one candidate
control for one pair
that had no control
within the caliper
(291 rows).

\end{shaded}

The candidate
controls for the
pairs that had none
in the caliper,
which were selected
(54 pairs, 53 new
predictions), their
distances and
burial, and the
reason when none was
selected (no
candidate in the
caliper for 213
pairs; no reference
distance for 20; no
ligand in the
baseline prediction
for 4).

\textbf{Key
columns.}

{\small

\begin{longtable}[]{@{}
  >{\raggedright\arraybackslash}p{(\columnwidth - 2\tabcolsep) * \real{0.2909}}
  >{\raggedright\arraybackslash}p{(\columnwidth - 2\tabcolsep) * \real{0.7091}}@{}}
\toprule\noalign{}
\begin{minipage}[b]{\linewidth}\raggedright
Column
\end{minipage} &
\begin{minipage}[b]{\linewidth}\raggedright
Meaning
\end{minipage} \\
\midrule\noalign{}
\endhead
\bottomrule\noalign{}
\endlastfoot
\texttt{base\_system}
/ \texttt{source} /
\texttt{level} &
system, natural or
de novo, step \\
\texttt{wt} /
\texttt{new\_aa} /
\texttt{cat\_position}
/
\texttt{orig\_ctl\_position}
/
\texttt{new\_ctl\_position}
& the substitution
and the original and
new control
positions \\
\texttt{d\_lig\_cat\_ref}
/
\texttt{d\_lig\_ctl\_new}
/
\texttt{abs\_diff\_d}
& distances to the
ligand and their
difference \\
\texttt{cat\_rel\_sasa}
/
\texttt{ctl\_rel\_sasa\_new}
/
\texttt{abs\_diff\_rel}
& burial and its
difference \\
\texttt{n\_candidates}
/ \texttt{selected}
/ \texttt{reason} &
number of
candidates, 1 if
selected, and the
reason when not \\
\end{longtable}

}

\textbf{Excerpt}
(first 3 rows where
selected = 1).

{\footnotesize

\begin{longtable}[]{@{}
  >{\raggedright\arraybackslash}p{(\columnwidth - 12\tabcolsep) * \real{0.1429}}
  >{\raggedright\arraybackslash}p{(\columnwidth - 12\tabcolsep) * \real{0.1429}}
  >{\raggedright\arraybackslash}p{(\columnwidth - 12\tabcolsep) * \real{0.1429}}
  >{\raggedright\arraybackslash}p{(\columnwidth - 12\tabcolsep) * \real{0.1429}}
  >{\raggedright\arraybackslash}p{(\columnwidth - 12\tabcolsep) * \real{0.1429}}
  >{\raggedright\arraybackslash}p{(\columnwidth - 12\tabcolsep) * \real{0.1429}}
  >{\raggedright\arraybackslash}p{(\columnwidth - 12\tabcolsep) * \real{0.1429}}@{}}
\toprule\noalign{}
\begin{minipage}[b]{\linewidth}\raggedright
\texttt{base\_system}
\end{minipage} &
\begin{minipage}[b]{\linewidth}\raggedright
\texttt{level}
\end{minipage} &
\begin{minipage}[b]{\linewidth}\raggedright
\texttt{new\_ctl\_position}
\end{minipage} &
\begin{minipage}[b]{\linewidth}\raggedright
\texttt{d\_lig\_cat\_ref}
\end{minipage} &
\begin{minipage}[b]{\linewidth}\raggedright
\texttt{d\_lig\_ctl\_new}
\end{minipage} &
\begin{minipage}[b]{\linewidth}\raggedright
\texttt{selected}
\end{minipage} &
\begin{minipage}[b]{\linewidth}\raggedright
\texttt{reason}
\end{minipage} \\
\midrule\noalign{}
\endhead
\bottomrule\noalign{}
\endlastfoot
\texttt{DN\_campaign1\_\_E4}
& \texttt{ala} &
\texttt{30} &
\texttt{2.16} &
\texttt{2.6964} &
\texttt{1} &
\texttt{selected} \\
\texttt{M0024\_1nzy\_a2}
& \texttt{ala} &
\texttt{82} &
\texttt{3.433} &
\texttt{4.512} &
\texttt{1} &
\texttt{selected} \\
\texttt{M0050\_1dbt\_a1}
& \texttt{isosteric}
& \texttt{193} &
\texttt{3.703} &
\texttt{5.0744} &
\texttt{1} &
\texttt{selected} \\
\end{longtable}

}

\hypertarget{S39}{%
\subsubsection{Table
S39. Re-predicted
ladder: specificity
ratio in strata of
the distance to the
ligand}\label{S39}}

\begin{shaded}\small

\textbf{File.}
\texttt{supplementary/S39\_ladder\_distance\_strata.csv}
· 203 rows × 27
columns · public

\textbf{Supports.}
Methods 2.5.1;
Results 3.4.1.
\textbf{Cited in.}
Methods 2.5.1;
Results 3.4.1.

\textbf{A row is}
one metric in one
distance stratum,
for one source (203
rows).

\end{shaded}

The specificity
ratio of each metric
in the near, mid and
far third of the
pairs by the
distance of the
substituted residue
to the ligand (edges
2.73 and 3.66 A),
and on all pairs,
with the median
distances of the two
arms.

\textbf{Key
columns.}

{\small

\begin{longtable}[]{@{}
  >{\raggedright\arraybackslash}p{(\columnwidth - 2\tabcolsep) * \real{0.2909}}
  >{\raggedright\arraybackslash}p{(\columnwidth - 2\tabcolsep) * \real{0.7091}}@{}}
\toprule\noalign{}
\begin{minipage}[b]{\linewidth}\raggedright
Column
\end{minipage} &
\begin{minipage}[b]{\linewidth}\raggedright
Meaning
\end{minipage} \\
\midrule\noalign{}
\endhead
\bottomrule\noalign{}
\endlastfoot
\texttt{source} /
\texttt{metric\_name}
/ \texttt{role} /
\texttt{applicability\_status}
/ \texttt{headline}
& the metric,
interface or global,
applicability status
and whether it is a
headline row \\
\texttt{stratum} /
\texttt{edge\_lo} /
\texttt{edge\_hi} /
\texttt{k\_pairs} /
\texttt{k\_systems}
& the stratum and
its size \\
SR\_mean (+ lo/hi 95
and 90) / SR\_median
(+ lo/hi 95) &
specificity ratio of
the mean and of the
median change with
intervals \\
\texttt{mean\_abs\_delta\_cat}
/
\texttt{mean\_abs\_delta\_ctl}
& mean absolute
change in the two
arms \\
\texttt{median\_d\_lig\_cat}
/
\texttt{median\_d\_lig\_ctl}
/
\texttt{median\_abs\_diff\_d\_lig}
& median distances
to the ligand \\
\end{longtable}

}

\textbf{Excerpt}
(first 3 rows where
source = natural,
metric\_name =
chai\_iptm\_mean).

{\footnotesize

\begin{longtable}[]{@{}
  >{\raggedright\arraybackslash}p{(\columnwidth - 12\tabcolsep) * \real{0.1429}}
  >{\raggedright\arraybackslash}p{(\columnwidth - 12\tabcolsep) * \real{0.1429}}
  >{\raggedright\arraybackslash}p{(\columnwidth - 12\tabcolsep) * \real{0.1429}}
  >{\raggedright\arraybackslash}p{(\columnwidth - 12\tabcolsep) * \real{0.1429}}
  >{\raggedright\arraybackslash}p{(\columnwidth - 12\tabcolsep) * \real{0.1429}}
  >{\raggedright\arraybackslash}p{(\columnwidth - 12\tabcolsep) * \real{0.1429}}
  >{\raggedright\arraybackslash}p{(\columnwidth - 12\tabcolsep) * \real{0.1429}}@{}}
\toprule\noalign{}
\begin{minipage}[b]{\linewidth}\raggedright
\texttt{metric\_name}
\end{minipage} &
\begin{minipage}[b]{\linewidth}\raggedright
\texttt{stratum}
\end{minipage} &
\begin{minipage}[b]{\linewidth}\raggedright
\texttt{k\_pairs}
\end{minipage} &
\begin{minipage}[b]{\linewidth}\raggedright
\texttt{k\_systems}
\end{minipage} &
\begin{minipage}[b]{\linewidth}\raggedright
\texttt{SR\_mean}
\end{minipage} &
\begin{minipage}[b]{\linewidth}\raggedright
\texttt{SR\_mean\_lo95}
\end{minipage} &
\begin{minipage}[b]{\linewidth}\raggedright
\texttt{SR\_mean\_hi95}
\end{minipage} \\
\midrule\noalign{}
\endhead
\bottomrule\noalign{}
\endlastfoot
\texttt{chai\_iptm\_mean}
&
\texttt{all\_pairs\_original\_gate}
& \texttt{286} &
\texttt{59} &
\texttt{2.0359} &
\texttt{1.4653} &
\texttt{2.8031} \\
\texttt{chai\_iptm\_mean}
&
\texttt{all\_pairs}
& \texttt{281} &
\texttt{58} &
\texttt{2.0359} &
\texttt{1.4532} &
\texttt{2.833} \\
\texttt{chai\_iptm\_mean}
&
\texttt{near\_third}
& \texttt{89} &
\texttt{27} &
\texttt{2.4894} &
\texttt{1.2129} &
\texttt{4.7288} \\
\end{longtable}

}

\hypertarget{S40}{%
\subsubsection{Table
S40. Re-predicted
ladder:
caliper-matched
subsets and the rule
labels}\label{S40}}

\begin{shaded}\small

\textbf{File.}
\texttt{supplementary/S40\_ladder\_matched\_subsets.csv}
· 348 rows × 25
columns · public

\textbf{Supports.}
Methods 2.5.1;
Results 3.4.1 (68 to
69 pairs in 21 to 22
systems).
\textbf{Cited in.}
Methods 2.5.1;
Results 3.4.1.

\textbf{A row is}
one metric in one
matched subset for
one source (348
rows).

\end{shaded}

The specificity
ratio in
caliper-matched
subsets:
original\_controls
(original controls
within the caliper),
original\_plus\_new\_controls
(with the new
controls added, the
subset of Table 10)
and the sensitivity
sets
distance\_in\_prediction,
distance\_to\_organic\_atoms,
caliper\_1A and
caliper\_3A that use
other distance
definitions and
calipers, with the
label from the
pre-specified rule
(lower 90\% bound
above 1.2 with at
least 30 pairs in 15
systems).

\textbf{Key
columns.}

{\small

\begin{longtable}[]{@{}
  >{\raggedright\arraybackslash}p{(\columnwidth - 2\tabcolsep) * \real{0.2909}}
  >{\raggedright\arraybackslash}p{(\columnwidth - 2\tabcolsep) * \real{0.7091}}@{}}
\toprule\noalign{}
\begin{minipage}[b]{\linewidth}\raggedright
Column
\end{minipage} &
\begin{minipage}[b]{\linewidth}\raggedright
Meaning
\end{minipage} \\
\midrule\noalign{}
\endhead
\bottomrule\noalign{}
\endlastfoot
\texttt{source} /
\texttt{metric\_name}
/ \texttt{role} /
\texttt{applicability\_status}
& the metric and its
status \\
\texttt{set\_id} /
\texttt{control} &
the matched subset
and original or
original+new
controls \\
\texttt{k\_pairs} /
\texttt{k\_systems}
/ \texttt{N\_pairs}
/
\texttt{N\_systems}
& subset size and
total \\
SR\_mean (+ lo/hi) /
SR\_median (+ lo/hi)
& specificity ratio
with intervals \\
\texttt{verdict} /
\texttt{rule\_applies}
& the label (for
example
not\_testable
(coverage),
global\_comparator\_no\_verdict)
and whether the
pre-specified rule
applies \\
\end{longtable}

}

\textbf{Excerpt}
(first 3 rows where
source = natural,
metric\_name =
chai\_iptm\_mean).

{\footnotesize

\begin{longtable}[]{@{}
  >{\raggedright\arraybackslash}p{(\columnwidth - 10\tabcolsep) * \real{0.1667}}
  >{\raggedright\arraybackslash}p{(\columnwidth - 10\tabcolsep) * \real{0.1667}}
  >{\raggedright\arraybackslash}p{(\columnwidth - 10\tabcolsep) * \real{0.1667}}
  >{\raggedright\arraybackslash}p{(\columnwidth - 10\tabcolsep) * \real{0.1667}}
  >{\raggedright\arraybackslash}p{(\columnwidth - 10\tabcolsep) * \real{0.1667}}
  >{\raggedright\arraybackslash}p{(\columnwidth - 10\tabcolsep) * \real{0.1667}}@{}}
\toprule\noalign{}
\begin{minipage}[b]{\linewidth}\raggedright
\texttt{metric\_name}
\end{minipage} &
\begin{minipage}[b]{\linewidth}\raggedright
\texttt{set\_id}
\end{minipage} &
\begin{minipage}[b]{\linewidth}\raggedright
\texttt{k\_pairs}
\end{minipage} &
\begin{minipage}[b]{\linewidth}\raggedright
\texttt{k\_systems}
\end{minipage} &
\begin{minipage}[b]{\linewidth}\raggedright
\texttt{SR\_mean}
\end{minipage} &
\begin{minipage}[b]{\linewidth}\raggedright
\texttt{verdict}
\end{minipage} \\
\midrule\noalign{}
\endhead
\bottomrule\noalign{}
\endlastfoot
\texttt{chai\_iptm\_mean}
&
\texttt{original\_controls}
& \texttt{17} &
\texttt{16} &
\texttt{0.3831} &
\texttt{not\_testable\ (coverage)} \\
\texttt{chai\_iptm\_mean}
&
\texttt{distance\_in\_prediction}
& \texttt{23} &
\texttt{22} &
\texttt{0.4155} &
\texttt{sensitivity:not\_testable\ (coverag…} \\
\texttt{chai\_iptm\_mean}
&
\texttt{distance\_to\_organic\_atoms}
& \texttt{16} &
\texttt{15} &
\texttt{0.369} &
\texttt{sensitivity:not\_testable\ (coverag…} \\
\end{longtable}

}

\hypertarget{S41}{%
\subsubsection{Table
S41. Re-predicted
ladder:
regression-adjusted
ratios}\label{S41}}

\begin{shaded}\small

\textbf{File.}
\texttt{supplementary/S41\_ladder\_regression\_adjusted.csv}
· 116 rows × 28
columns · public

\textbf{Supports.}
Methods 2.5.1 (Eq.
24); Results 3.4.1.
\textbf{Cited in.}
Methods 2.5.1;
Results 3.4.1.

\textbf{A row is}
one metric with the
original or the
extended set of
controls, for one
source (116 rows).

\end{shaded}

The ratio of
catalytic to control
change adjusted for
the distance to the
ligand and the
burial,
log(\textbar delta\textbar{}
+ eps)
\textasciitilde{}
arm + d\_lig +
relSASA, beside the
unadjusted ratio on
the same pairs, and
the regression
coefficients.

\textbf{Key
columns.}

{\small

\begin{longtable}[]{@{}
  >{\raggedright\arraybackslash}p{(\columnwidth - 2\tabcolsep) * \real{0.2909}}
  >{\raggedright\arraybackslash}p{(\columnwidth - 2\tabcolsep) * \real{0.7091}}@{}}
\toprule\noalign{}
\begin{minipage}[b]{\linewidth}\raggedright
Column
\end{minipage} &
\begin{minipage}[b]{\linewidth}\raggedright
Meaning
\end{minipage} \\
\midrule\noalign{}
\endhead
\bottomrule\noalign{}
\endlastfoot
\texttt{source} /
\texttt{metric\_name}
/ \texttt{role} /
\texttt{applicability\_status}
& the metric \\
\texttt{model} /
\texttt{control} &
the regression and
original\_controls
or
plus\_new\_controls \\
\texttt{k\_pairs} /
\texttt{k\_systems}
/ \texttt{n\_rows} /
\texttt{eps} & size
and the small
constant that keeps
zero changes
finite \\
adj\_GMR\_cat\_over\_ctl
(+ lo/hi95) &
adjusted
geometric-mean
ratio \\
unadj\_GMR\_cat\_over\_ctl
(+ lo/hi95) &
unadjusted ratio on
the same pairs \\
\texttt{beta\_d\_lig\_per\_A}
/ beta\_relSASA (+
lo/hi95) &
coefficients per A
of distance and per
unit of relative
accessible area \\
\end{longtable}

}

\textbf{Excerpt}
(first 3 rows where
source = natural,
metric\_name =
chai\_iptm\_mean).

{\footnotesize

\begin{longtable}[]{@{}
  >{\raggedright\arraybackslash}p{(\columnwidth - 8\tabcolsep) * \real{0.2000}}
  >{\raggedright\arraybackslash}p{(\columnwidth - 8\tabcolsep) * \real{0.2000}}
  >{\raggedright\arraybackslash}p{(\columnwidth - 8\tabcolsep) * \real{0.2000}}
  >{\raggedright\arraybackslash}p{(\columnwidth - 8\tabcolsep) * \real{0.2000}}
  >{\raggedright\arraybackslash}p{(\columnwidth - 8\tabcolsep) * \real{0.2000}}@{}}
\toprule\noalign{}
\begin{minipage}[b]{\linewidth}\raggedright
\texttt{metric\_name}
\end{minipage} &
\begin{minipage}[b]{\linewidth}\raggedright
\texttt{control}
\end{minipage} &
\begin{minipage}[b]{\linewidth}\raggedright
\texttt{k\_pairs}
\end{minipage} &
\begin{minipage}[b]{\linewidth}\raggedright
\texttt{adj\_GMR\_cat\_over\_ctl}
\end{minipage} &
\begin{minipage}[b]{\linewidth}\raggedright
\texttt{unadj\_GMR\_cat\_over\_ctl}
\end{minipage} \\
\midrule\noalign{}
\endhead
\bottomrule\noalign{}
\endlastfoot
\texttt{chai\_iptm\_mean}
&
\texttt{original\_controls}
& \texttt{261} &
\texttt{1.1785} &
\texttt{2.0198} \\
\texttt{chai\_iptm\_mean}
&
\texttt{plus\_new\_controls}
& \texttt{261} &
\texttt{1.0882} &
\texttt{1.7922} \\
\end{longtable}

}

\hypertarget{S42}{%
\subsubsection{Table
S42. Panel
equivalence
statistic: earlier
whole panel against
the main-set
metrics}\label{S42}}

\begin{shaded}\small

\textbf{File.}
\texttt{supplementary/S42\_equivalence\_earlier\_vs\_main\_panel.csv}
· 6 rows × 10
columns · public

\textbf{Supports.}
Results 3.4.1;
Methods 2.5.1
(equivalence to 1).
\textbf{Cited in.}
Methods 2.5.1;
Results 3.4.1.

\textbf{A row is}
one analysis and set
of metrics (6 rows).

\end{shaded}

The geometric-mean
specificity ratio of
the panel with a
90\% interval
(nested bootstrap
that recomputes
every metric inside
each draw), for the
earlier whole panel
(100 to 103 metrics,
including
composition metrics
that have a ratio of
1 by arithmetic) and
for the 27 main-set
metrics with a
defined ratio
(isosteric step:
1.075 {[}0.958,
1.178{]}).

\textbf{Key
columns.}

{\small

\begin{longtable}[]{@{}
  >{\raggedright\arraybackslash}p{(\columnwidth - 2\tabcolsep) * \real{0.2909}}
  >{\raggedright\arraybackslash}p{(\columnwidth - 2\tabcolsep) * \real{0.7091}}@{}}
\toprule\noalign{}
\begin{minipage}[b]{\linewidth}\raggedright
Column
\end{minipage} &
\begin{minipage}[b]{\linewidth}\raggedright
Meaning
\end{minipage} \\
\midrule\noalign{}
\endhead
\bottomrule\noalign{}
\endlastfoot
\texttt{analysis\_set}
/
\texttt{metric\_set}
/
\texttt{n\_metrics}
& the analysis, the
set of metrics and
its size \\
\texttt{geomean\_sr}
/ \texttt{ci90\_lo}
/ \texttt{ci90\_hi}
/
\texttt{median\_sr}
& geometric-mean
ratio, 90\% interval
and median \\
\texttt{n\_clusters}
/ \texttt{status} /
\texttt{note} &
resampling clusters,
status and
remarks \\
\end{longtable}

}

\textbf{Excerpt}
(first 3 rows).

{\footnotesize

\begin{longtable}[]{@{}
  >{\raggedright\arraybackslash}p{(\columnwidth - 10\tabcolsep) * \real{0.1667}}
  >{\raggedright\arraybackslash}p{(\columnwidth - 10\tabcolsep) * \real{0.1667}}
  >{\raggedright\arraybackslash}p{(\columnwidth - 10\tabcolsep) * \real{0.1667}}
  >{\raggedright\arraybackslash}p{(\columnwidth - 10\tabcolsep) * \real{0.1667}}
  >{\raggedright\arraybackslash}p{(\columnwidth - 10\tabcolsep) * \real{0.1667}}
  >{\raggedright\arraybackslash}p{(\columnwidth - 10\tabcolsep) * \real{0.1667}}@{}}
\toprule\noalign{}
\begin{minipage}[b]{\linewidth}\raggedright
\texttt{analysis\_set}
\end{minipage} &
\begin{minipage}[b]{\linewidth}\raggedright
\texttt{metric\_set}
\end{minipage} &
\begin{minipage}[b]{\linewidth}\raggedright
\texttt{n\_metrics}
\end{minipage} &
\begin{minipage}[b]{\linewidth}\raggedright
\texttt{geomean\_sr}
\end{minipage} &
\begin{minipage}[b]{\linewidth}\raggedright
\texttt{ci90\_lo}
\end{minipage} &
\begin{minipage}[b]{\linewidth}\raggedright
\texttt{ci90\_hi}
\end{minipage} \\
\midrule\noalign{}
\endhead
\bottomrule\noalign{}
\endlastfoot
\texttt{pilot\ set}
&
\texttt{earlier\ whole\ panel\ (all\ metrics\ …}
& \texttt{103} &
\texttt{0.9642} &
\texttt{0.7479} &
\texttt{1.2195} \\
\texttt{main\ set} &
\texttt{earlier\ whole\ panel\ (all\ metrics\ …}
& \texttt{102} &
\texttt{0.9575} &
\texttt{0.8969} &
\texttt{1.0839} \\
\texttt{main\ set,\ packing-matched\ controls}
&
\texttt{earlier\ whole\ panel\ (all\ metrics\ …}
& \texttt{102} &
\texttt{0.9947} &
\texttt{0.9374} &
\texttt{1.1342} \\
\end{longtable}

}

\hypertarget{S43}{%
\subsubsection{Table
S43. Re-predicted
ladder: specificity
ratio per step for
all
metrics}\label{S43}}

\begin{shaded}\small

\textbf{File.}
\texttt{supplementary/S43\_ladder\_re\_predicted\_per\_level.csv}
· 231 rows × 16
columns · public

\textbf{Supports.}
Methods 2.5.1;
Results 3.4.1.
\textbf{Cited in.}
Methods 2.5.1;
Results 3.4.1.

\textbf{A row is}
one metric at one
step for one source
(231 rows).

\end{shaded}

The re-predicted
catalytic-lesion
ladder per step (ALL
steps pooled,
isosteric,
non-isosteric, Ala,
Gly) for all 44
metrics, natural and
de novo: pairs and
systems, specificity
ratio (median and
mean) with
intervals, and the
verdict of the
earlier rule.

\textbf{Key
columns.}

{\small

\begin{longtable}[]{@{}
  >{\raggedright\arraybackslash}p{(\columnwidth - 2\tabcolsep) * \real{0.2909}}
  >{\raggedright\arraybackslash}p{(\columnwidth - 2\tabcolsep) * \real{0.7091}}@{}}
\toprule\noalign{}
\begin{minipage}[b]{\linewidth}\raggedright
Column
\end{minipage} &
\begin{minipage}[b]{\linewidth}\raggedright
Meaning
\end{minipage} \\
\midrule\noalign{}
\endhead
\bottomrule\noalign{}
\endlastfoot
\texttt{source} /
\texttt{metric\_name}
/
\texttt{canonical\_key}
/
\texttt{input\_class}
/
\texttt{ame\_flavour}
& the metric \\
\texttt{level} /
\texttt{n\_pairs} /
\texttt{n\_systems}
& step and size \\
\texttt{sr\_median}
/ \texttt{sr\_lo} /
\texttt{sr\_hi} &
ratio of medians
with interval \\
\texttt{sr\_mean} /
\texttt{sr\_mean\_lo}
/
\texttt{sr\_mean\_hi}
& ratio of means
with interval \\
\texttt{verdict\_earlier\_rule}
/
\texttt{applicability\_status}
& verdict of the
earlier rule and
applicability
status \\
\end{longtable}

}

\textbf{Excerpt}
(first 3 rows where
source = natural,
level = isosteric,
canonical\_key =
chai\_iptm).

{\footnotesize

\begin{longtable}[]{@{}
  >{\raggedright\arraybackslash}p{(\columnwidth - 12\tabcolsep) * \real{0.1429}}
  >{\raggedright\arraybackslash}p{(\columnwidth - 12\tabcolsep) * \real{0.1429}}
  >{\raggedright\arraybackslash}p{(\columnwidth - 12\tabcolsep) * \real{0.1429}}
  >{\raggedright\arraybackslash}p{(\columnwidth - 12\tabcolsep) * \real{0.1429}}
  >{\raggedright\arraybackslash}p{(\columnwidth - 12\tabcolsep) * \real{0.1429}}
  >{\raggedright\arraybackslash}p{(\columnwidth - 12\tabcolsep) * \real{0.1429}}
  >{\raggedright\arraybackslash}p{(\columnwidth - 12\tabcolsep) * \real{0.1429}}@{}}
\toprule\noalign{}
\begin{minipage}[b]{\linewidth}\raggedright
\texttt{metric\_name}
\end{minipage} &
\begin{minipage}[b]{\linewidth}\raggedright
\texttt{level}
\end{minipage} &
\begin{minipage}[b]{\linewidth}\raggedright
\texttt{n\_pairs}
\end{minipage} &
\begin{minipage}[b]{\linewidth}\raggedright
\texttt{n\_systems}
\end{minipage} &
\begin{minipage}[b]{\linewidth}\raggedright
\texttt{sr\_mean}
\end{minipage} &
\begin{minipage}[b]{\linewidth}\raggedright
\texttt{sr\_mean\_lo}
\end{minipage} &
\begin{minipage}[b]{\linewidth}\raggedright
\texttt{sr\_mean\_hi}
\end{minipage} \\
\midrule\noalign{}
\endhead
\bottomrule\noalign{}
\endlastfoot
\texttt{chai\_iptm\_max}
& \texttt{isosteric}
& \texttt{118} &
\texttt{59} &
\texttt{1.622} &
\texttt{1.058} &
\texttt{2.526} \\
\texttt{chai\_iptm\_mean}
& \texttt{isosteric}
& \texttt{118} &
\texttt{59} &
\texttt{1.963} &
\texttt{1.44} &
\texttt{2.707} \\
\texttt{chai\_iptm\_std}
& \texttt{isosteric}
& \texttt{118} &
\texttt{59} &
\texttt{1.046} &
\texttt{0.605} &
\texttt{1.878} \\
\end{longtable}

}

\hypertarget{S44}{%
\subsubsection{Table
S44. Second-shell
lesion: every metric
scored (Table 11
plus supplementary
comparators and
per-dose
flags)}\label{S44}}

\begin{shaded}\small

\textbf{File.}
\texttt{supplementary/S44\_second\_shell\_lesion\_all\_metrics\_ranked.csv}
· 37 rows × 59
columns · public

\textbf{Supports.}
Results 3.4.2; the
complete version of
Table 11.
\textbf{Cited in.}
Methods 2.5.2;
Results 3.4.2.

\textbf{A row is}
one metric scored on
second-shell removal
(37 in all).

\end{shaded}

Table 11 with the 8
supplementary
comparators and the
flags at each dose
(1, 2 and 4 residues
removed): R with
interval, the number
of informative
proteins, and
whether the metric
responds, whether
the control response
is above its noise
and whether it is
specific. No step
counts for any
metric, so no rank
on this test is
interpretable and no
reference level is
computed.

\textbf{Key
columns.}

{\small

\begin{longtable}[]{@{}
  >{\raggedright\arraybackslash}p{(\columnwidth - 2\tabcolsep) * \real{0.2909}}
  >{\raggedright\arraybackslash}p{(\columnwidth - 2\tabcolsep) * \real{0.7091}}@{}}
\toprule\noalign{}
\begin{minipage}[b]{\linewidth}\raggedright
Column
\end{minipage} &
\begin{minipage}[b]{\linewidth}\raggedright
Meaning
\end{minipage} \\
\midrule\noalign{}
\endhead
\bottomrule\noalign{}
\endlastfoot
\texttt{detection\_group}
/
\texttt{rank\_in\_group}
/ \texttt{rank} /
\texttt{rank\_in\_table\_11}
& group by first
dose detected, rank
in the group,
overall rank and
rank among the 29
items of Table 11 \\
\texttt{metric\_name}
/
\texttt{canonical\_key}
/ \texttt{source} /
\texttt{input\_class}
/ \texttt{role} /
\texttt{in\_ranked\_set}
& the metric and its
role \\
\texttt{n\_steps} /
\texttt{n\_steps\_interpretable}
& doses with a value
and counted (0 for
every metric) \\
r\_second\_shell\_ /
\_lo / \emph{hi /
n\_second\_shell} &
R, interval and
informative proteins
at dose 1, 2 and
4 \\
\texttt{responds\_}
/
\texttt{control\_above\_noise\_}
/
specific\_second\_shell\_
& per-dose flags \\
\texttt{r\_mean\_all\_steps}
/
\texttt{first\_step\_detected}
/
\texttt{first\_step\_specific}
& mean R over all
doses and the
onsets \\
\end{longtable}

}

The shared block of
12 provenance
columns (see
Conventions) is also
present and is not
repeated above.

\textbf{Excerpt}
(first 3 rows).

{\footnotesize

\begin{longtable}[]{@{}
  >{\raggedright\arraybackslash}p{(\columnwidth - 8\tabcolsep) * \real{0.2000}}
  >{\raggedright\arraybackslash}p{(\columnwidth - 8\tabcolsep) * \real{0.2000}}
  >{\raggedright\arraybackslash}p{(\columnwidth - 8\tabcolsep) * \real{0.2000}}
  >{\raggedright\arraybackslash}p{(\columnwidth - 8\tabcolsep) * \real{0.2000}}
  >{\raggedright\arraybackslash}p{(\columnwidth - 8\tabcolsep) * \real{0.2000}}@{}}
\toprule\noalign{}
\begin{minipage}[b]{\linewidth}\raggedright
\texttt{detection\_group}
\end{minipage} &
\begin{minipage}[b]{\linewidth}\raggedright
\texttt{rank\_in\_group}
\end{minipage} &
\begin{minipage}[b]{\linewidth}\raggedright
\texttt{metric\_name}
\end{minipage} &
\begin{minipage}[b]{\linewidth}\raggedright
\texttt{r\_mean\_all\_steps}
\end{minipage} &
\begin{minipage}[b]{\linewidth}\raggedright
\texttt{n\_steps\_interpretable}
\end{minipage} \\
\midrule\noalign{}
\endhead
\bottomrule\noalign{}
\endlastfoot
\texttt{second\_shell\_4}
& \texttt{1} &
\texttt{propka\_catalytic\_shift\_mean}
& \texttt{0.7005} &
\texttt{0} \\
\texttt{second\_shell\_4}
& \texttt{2} &
\texttt{propka\_catalytic\_pka\_mean}
& \texttt{0.6915} &
\texttt{0} \\
\texttt{never} &
\texttt{1} &
\texttt{catalytic\_frac\_buried}
& \texttt{0.781} &
\texttt{0} \\
\end{longtable}

}

\textbf{Notes.}

\begin{itemize}
\tightlist
\item
  Because both arms
  are scored over
  the catalytic
  residues, a high
  raw R on this test
  measures location.
\end{itemize}

\hypertarget{S45}{%
\subsubsection{Table
S45. Drift of the
current Chai-1
engine against the
earlier
predictions}\label{S45}}

\begin{shaded}\small

\textbf{File.}
\texttt{supplementary/S45\_engine\_drift\_check.csv}
· 176 rows × 5
columns · public

\textbf{Supports.}
Results 3.2;
Discussion
(limitations).
\textbf{Cited in.}
Results 3.2.

\textbf{A row is}
one metric on one
re-run prediction
(176 rows).

\end{shaded}

A drift check: four
predictions (two
cognate predictions
of the substrate
swap and two
predictions of the
catalytic-lesion
ladder) were re-run
with the current
Chai-1 engine, and
every metric is
compared with its
stored value. Most
differences are
small (median
absolute difference
0.0007) but the
largest is 0.91,
which is why Chai-1
results are reported
with seed noise and
are not claimed to
be bitwise
reproducible.

\textbf{Key
columns.}

{\small

\begin{longtable}[]{@{}
  >{\raggedright\arraybackslash}p{(\columnwidth - 2\tabcolsep) * \real{0.2909}}
  >{\raggedright\arraybackslash}p{(\columnwidth - 2\tabcolsep) * \real{0.7091}}@{}}
\toprule\noalign{}
\begin{minipage}[b]{\linewidth}\raggedright
Column
\end{minipage} &
\begin{minipage}[b]{\linewidth}\raggedright
Meaning
\end{minipage} \\
\midrule\noalign{}
\endhead
\bottomrule\noalign{}
\endlastfoot
\texttt{prediction}
& the prediction
that was re-run
(system, condition
and seed) \\
\texttt{metric\_name}
& the metric \\
\texttt{stored} /
\texttt{rerun} /
\texttt{abs\_diff} &
stored value, value
of the re-run and
their absolute
difference \\
\end{longtable}

}

\textbf{Excerpt}
(first 3 rows).

{\footnotesize

\begin{longtable}[]{@{}
  >{\raggedright\arraybackslash}p{(\columnwidth - 8\tabcolsep) * \real{0.2000}}
  >{\raggedright\arraybackslash}p{(\columnwidth - 8\tabcolsep) * \real{0.2000}}
  >{\raggedright\arraybackslash}p{(\columnwidth - 8\tabcolsep) * \real{0.2000}}
  >{\raggedright\arraybackslash}p{(\columnwidth - 8\tabcolsep) * \real{0.2000}}
  >{\raggedright\arraybackslash}p{(\columnwidth - 8\tabcolsep) * \real{0.2000}}@{}}
\toprule\noalign{}
\begin{minipage}[b]{\linewidth}\raggedright
\texttt{prediction}
\end{minipage} &
\begin{minipage}[b]{\linewidth}\raggedright
\texttt{metric\_name}
\end{minipage} &
\begin{minipage}[b]{\linewidth}\raggedright
\texttt{stored}
\end{minipage} &
\begin{minipage}[b]{\linewidth}\raggedright
\texttt{rerun}
\end{minipage} &
\begin{minipage}[b]{\linewidth}\raggedright
\texttt{abs\_diff}
\end{minipage} \\
\midrule\noalign{}
\endhead
\bottomrule\noalign{}
\endlastfoot
\texttt{M0112\_3dfr\_a1\_\_cognate\_\_s101}
&
\texttt{ame\_crystal\_clash\_max}
& \texttt{2.722} &
\texttt{2.7284} &
\texttt{0.0064} \\
\texttt{M0112\_3dfr\_a1\_\_cognate\_\_s101}
&
\texttt{ame\_crystal\_pass}
& \texttt{1.0} &
\texttt{1.0} &
\texttt{0.0} \\
\texttt{M0112\_3dfr\_a1\_\_cognate\_\_s101}
&
\texttt{ame\_crystal\_rmsd\_max}
& \texttt{0.8037} &
\texttt{0.8055} &
\texttt{0.0018} \\
\end{longtable}

}

\hypertarget{S46}{%
\subsubsection{Table
S46. Effective
dimensionality of
the
panel}\label{S46}}

\begin{shaded}\small

\textbf{File.}
\texttt{supplementary/S46\_effective\_dimensionality.csv}
· 4 rows × 10
columns · public

\textbf{Supports.}
Discussion
(limitations).
\textbf{Cited in.}
Results 3.4.1;
Discussion.

\textbf{A row is}
one set of metrics
and one panel (4
rows).

\end{shaded}

The participation
ratio and related
measures of the
effective number of
independent metrics,
for the earlier
whole panel and for
the 29
structure-space
main-set metrics,
for native values
(7.85 of 29) and for
isosteric responses
(10.5).

\textbf{Key
columns.}

{\small

\begin{longtable}[]{@{}
  >{\raggedright\arraybackslash}p{(\columnwidth - 2\tabcolsep) * \real{0.2909}}
  >{\raggedright\arraybackslash}p{(\columnwidth - 2\tabcolsep) * \real{0.7091}}@{}}
\toprule\noalign{}
\begin{minipage}[b]{\linewidth}\raggedright
Column
\end{minipage} &
\begin{minipage}[b]{\linewidth}\raggedright
Meaning
\end{minipage} \\
\midrule\noalign{}
\endhead
\bottomrule\noalign{}
\endlastfoot
\texttt{metric\_set}
/ \texttt{panel} /
\texttt{n\_metrics}
/
\texttt{n\_systems\_complete}
& the set, the panel
and sizes \\
\texttt{participation\_ratio}
/
\texttt{effective\_rank\_fraction}
& effective number
of independent
metrics and its
share \\
\texttt{var\_explained\_pc1}
/
\texttt{n\_pcs\_for\_90pct}
/
\texttt{median\_abs\_offdiag\_r}
& variance in the
first component,
components for 90\%
and median absolute
correlation between
metrics \\
\end{longtable}

}

\textbf{Excerpt}
(first 3 rows).

{\footnotesize

\begin{longtable}[]{@{}
  >{\raggedright\arraybackslash}p{(\columnwidth - 8\tabcolsep) * \real{0.2000}}
  >{\raggedright\arraybackslash}p{(\columnwidth - 8\tabcolsep) * \real{0.2000}}
  >{\raggedright\arraybackslash}p{(\columnwidth - 8\tabcolsep) * \real{0.2000}}
  >{\raggedright\arraybackslash}p{(\columnwidth - 8\tabcolsep) * \real{0.2000}}
  >{\raggedright\arraybackslash}p{(\columnwidth - 8\tabcolsep) * \real{0.2000}}@{}}
\toprule\noalign{}
\begin{minipage}[b]{\linewidth}\raggedright
\texttt{metric\_set}
\end{minipage} &
\begin{minipage}[b]{\linewidth}\raggedright
\texttt{panel}
\end{minipage} &
\begin{minipage}[b]{\linewidth}\raggedright
\texttt{n\_metrics}
\end{minipage} &
\begin{minipage}[b]{\linewidth}\raggedright
\texttt{participation\_ratio}
\end{minipage} &
\begin{minipage}[b]{\linewidth}\raggedright
\texttt{var\_explained\_pc1}
\end{minipage} \\
\midrule\noalign{}
\endhead
\bottomrule\noalign{}
\endlastfoot
\texttt{earlier\ whole\ panel}
&
\texttt{native\ values\ across\ systems}
& \texttt{97} &
\texttt{12.12} &
\texttt{0.195} \\
\texttt{earlier\ whole\ panel}
&
\texttt{isosteric\ responses\ (catalytic\ ar…}
& \texttt{76} &
\texttt{14.2} &
\texttt{0.155} \\
\texttt{main-set\ metrics\ (one\ member\ per\ …}
&
\texttt{native\ values\ across\ systems}
& \texttt{29} &
\texttt{7.85} &
\texttt{0.207} \\
\end{longtable}

}

\hypertarget{S47}{%
\subsubsection{Table
S47. De novo
designs: burial
equivalence of the
created-site
control}\label{S47}}

\begin{shaded}\small

\textbf{File.}
\texttt{supplementary/S47\_denovo\_created\_control.csv}
· 5 rows × 11
columns · public

\textbf{Supports.}
Results 3.4.1;
Discussion
(limitations).
\textbf{Cited in.}
Results 3.4.1;
Discussion.

\textbf{A row is}
one burial metric (5
rows).

\end{shaded}

In the 30 de novo
designs the control
is a position
mutated to the
catalytic residue
type. For the five
burial metrics the
table gives the
catalytic-arm and
control-arm medians,
the paired
difference with 90\%
interval and whether
they differ.

\textbf{Key
columns.}

{\small

\begin{longtable}[]{@{}
  >{\raggedright\arraybackslash}p{(\columnwidth - 2\tabcolsep) * \real{0.2909}}
  >{\raggedright\arraybackslash}p{(\columnwidth - 2\tabcolsep) * \real{0.7091}}@{}}
\toprule\noalign{}
\begin{minipage}[b]{\linewidth}\raggedright
Column
\end{minipage} &
\begin{minipage}[b]{\linewidth}\raggedright
Meaning
\end{minipage} \\
\midrule\noalign{}
\endhead
\bottomrule\noalign{}
\endlastfoot
\texttt{metric} /
\texttt{n\_systems}
/
\texttt{n\_clusters}
& the burial metric,
designs and backbone
clusters \\
\texttt{catalytic\_arm\_median}
/
\texttt{control\_arm\_median}
& medians in the two
arms \\
\texttt{paired\_diff\_median}
/
\texttt{paired\_diff\_ci90\_lo}
/
\texttt{paired\_diff\_ci90\_hi}
/
\texttt{ratio\_median}
& paired difference
with interval and
ratio \\
\texttt{differs} /
\texttt{note} & 1 if
the arms differ \\
\end{longtable}

}

\textbf{Excerpt}
(first 3 rows).

{\footnotesize

\begin{longtable}[]{@{}
  >{\raggedright\arraybackslash}p{(\columnwidth - 10\tabcolsep) * \real{0.1667}}
  >{\raggedright\arraybackslash}p{(\columnwidth - 10\tabcolsep) * \real{0.1667}}
  >{\raggedright\arraybackslash}p{(\columnwidth - 10\tabcolsep) * \real{0.1667}}
  >{\raggedright\arraybackslash}p{(\columnwidth - 10\tabcolsep) * \real{0.1667}}
  >{\raggedright\arraybackslash}p{(\columnwidth - 10\tabcolsep) * \real{0.1667}}
  >{\raggedright\arraybackslash}p{(\columnwidth - 10\tabcolsep) * \real{0.1667}}@{}}
\toprule\noalign{}
\begin{minipage}[b]{\linewidth}\raggedright
\texttt{metric}
\end{minipage} &
\begin{minipage}[b]{\linewidth}\raggedright
\texttt{n\_systems}
\end{minipage} &
\begin{minipage}[b]{\linewidth}\raggedright
\texttt{catalytic\_arm\_median}
\end{minipage} &
\begin{minipage}[b]{\linewidth}\raggedright
\texttt{control\_arm\_median}
\end{minipage} &
\begin{minipage}[b]{\linewidth}\raggedright
\texttt{paired\_diff\_median}
\end{minipage} &
\begin{minipage}[b]{\linewidth}\raggedright
\texttt{differs}
\end{minipage} \\
\midrule\noalign{}
\endhead
\bottomrule\noalign{}
\endlastfoot
\texttt{catalytic\_rel\_sasa\_mean}
& \texttt{30} &
\texttt{0.1433} &
\texttt{0.1315} &
\texttt{0.0278} &
\texttt{0} \\
\texttt{catalytic\_rel\_sasa\_max}
& \texttt{30} &
\texttt{0.2157} &
\texttt{0.2042} &
\texttt{0.0256} &
\texttt{0} \\
\texttt{catalytic\_sasa\_mean}
& \texttt{30} &
\texttt{32.1092} &
\texttt{29.456} &
\texttt{6.2169} &
\texttt{0} \\
\end{longtable}

}

\hypertarget{S48}{%
\subsubsection{Table
S48. Experiments
that could not be
made or were
retired}\label{S48}}

\begin{shaded}\small

\textbf{File.}
\texttt{supplementary/S48\_not\_measurable.csv}
· 9 rows × 4 columns
· public

\textbf{Supports.}
Discussion
(limitations).
\textbf{Cited in.}
Results 3.4.1.

\textbf{A row is}
one experiment (9
rows).

\end{shaded}

The deformation
experiment (apparent
specific cells track
the clash burden
present before
repacking),
oxyanion-hole
removal (the step
that was run located
no bound ligand in
any of the 143
structures and
removed a proxy
residue, so it did
not test the named
change), metal
removal (no cell
could be scored on
any metric), the de
novo second shell
and deformation
experiments, PLACER
on the isosteric
step, the
unrelated-protein
sanity floor and the
two trapping-mutant
tiers.

\textbf{Key
columns.}

{\small

\begin{longtable}[]{@{}
  >{\raggedright\arraybackslash}p{(\columnwidth - 2\tabcolsep) * \real{0.2909}}
  >{\raggedright\arraybackslash}p{(\columnwidth - 2\tabcolsep) * \real{0.7091}}@{}}
\toprule\noalign{}
\begin{minipage}[b]{\linewidth}\raggedright
Column
\end{minipage} &
\begin{minipage}[b]{\linewidth}\raggedright
Meaning
\end{minipage} \\
\midrule\noalign{}
\endhead
\bottomrule\noalign{}
\endlastfoot
\texttt{reader\_name}
/
\texttt{claim\_type}
& the experiment and
the kind of claim it
would have
supported \\
\texttt{role} &
retired or
not\_measurable \\
\texttt{mechanism\_or\_caveat}
& why it could not
be made \\
\end{longtable}

}

\textbf{Excerpt}
(first 3 rows).

{\footnotesize

\begin{longtable}[]{@{}
  >{\raggedright\arraybackslash}p{(\columnwidth - 4\tabcolsep) * \real{0.3333}}
  >{\raggedright\arraybackslash}p{(\columnwidth - 4\tabcolsep) * \real{0.3333}}
  >{\raggedright\arraybackslash}p{(\columnwidth - 4\tabcolsep) * \real{0.3333}}@{}}
\toprule\noalign{}
\begin{minipage}[b]{\linewidth}\raggedright
\texttt{reader\_name}
\end{minipage} &
\begin{minipage}[b]{\linewidth}\raggedright
\texttt{role}
\end{minipage} &
\begin{minipage}[b]{\linewidth}\raggedright
\texttt{mechanism\_or\_caveat}
\end{minipage} \\
\midrule\noalign{}
\endhead
\bottomrule\noalign{}
\endlastfoot
\texttt{sanity\ floor:\ unrelated\ protein}
& \texttt{retired} &
\texttt{signed\ paired-median\ statistic\ ca…} \\
\texttt{oxyanion-hole\ step}
& \texttt{retired} &
\texttt{well-defined\ perturbation,\ but\ no…} \\
\texttt{metal\ removal}
&
\texttt{not\_measurable}
&
\texttt{no\ scoreable\ cell\ on\ any\ metric} \\
\end{longtable}

}

\hypertarget{S49}{%
\subsubsection{Table
S49. PLACER ensemble
metrics on the
isosteric step (not
interpretable)}\label{S49}}

\begin{shaded}\small

\textbf{File.}
\texttt{supplementary/S49\_placer\_isosteric.csv}
· 50 rows × 12
columns · public

\textbf{Supports.}
Results 3.4.1;
Discussion
(limitations).
\textbf{Cited in.}
Results 3.4.1.

\textbf{A row is}
one PLACER metric in
one analysis (50
rows).

\end{shaded}

The specificity
ratio, verdict and
applicability of 25
PLACER ensemble
metrics on the
isosteric step, in
two analyses. They
are not
interpretable: the
control arm is
covered 1.6 to 3.2
times less than the
catalytic arm in
every crop
configuration, so
the count of
specific metrics
moves with the crop
anchor alone.

\textbf{Key
columns.}

{\small

\begin{longtable}[]{@{}
  >{\raggedright\arraybackslash}p{(\columnwidth - 2\tabcolsep) * \real{0.2909}}
  >{\raggedright\arraybackslash}p{(\columnwidth - 2\tabcolsep) * \real{0.7091}}@{}}
\toprule\noalign{}
\begin{minipage}[b]{\linewidth}\raggedright
Column
\end{minipage} &
\begin{minipage}[b]{\linewidth}\raggedright
Meaning
\end{minipage} \\
\midrule\noalign{}
\endhead
\bottomrule\noalign{}
\endlastfoot
\texttt{analysis\_set}
/
\texttt{metric\_name}
/ \texttt{level} /
\texttt{n\_systems}
& the analysis, the
metric, the step and
enzymes \\
\texttt{sr\_median}
/
\texttt{sr\_ci\_lo}
/
\texttt{sr\_ci\_hi}
& specificity ratio
with 95\%
interval \\
\texttt{responds} /
\texttt{specific} /
\texttt{verdict} /
\texttt{applicability\_status}
& flags, verdict
(blind,
non\_specific,
specific,
no\_dynamic\_range)
and status (GLOBAL
or X-CONFOUND) \\
\end{longtable}

}

\textbf{Excerpt}
(first 3 rows).

{\footnotesize

\begin{longtable}[]{@{}
  >{\raggedright\arraybackslash}p{(\columnwidth - 10\tabcolsep) * \real{0.1667}}
  >{\raggedright\arraybackslash}p{(\columnwidth - 10\tabcolsep) * \real{0.1667}}
  >{\raggedright\arraybackslash}p{(\columnwidth - 10\tabcolsep) * \real{0.1667}}
  >{\raggedright\arraybackslash}p{(\columnwidth - 10\tabcolsep) * \real{0.1667}}
  >{\raggedright\arraybackslash}p{(\columnwidth - 10\tabcolsep) * \real{0.1667}}
  >{\raggedright\arraybackslash}p{(\columnwidth - 10\tabcolsep) * \real{0.1667}}@{}}
\toprule\noalign{}
\begin{minipage}[b]{\linewidth}\raggedright
\texttt{analysis\_set}
\end{minipage} &
\begin{minipage}[b]{\linewidth}\raggedright
\texttt{metric\_name}
\end{minipage} &
\begin{minipage}[b]{\linewidth}\raggedright
\texttt{n\_systems}
\end{minipage} &
\begin{minipage}[b]{\linewidth}\raggedright
\texttt{sr\_median}
\end{minipage} &
\begin{minipage}[b]{\linewidth}\raggedright
\texttt{verdict}
\end{minipage} &
\begin{minipage}[b]{\linewidth}\raggedright
\texttt{applicability\_status}
\end{minipage} \\
\midrule\noalign{}
\endhead
\bottomrule\noalign{}
\endlastfoot
\texttt{main\ set,\ packing-matched\ controls}
&
\texttt{placer\_catalytic\_rmsf\_max}
& \texttt{132} &
\texttt{1.0106} &
\texttt{blind} &
\texttt{X-CONFOUND} \\
\texttt{main\ set,\ packing-matched\ controls}
&
\texttt{placer\_catalytic\_rmsf\_mean}
& \texttt{132} &
\texttt{0.9058} &
\texttt{blind} &
\texttt{X-CONFOUND} \\
\texttt{main\ set,\ packing-matched\ controls}
&
\texttt{placer\_centroid\_dev\_mean}
& \texttt{132} &
\texttt{0.9237} &
\texttt{blind} &
\texttt{X-CONFOUND} \\
\end{longtable}

}

\section{Appendix A.
Main-text tables}

The seven data
tables of the
Results (Tables 5 to
11), with the
supplementary tables
that hold the
complete version or
more detail. Tables
1 to 4 of the
Methods are
descriptive and are
not released as
files.

\begin{longtable}[]{@{}
  >{\raggedright\arraybackslash}p{(\columnwidth - 6\tabcolsep) * \real{0.0862}}
  >{\raggedright\arraybackslash}p{(\columnwidth - 6\tabcolsep) * \real{0.5172}}
  >{\raggedright\arraybackslash}p{(\columnwidth - 6\tabcolsep) * \real{0.1379}}
  >{\raggedright\arraybackslash}p{(\columnwidth - 6\tabcolsep) * \real{0.2586}}@{}}
\toprule\noalign{}
\begin{minipage}[b]{\linewidth}\raggedright
Table
\end{minipage} &
\begin{minipage}[b]{\linewidth}\raggedright
File
\end{minipage} &
\begin{minipage}[b]{\linewidth}\raggedright
Rows × columns
\end{minipage} &
\begin{minipage}[b]{\linewidth}\raggedright
More detail in
\end{minipage} \\
\midrule\noalign{}
\endhead
\bottomrule\noalign{}
\endlastfoot
\textbf{T5} &
\texttt{main/T5\_dead\_vs\_active.csv}
& 30 × 26 &
\protect\hyperlink{S11}{S11},
\protect\hyperlink{S7}{S7},
\protect\hyperlink{S8}{S8} \\
\textbf{T6} &
\texttt{main/T6\_plated\_designs\_auroc.csv}
& 30 × 28 &
\protect\hyperlink{S13}{S13} \\
\textbf{T7} &
\texttt{main/T7\_substrate\_swap.csv}
& 30 × 32 &
\protect\hyperlink{S14}{S14},
\protect\hyperlink{S15}{S15} \\
\textbf{T8} &
\texttt{main/T8\_bglb\_variants.csv}
& 30 × 27 &
\protect\hyperlink{S16}{S16},
\protect\hyperlink{S17}{S17},
\protect\hyperlink{S18}{S18},
\protect\hyperlink{S19}{S19} \\
\textbf{T9} &
\texttt{main/T9\_plated\_designs.csv}
& 30 × 26 &
\protect\hyperlink{S21}{S21},
\protect\hyperlink{S20}{S20},
\protect\hyperlink{S23}{S23},
\protect\hyperlink{S22}{S22},
\protect\hyperlink{S24}{S24} \\
\textbf{T10} &
\texttt{main/T10\_specificity\_catalytic\_lesion.csv}
& 35 × 36 &
\protect\hyperlink{S28}{S28},
\protect\hyperlink{S29}{S29},
\protect\hyperlink{S30}{S30} \\
\textbf{T11} &
\texttt{main/T11\_specificity\_second\_shell\_lesion.csv}
& 29 × 30 &
\protect\hyperlink{S44}{S44},
\protect\hyperlink{S29}{S29} \\
\end{longtable}

\section{Appendix B.
Notes on the data}

\begin{itemize}
\tightlist
\item
  Two sets of
  numbers for
  natural enzymes in
  the substrate
  swap:
  \protect\hyperlink{S14}{S14}
  uses all 59
  natural enzymes,
  Table 7 uses the
  evaluation set of
  55.
\item
  Two forms of the
  rank statistic R:
  \protect\hyperlink{S36}{S36}
  is R on the 143
  main enzymes with
  90\% intervals;
  Table 10 and
  \protect\hyperlink{S28}{S28}
  are the pooled
  analysis on 195
  enzymes with 95\%
  intervals and the
  counted-step rule.
  Values for the
  same metric
  differ.
\item
  Among the 29
  prediction-based
  metrics of the
  ladder with
  ligand-distance-matched
  controls, the
  cells of
  \protect\hyperlink{S3}{S3}
  keep the label
  NOT-RUN although
  the analysis was
  carried out
  (results in
  \protect\hyperlink{S40}{S40}
  and
  \protect\hyperlink{S41}{S41}).
\end{itemize}

\end{document}
